%% prochand.tex
%%
%% Copyright (C) 2000-2018 by Simone Piccardi.  Permission is granted to
%% copy, distribute and/or modify this document under the terms of the GNU Free
%% Documentation License, Version 1.1 or any later version published by the
%% Free Software Foundation; with the Invariant Sections being "Un preambolo",
%% with no Front-Cover Texts, and with no Back-Cover Texts.  A copy of the
%% license is included in the section entitled "GNU Free Documentation
%% License".
%%

\chapter{La gestione dei processi}
\label{cha:process_handling}

Come accennato nell'introduzione in un sistema unix-like tutte le operazioni
vengono svolte tramite opportuni processi.  In sostanza questi ultimi vengono
a costituire l'unità base per l'allocazione e l'uso delle risorse del sistema.

Nel precedente capitolo abbiamo esaminato il funzionamento di un processo come
unità a se stante, in questo esamineremo il funzionamento dei processi
all'interno del sistema. Saranno cioè affrontati i dettagli della creazione e
della terminazione dei processi, della gestione dei loro attributi e
privilegi, e di tutte le funzioni a questo connesse. Infine nella sezione
finale introdurremo alcune problematiche generiche della programmazione in
ambiente \textit{multitasking}.


\section{Le funzioni di base della gestione dei processi}
\label{sec:proc_handling}

In questa sezione tratteremo le problematiche della gestione dei processi
all'interno del sistema, illustrandone tutti i dettagli.  Inizieremo con una
panoramica dell'architettura dei processi, tratteremo poi le funzioni
elementari che permettono di leggerne gli identificatori, per poi passare alla
spiegazione delle funzioni base che si usano per la creazione e la
terminazione dei processi, e per la messa in esecuzione degli altri programmi.


\subsection{L'architettura della gestione dei processi}
\label{sec:proc_hierarchy}

A differenza di quanto avviene in altri sistemi, ad esempio nel VMS, dove la
generazione di nuovi processi è un'operazione privilegiata, una delle
caratteristiche fondanti di Unix, che esamineremo in dettaglio più avanti, è
che qualunque processo può a sua volta generarne altri. Ogni processo è
identificato presso il sistema da un numero univoco, il cosiddetto
\textit{Process ID}, o più brevemente \ids{PID}, assegnato in forma
progressiva (vedi sez.~\ref{sec:proc_pid}) quando il processo viene creato.

Una seconda caratteristica di un sistema unix-like è che la generazione di un
processo è un'operazione separata rispetto al lancio di un programma. In
genere la sequenza è sempre quella di creare un nuovo processo, il quale
eseguirà, in un passo successivo, il programma desiderato: questo è ad esempio
quello che fa la shell quando mette in esecuzione il programma che gli
indichiamo nella linea di comando.

Una terza caratteristica del sistema è che ogni processo è sempre stato
generato da un altro processo, il processo generato viene chiamato
\textit{processo figlio} (\textit{child process}) mentre quello che lo ha
generato viene chiamato \textsl{processo padre} (\textit{parent
  process}). Questo vale per tutti i processi, con una sola eccezione; dato
che ci deve essere un punto di partenza esiste un processo iniziale (che
normalmente è \cmd{/sbin/init}), che come accennato in
sez.~\ref{sec:intro_kern_and_sys} viene lanciato dal kernel alla conclusione
della fase di avvio. Essendo questo il primo processo lanciato dal sistema ha
sempre \ids{PID} uguale a 1 e non è figlio di nessun altro processo.

Ovviamente \cmd{init} è un processo particolare che in genere si occupa di
lanciare tutti gli altri processi necessari al funzionamento del sistema,
inoltre \cmd{init} è essenziale per svolgere una serie di compiti
amministrativi nelle operazioni ordinarie del sistema (torneremo su alcuni di
essi in sez.~\ref{sec:proc_termination}) e non può mai essere terminato. La
struttura del sistema comunque consente di lanciare al posto di \cmd{init}
qualunque altro programma, e in casi di emergenza (ad esempio se il file di
\cmd{init} si fosse corrotto) è ad esempio possibile lanciare una shell al suo
posto.\footnote{la cosa si fa passando la riga \cmd{init=/bin/sh} come
  parametro di avvio del kernel, l'argomento è di natura sistemistica e
  trattato in sez.~5.3 di \cite{AGL}.}

\begin{figure}[!htb]
  \footnotesize
\begin{Console}
[piccardi@gont piccardi]$ \textbf{pstree -n} 
init-+-keventd
     |-kapm-idled
     |-kreiserfsd
     |-portmap
     |-syslogd
     |-klogd
     |-named
     |-rpc.statd
     |-gpm
     |-inetd
     |-junkbuster
     |-master-+-qmgr
     |        `-pickup
     |-sshd
     |-xfs
     |-cron
     |-bash---startx---xinit-+-XFree86
     |                       `-WindowMaker-+-ssh-agent
     |                                     |-wmtime
     |                                     |-wmmon
     |                                     |-wmmount
     |                                     |-wmppp
     |                                     |-wmcube
     |                                     |-wmmixer
     |                                     |-wmgtemp
     |                                     |-wterm---bash---pstree
     |                                     `-wterm---bash-+-emacs
     |                                                    `-man---pager
     |-5*[getty]
     |-snort
     `-wwwoffled
\end{Console}
%$
  \caption{L'albero dei processi, così come riportato dal comando
    \cmd{pstree}.}
  \label{fig:proc_tree}
\end{figure}

Dato che tutti i processi attivi nel sistema sono comunque generati da
\cmd{init} o da uno dei suoi figli si possono classificare i processi con la
relazione padre/figlio in un'organizzazione gerarchica ad albero. In
fig.~\ref{fig:proc_tree} si è mostrato il risultato del comando \cmd{pstree}
che permette di visualizzare questa struttura, alla cui base c'è \cmd{init}
che è progenitore di tutti gli altri processi.\footnote{in realtà questo non è
  del tutto vero, in Linux, specialmente nelle versioni più recenti del
  kernel, ci sono alcuni processi speciali (come \cmd{keventd}, \cmd{kswapd},
  ecc.) che pur comparendo nei comandi come figli di \cmd{init}, o con
  \ids{PID} successivi ad uno, sono in realtà processi interni al kernel e che
  non rientrano in questa classificazione.}

\itindbeg{process~table}

Il kernel mantiene una tabella dei processi attivi, la cosiddetta
\textit{process table}. Per ciascun processo viene mantenuta una voce in
questa tabella, costituita da una struttura \kstruct{task\_struct}, che
contiene tutte le informazioni rilevanti per quel processo. Tutte le strutture
usate a questo scopo sono dichiarate nell'\textit{header file}
\file{linux/sched.h}, ed in fig.~\ref{fig:proc_task_struct} si è riportato uno
schema semplificato che mostra la struttura delle principali informazioni
contenute nella \texttt{task\_struct}, che in seguito incontreremo a più
riprese.

\begin{figure}[!htb]
  \centering \includegraphics[width=14cm]{img/task_struct}
  \caption{Schema semplificato dell'architettura delle strutture (
    \kstructd{task\_struct}, \kstructd{fs\_struct}, \kstructd{file\_struct})
    usate dal kernel nella gestione dei processi.}
  \label{fig:proc_task_struct}
\end{figure}

\itindend{process~table}

% TODO la task_struct è cambiata per qualche dettaglio vedi anche
% http://www.ibm.com/developerworks/linux/library/l-linux-process-management/
% TODO completare la parte su quando viene chiamato lo scheduler.

\itindbeg{scheduler}

Come accennato in sez.~\ref{sec:intro_unix_struct} è lo \textit{scheduler} che
decide quale processo mettere in esecuzione; esso viene eseguito in occasione
di dell'invocazione di ogni \textit{system call} ed per ogni interrupt
dall'hardware oltre che in una serie di altre occasioni, e può essere anche
attivato esplicitamente. Il timer di sistema provvede comunque a che esso sia
invocato periodicamente, generando un interrupt periodico secondo una
frequenza predeterminata, specificata dalla costante \const{HZ} del kernel
(torneremo su questo argomento in sez.~\ref{sec:sys_unix_time}), che assicura
che lo \textit{scheduler} venga comunque eseguito ad intervalli regolari e
possa prendere le sue decisioni.

A partire dal kernel 2.6.21 è stato introdotto anche un meccanismo
completamente diverso, detto \textit{tickless}, in cui non c'è più una
interruzione periodica con frequenza prefissata, ma ad ogni chiamata del timer
viene programmata l'interruzione successiva sulla base di una stima; in questo
modo si evita di dover eseguire un migliaio di interruzioni al secondo anche
su macchine che non stanno facendo nulla, con un forte risparmio nell'uso
dell'energia da parte del processore che può essere messo in stato di
sospensione anche per lunghi periodi di tempo.

Ma, indipendentemente dalle motivazioni per cui questo avviene, ogni volta che
viene eseguito lo \textit{scheduler} effettua il calcolo delle priorità dei
vari processi attivi (torneremo su questo in sez.~\ref{sec:proc_priority}) e
stabilisce quale di essi debba essere posto in esecuzione fino alla successiva
invocazione.

\itindend{scheduler}

\subsection{Gli identificatori dei processi}
\label{sec:proc_pid}

\itindbeg{Process~ID~(PID)}

Come accennato nella sezione precedente ogni processo viene identificato dal
sistema da un numero identificativo univoco, il \textit{process ID} o
\ids{PID}. Questo è un tipo di dato standard, \type{pid\_t}, che in genere è un
intero con segno (nel caso di Linux e della \acr{glibc} il tipo usato è
\ctyp{int}).

Il \ids{PID} viene assegnato in forma progressiva ogni volta che un nuovo
processo viene creato,\footnote{in genere viene assegnato il numero successivo
  a quello usato per l'ultimo processo creato, a meno che questo numero non
  sia già utilizzato per un altro \ids{PID}, \acr{pgid} o \acr{sid} (vedi
  sez.~\ref{sec:sess_proc_group}).} fino ad un limite che, essendo il
tradizionalmente il \ids{PID} un numero positivo memorizzato in un intero a 16
bit, arriva ad un massimo di 32768.  Oltre questo valore l'assegnazione
riparte dal numero più basso disponibile a partire da un minimo di
300,\footnote{questi valori, fino al kernel 2.4.x, erano definiti dalla macro
  \constd{PID\_MAX} nei file \file{threads.h} e \file{fork.c} dei sorgenti del
  kernel, con il 2.6.x e la nuova interfaccia per i \textit{thread} anche il
  meccanismo di allocazione dei \ids{PID} è stato modificato ed il valore
  massimo è impostabile attraverso il file \sysctlfiled{kernel/pid\_max} e di
  default vale 32768.} che serve a riservare i \ids{PID} più bassi ai processi
eseguiti direttamente dal kernel.  Per questo motivo, come visto in
sez.~\ref{sec:proc_hierarchy}, il processo di avvio (\cmd{init}) ha sempre il
\ids{PID} uguale a uno.

\itindbeg{Parent~Process~ID~(PPID)} 

Tutti i processi inoltre memorizzano anche il \ids{PID} del genitore da cui
sono stati creati, questo viene chiamato in genere \ids{PPID} (da
\textit{Parent Process ID}).  Questi due identificativi possono essere
ottenuti usando le due funzioni di sistema \funcd{getpid} e \funcd{getppid}, i
cui prototipi sono:

\begin{funcproto}{ 
\fhead{sys/types.h}
\fhead{unistd.h}
\fdecl{pid\_t getpid(void)}
\fdesc{Restituisce il \ids{PID} del processo corrente..} 
\fdecl{pid\_t getppid(void)}
\fdesc{Restituisce il \ids{PID} del padre del processo corrente.} 
}
{Entrambe le funzioni non riportano condizioni di errore.}   
\end{funcproto}

\noindent esempi dell'uso di queste funzioni sono riportati in
fig.~\ref{fig:proc_fork_code}, nel programma \file{fork\_test.c}.

Il fatto che il \ids{PID} sia un numero univoco per il sistema lo rende un
candidato per generare ulteriori indicatori associati al processo di cui
diventa possibile garantire l'unicità: ad esempio in alcune implementazioni la
funzione \func{tempnam} (si veda sez.~\ref{sec:file_temp_file}) usa il
\ids{PID} per generare un \textit{pathname} univoco, che non potrà essere
replicato da un altro processo che usi la stessa funzione. Questo utilizzo
però può risultare pericoloso, un \ids{PID} infatti è univoco solo fintanto
che un processo è attivo, una volta terminato esso potrà essere riutilizzato
da un processo completamente diverso, e di questo bisogna essere ben
consapevoli.

Tutti i processi figli dello stesso processo padre sono detti
\textit{sibling}, questa è una delle relazioni usate nel \textsl{controllo di
  sessione}, in cui si raggruppano i processi creati su uno stesso terminale,
o relativi allo stesso login. Torneremo su questo argomento in dettaglio in
cap.~\ref{cha:session}, dove esamineremo gli altri identificativi associati ad
un processo e le varie relazioni fra processi utilizzate per definire una
sessione.

Oltre al \ids{PID} e al \ids{PPID}, e a quelli che vedremo in
sez.~\ref{sec:sess_proc_group}, relativi al controllo di sessione, ad ogni
processo vengono associati degli ulteriori identificatori ed in particolare
quelli che vengono usati per il controllo di accesso.  Questi servono per
determinare se un processo può eseguire o meno le operazioni richieste, a
seconda dei privilegi e dell'identità di chi lo ha posto in esecuzione;
l'argomento è complesso e sarà affrontato in dettaglio in
sez.~\ref{sec:proc_perms}.

\itindend{Process~ID~(PID)}
\itindend{Parent~Process~ID~(PPID)} 

\subsection{La funzione \func{fork} e le funzioni di creazione dei processi}
\label{sec:proc_fork}

La funzione di sistema \func{fork} è la funzione fondamentale della gestione
dei processi: come si è detto tradizionalmente l'unico modo di creare un nuovo
processo era attraverso l'uso di questa funzione,\footnote{in realtà oggi la
  \textit{system call} usata da Linux per creare nuovi processi è \func{clone}
  (vedi \ref{sec:process_clone}), anche perché a partire dalla \acr{glibc}
  2.3.3 non viene più usata la \textit{system call} originale, ma la stessa
  \func{fork} viene implementata tramite \func{clone}, cosa che consente una
  migliore interazione coi \textit{thread}.} essa quindi riveste un ruolo
centrale tutte le volte che si devono scrivere programmi che usano il
\textit{multitasking}.\footnote{oggi questa rilevanza, con la diffusione
  dell'uso dei \textit{thread}\unavref{ che tratteremo al
    cap.~\ref{cha:threads}}, è in parte minore, ma \func{fork} resta comunque
  la funzione principale per la creazione di processi.} Il prototipo di
\funcd{fork} è:

\begin{funcproto}{ 
\fhead{unistd.h}
\fdecl{pid\_t fork(void)}
\fdesc{Crea un nuovo processo.} 
}

{La funzione ritorna in caso di successo il \ids{PID} del figlio nel padre e
  $0$ nel figlio mentre ritorna $-1$ nel padre, senza creare il figlio, per un
  errore, al caso \var{errno} assumerà uno dei valori:
  \begin{errlist}
  \item[\errcode{EAGAIN}] non ci sono risorse sufficienti per creare un altro
    processo (per allocare la tabella delle pagine e le strutture del task) o
    si è esaurito il numero di processi disponibili.
  \item[\errcode{ENOMEM}] non è stato possibile allocare la memoria per le
    strutture necessarie al kernel per creare il nuovo processo.
  \end{errlist}}
\end{funcproto}

Dopo il successo dell'esecuzione di una \func{fork} sia il processo padre che
il processo figlio continuano ad essere eseguiti normalmente, a partire
dall'istruzione successiva alla \func{fork}. Il processo figlio è una copia
del padre, e riceve una copia dei segmenti di testo, dati e dello
\textit{stack} (vedi sez.~\ref{sec:proc_mem_layout}), ed esegue esattamente lo
stesso codice del padre. Si tenga presente però che la memoria è copiata e non
condivisa, pertanto padre e figlio vedranno variabili diverse e le eventuali
modifiche saranno totalmente indipendenti.

\itindbeg{copy~on~write}

Per quanto riguarda la gestione della memoria, in generale il segmento di
testo, che è identico per i due processi, è condiviso e tenuto in sola lettura
per il padre e per i figli. Per gli altri segmenti Linux utilizza la tecnica
del \textit{copy on write}. Questa tecnica comporta che una pagina di memoria
viene effettivamente copiata per il nuovo processo solo quando ci viene
effettuata sopra una scrittura, e si ha quindi una reale differenza fra padre
e figlio.  In questo modo si rende molto più efficiente il meccanismo della
creazione di un nuovo processo, non essendo più necessaria la copia di tutto
lo spazio degli indirizzi virtuali del padre, ma solo delle pagine di memoria
che sono state modificate, e solo al momento della modifica stessa.

\itindend{copy~on~write}

La differenza che si ha nei due processi è che nel processo padre il valore di
ritorno della funzione \func{fork} è il \ids{PID} del processo figlio, mentre
nel figlio è zero; in questo modo il programma può identificare se viene
eseguito dal padre o dal figlio.  Si noti come la funzione \func{fork} ritorni
due volte, una nel padre e una nel figlio.

La scelta di questi valori di ritorno non è casuale, un processo infatti può
avere più figli, ed il valore di ritorno di \func{fork} è l'unico che gli
permette di identificare qual è quello appena creato. Al contrario un figlio
ha sempre un solo padre il cui \ids{PID}, come spiegato in
sez.~\ref{sec:proc_pid}, può sempre essere ottenuto con \func{getppid}; per
questo si ritorna un valore nullo, che non è il \ids{PID} di nessun processo.

Normalmente la chiamata a \func{fork} può fallire solo per due ragioni: o ci
sono già troppi processi nel sistema, il che di solito è sintomo che
qualcos'altro non sta andando per il verso giusto, o si è ecceduto il limite
sul numero totale di processi permessi all'utente, argomento che tratteremo in
dettaglio in sez.~\ref{sec:sys_resource_limit}.

L'uso di \func{fork} avviene secondo due modalità principali; la prima è
quella in cui all'interno di un programma si creano processi figli cui viene
affidata l'esecuzione di una certa sezione di codice, mentre il processo padre
ne esegue un'altra. È il caso tipico dei programmi server (il modello
\textit{client-server} è illustrato in sez.~\ref{sec:net_cliserv}) in cui il
padre riceve ed accetta le richieste da parte dei programmi client, per
ciascuna delle quali pone in esecuzione un figlio che è incaricato di fornire
le risposte associate al servizio.

La seconda modalità è quella in cui il processo vuole eseguire un altro
programma; questo è ad esempio il caso della shell. In questo caso il processo
crea un figlio la cui unica operazione è quella di fare una \func{exec} (di
cui parleremo in sez.~\ref{sec:proc_exec}) subito dopo la \func{fork}.

Alcuni sistemi operativi (il VMS ad esempio) combinano le operazioni di questa
seconda modalità (una \func{fork} seguita da una \func{exec}) in un'unica
operazione che viene chiamata \textit{spawn}. Nei sistemi unix-like è stato
scelto di mantenere questa separazione, dato che, come per la prima modalità
d'uso, esistono numerosi scenari in cui si può usare una \func{fork} senza
aver bisogno di eseguire una \func{exec}. 

Inoltre, anche nel caso della seconda modalità d'uso, avere le due funzioni
separate permette al figlio di cambiare alcune caratteristiche del processo
(maschera dei segnali, redirezione dell'output, utente per conto del cui viene
eseguito, e molto altro su cui torneremo in seguito) prima della \func{exec},
rendendo così relativamente facile intervenire sulle le modalità di esecuzione
del nuovo programma.

\begin{figure}[!htb]
  \footnotesize \centering
  \begin{minipage}[c]{\codesamplewidth}
  \includecodesample{listati/fork_test.c}
  \end{minipage}
  \normalsize
  \caption{Esempio di codice per la creazione di nuovi processi (da
    \file{fork\_test.c}).}
  \label{fig:proc_fork_code}
\end{figure}

In fig.~\ref{fig:proc_fork_code} è riportato il corpo del codice del programma
di esempio \cmd{forktest}, che permette di illustrare molte caratteristiche
dell'uso della funzione \func{fork}. Il programma crea un numero di figli
specificato da linea di comando, e prende anche alcune opzioni per indicare
degli eventuali tempi di attesa in secondi (eseguiti tramite la funzione
\func{sleep}) per il padre ed il figlio (con \cmd{forktest -h} si ottiene la
descrizione delle opzioni). Il codice completo, compresa la parte che gestisce
le opzioni a riga di comando, è disponibile nel file \file{fork\_test.c},
distribuito insieme agli altri sorgenti degli esempi della guida su
\url{http://gapil.gnulinux.it}.

Decifrato il numero di figli da creare, il ciclo principale del programma
(\texttt{\small 24-40}) esegue in successione la creazione dei processi figli
controllando il successo della chiamata a \func{fork} (\texttt{\small
  25-29}); ciascun figlio (\texttt{\small 31-34}) si limita a stampare il
suo numero di successione, eventualmente attendere il numero di secondi
specificato e scrivere un messaggio prima di uscire. Il processo padre invece
(\texttt{\small 36-38}) stampa un messaggio di creazione, eventualmente
attende il numero di secondi specificato, e procede nell'esecuzione del ciclo;
alla conclusione del ciclo, prima di uscire, può essere specificato un altro
periodo di attesa.

Se eseguiamo il comando, che è preceduto dall'istruzione \code{export
  LD\_LIBRARY\_PATH=./} per permettere l'uso delle librerie dinamiche, senza
specificare attese (come si può notare in (\texttt{\small 17-19}) i valori
predefiniti specificano di non attendere), otterremo come risultato sul
terminale:
\begin{Console}
[piccardi@selidor sources]$ \textbf{export LD_LIBRARY_PATH=./; ./forktest 3}
Process 1963: forking 3 child
Spawned 1 child, pid 1964 
Child 1 successfully executing
Child 1, parent 1963, exiting
Go to next child 
Spawned 2 child, pid 1965 
Child 2 successfully executing
Child 2, parent 1963, exiting
Go to next child 
Child 3 successfully executing
Child 3, parent 1963, exiting
Spawned 3 child, pid 1966 
Go to next child 
\end{Console}
%$

Esaminiamo questo risultato: una prima conclusione che si può trarre è che non
si può dire quale processo fra il padre ed il figlio venga eseguito per primo
dopo la chiamata a \func{fork}; dall'esempio si può notare infatti come nei
primi due cicli sia stato eseguito per primo il padre (con la stampa del
\ids{PID} del nuovo processo) per poi passare all'esecuzione del figlio
(completata con i due avvisi di esecuzione ed uscita), e tornare
all'esecuzione del padre (con la stampa del passaggio al ciclo successivo),
mentre la terza volta è stato prima eseguito il figlio (fino alla conclusione)
e poi il padre.

In generale l'ordine di esecuzione dipenderà, oltre che dall'algoritmo di
\textit{scheduling} usato dal kernel, dalla particolare situazione in cui si
trova la macchina al momento della chiamata, risultando del tutto
impredicibile.  Eseguendo più volte il programma di prova e producendo un
numero diverso di figli, si sono ottenute situazioni completamente diverse,
compreso il caso in cui il processo padre ha eseguito più di una \func{fork}
prima che uno dei figli venisse messo in esecuzione.

Pertanto non si può fare nessuna assunzione sulla sequenza di esecuzione delle
istruzioni del codice fra padre e figli, né sull'ordine in cui questi potranno
essere messi in esecuzione. Se è necessaria una qualche forma di precedenza
occorrerà provvedere ad espliciti meccanismi di sincronizzazione, pena il
rischio di incorrere nelle cosiddette \textit{race condition} (vedi
sez.~\ref{sec:proc_race_cond}).

In realtà con l'introduzione dei kernel della serie 2.6 lo \textit{scheduler}
è stato modificato per eseguire sempre per primo il figlio.\footnote{i
  risultati precedenti infatti sono stati ottenuti usando un kernel della
  serie 2.4.}  Questa è una ottimizzazione adottata per evitare che il padre,
effettuando per primo una operazione di scrittura in memoria, attivasse il
meccanismo del \textit{copy on write}, operazione inutile quando il figlio
viene creato solo per eseguire una \func{exec} per lanciare un altro programma
che scarta completamente lo spazio degli indirizzi e rende superflua la copia
della memoria modificata dal padre. Eseguendo sempre per primo il figlio la
\func{exec} verrebbe effettuata subito, con la certezza di utilizzare il
\textit{copy on write} solo quando necessario.

Con il kernel 2.6.32 però il comportamento è stato nuovamente cambiato,
stavolta facendo eseguire per primo sempre il padre. Si è realizzato infatti
che l'eventualità prospettata per la scelta precedente era comunque poco
probabile, mentre l'esecuzione immediata del padre presenta sempre il
vantaggio di poter utilizzare immediatamente tutti i dati che sono nella cache
della CPU e nell'unità di gestione della memoria virtuale, senza doverli
invalidare, cosa che per i processori moderni, che hanno linee di cache
interne molto profonde, avrebbe un forte impatto sulle prestazioni.

Allora anche se quanto detto in precedenza si verifica nel comportamento
effettivo dei programmi soltanto per i kernel fino alla serie 2.4, per
mantenere la portabilità con altri kernel unix-like e con i diversi
comportamenti adottati dalle Linux nella sua evoluzione, è comunque opportuno
non fare nessuna assunzione sull'ordine di esecuzione di padre e figlio dopo
la chiamata a \func{fork}.

Si noti infine come dopo la \func{fork}, essendo i segmenti di memoria
utilizzati dai singoli processi completamente indipendenti, le modifiche delle
variabili nei processi figli, come l'incremento di \var{i} in (\texttt{\small
  31}), sono visibili solo a loro, (ogni processo vede solo la propria copia
della memoria), e non hanno alcun effetto sul valore che le stesse variabili
hanno nel processo padre ed in eventuali altri processi figli che eseguano lo
stesso codice.

Un secondo aspetto molto importante nella creazione dei processi figli è
quello dell'interazione dei vari processi con i file. Ne parleremo qui anche
se buona parte dei concetti relativi ai file verranno trattati più avanti
(principalmente in sez.~\ref{sec:file_unix_interface}). Per illustrare meglio
quello che avviene si può redirigere su un file l'output del programma di
test, quello che otterremo è:
\begin{Console}
[piccardi@selidor sources]$ \textbf{./forktest 3 > output}
[piccardi@selidor sources]$ \textbf{cat output}
Process 1967: forking 3 child
Child 1 successfully executing
Child 1, parent 1967, exiting
Test for forking 3 child
Spawned 1 child, pid 1968 
Go to next child 
Child 2 successfully executing
Child 2, parent 1967, exiting
Test for forking 3 child
Spawned 1 child, pid 1968 
Go to next child 
Spawned 2 child, pid 1969 
Go to next child 
Child 3 successfully executing
Child 3, parent 1967, exiting
Test for forking 3 child
Spawned 1 child, pid 1968 
Go to next child 
Spawned 2 child, pid 1969 
Go to next child 
Spawned 3 child, pid 1970 
Go to next child 
\end{Console}
che come si vede è completamente diverso da quanto ottenevamo sul terminale.

Il comportamento delle varie funzioni di interfaccia con i file è analizzato
in gran dettaglio in sez.~\ref{sec:file_unix_interface} per l'interfaccia
nativa Unix ed in sez.~\ref{sec:files_std_interface} per la standardizzazione
adottata nelle librerie del linguaggio C, valida per qualunque sistema
operativo.

Qui basta accennare che si sono usate le funzioni standard della libreria del
C che prevedono l'output bufferizzato. Il punto è che questa bufferizzazione
(che tratteremo in dettaglio in sez.~\ref{sec:file_buffering}) varia a seconda
che si tratti di un file su disco, in cui il buffer viene scaricato su disco
solo quando necessario, o di un terminale, in cui il buffer viene scaricato ad
ogni carattere di ``a capo''.

Nel primo esempio allora avevamo che, essendovi un a capo nella stringa
stampata, ad ogni chiamata a \func{printf} il buffer veniva scaricato, per cui
le singole righe comparivano a video subito dopo l'esecuzione della
\func{printf}. Ma con la redirezione su file la scrittura non avviene più alla
fine di ogni riga e l'output resta nel buffer. 

Dato che ogni figlio riceve una copia della memoria del padre, esso riceverà
anche quanto c'è nel buffer delle funzioni di I/O, comprese le linee scritte
dal padre fino allora. Così quando il buffer viene scritto su disco all'uscita
del figlio, troveremo nel file anche tutto quello che il processo padre aveva
scritto prima della sua creazione. E alla fine del file (dato che in questo
caso il padre esce per ultimo) troveremo anche l'output completo del padre.

L'esempio ci mostra un altro aspetto fondamentale dell'interazione con i file,
valido anche per l'esempio precedente, ma meno evidente: il fatto cioè che non
solo processi diversi possono scrivere in contemporanea sullo stesso file
(l'argomento dell'accesso concorrente ai file è trattato in dettaglio in
sez.~\ref{sec:file_shared_access}), ma anche che, a differenza di quanto
avviene per le variabili in memoria, la posizione corrente sul file è
condivisa fra il padre e tutti i processi figli.

Quello che succede è che quando lo \textit{standard output}\footnote{si chiama
  così il file su cui di default un programma scrive i suoi dati in uscita,
  tratteremo l'argomento in dettaglio in sez.~\ref{sec:file_fd}.} del padre
viene rediretto come si è fatto nell'esempio, lo stesso avviene anche per
tutti i figli. La funzione \func{fork} infatti ha la caratteristica di
duplicare nei processi figli tutti i \textit{file descriptor} (vedi
sez.~\ref{sec:file_fd}) dei file aperti nel processo padre (allo stesso modo
in cui lo fa la funzione \func{dup}, trattata in sez.~\ref{sec:file_dup}). Ciò
fa sì che padre e figli condividano le stesse voci della \textit{file table}
(tratteremo in dettaglio questi termini in sez.~\ref{sec:file_fd} e
sez.~\ref{sec:file_shared_access}) fra le quali c'è anche la posizione
corrente nel file.

Quando un processo scrive su un file la posizione corrente viene aggiornata
sulla \textit{file table}, e tutti gli altri processi, che vedono la stessa
\textit{file table}, vedranno il nuovo valore. In questo modo si evita, in
casi come quello appena mostrato in cui diversi figli scrivono sullo stesso
file usato dal padre, che una scrittura eseguita in un secondo tempo da un
processo vada a sovrapporsi a quelle precedenti: l'output potrà risultare
mescolato, ma non ci saranno parti perdute per via di una sovrascrittura.

Questo tipo di comportamento è essenziale in tutti quei casi in cui il padre
crea un figlio e attende la sua conclusione per proseguire, ed entrambi
scrivono sullo stesso file. Un caso tipico di questo comportamento è la shell
quando lancia un programma.  In questo modo, anche se lo standard output viene
rediretto, il padre potrà sempre continuare a scrivere in coda a quanto
scritto dal figlio in maniera automatica; se così non fosse ottenere questo
comportamento sarebbe estremamente complesso, necessitando di una qualche
forma di comunicazione fra i due processi per far riprendere al padre la
scrittura al punto giusto.

In generale comunque non è buona norma far scrivere più processi sullo stesso
file senza una qualche forma di sincronizzazione in quanto, come visto anche
con il nostro esempio, le varie scritture risulteranno mescolate fra loro in
una sequenza impredicibile. Per questo le modalità con cui in genere si usano
i file dopo una \func{fork} sono sostanzialmente due:
\begin{enumerate*}
\item Il processo padre aspetta la conclusione del figlio. In questo caso non
  è necessaria nessuna azione riguardo ai file, in quanto la sincronizzazione
  della posizione corrente dopo eventuali operazioni di lettura e scrittura
  effettuate dal figlio è automatica.
\item L'esecuzione di padre e figlio procede indipendentemente. In questo caso
  ciascuno dei due processi deve chiudere i file che non gli servono una volta
  che la \func{fork} è stata eseguita, per evitare ogni forma di interferenza.
\end{enumerate*}

Oltre ai file aperti i processi figli ereditano dal padre una serie di altre
proprietà; la lista dettagliata delle proprietà che padre e figlio hanno in
comune dopo l'esecuzione di una \func{fork} è la seguente:
\begin{itemize*}
\item i file aperti e gli eventuali flag di \textit{close-on-exec} impostati
  (vedi sez.~\ref{sec:proc_exec} e sez.~\ref{sec:file_fcntl_ioctl});
\item gli identificatori per il controllo di accesso: l'\textsl{user-ID
    reale}, il \textsl{group-ID reale}, l'\textsl{user-ID effettivo}, il
  \textsl{group-ID effettivo} ed i \textsl{group-ID supplementari} (vedi
  sez.~\ref{sec:proc_access_id});
\item gli identificatori per il controllo di sessione: il \textit{process
    group-ID} e il \textit{session id} ed il terminale di controllo (vedi
  sez.~\ref{sec:sess_proc_group});
\item la directory di lavoro (vedi sez.~\ref{sec:file_work_dir}) e la
  directory radice (vedi sez.~\ref{sec:file_chroot});
\item la maschera dei permessi di creazione dei file (vedi
  sez.~\ref{sec:file_perm_management});
\item la maschera dei segnali bloccati (vedi
  sez.~\ref{sec:sig_sigmask}) e le azioni installate (vedi
  sez.~\ref{sec:sig_gen_beha});
\item i segmenti di memoria condivisa agganciati al processo (vedi
  sez.~\ref{sec:ipc_sysv_shm});
\item i limiti sulle risorse (vedi sez.~\ref{sec:sys_resource_limit});
\item il valori di \textit{nice}, le priorità \textit{real-time} e le affinità
  di processore (vedi sez.~\ref{sec:proc_sched_stand},
  sez.~\ref{sec:proc_real_time} e sez.~\ref{sec:proc_sched_multiprocess});
\item le variabili di ambiente (vedi sez.~\ref{sec:proc_environ}).
\item l'insieme dei descrittori associati alle code di messaggi POSIX (vedi
  sez.~\ref{sec:ipc_posix_mq}) che vengono copiate come i \textit{file
    descriptor}, questo significa che entrambi condivideranno gli stessi flag.
\end{itemize*}

Oltre a quelle relative ad un diverso spazio degli indirizzi (e una memoria
totalmente indipendente) le differenze fra padre e figlio dopo l'esecuzione di
una \func{fork} invece sono:\footnote{a parte le ultime quattro, relative a
  funzionalità specifiche di Linux, le altre sono esplicitamente menzionate
  dallo standard POSIX.1-2001.}
\begin{itemize*}
\item il valore di ritorno di \func{fork};
\item il \ids{PID} (\textit{process id}), quello del figlio viene assegnato ad
  un nuovo valore univoco;
\item il \ids{PPID} (\textit{parent process id}), quello del figlio viene
  impostato al \ids{PID} del padre;
\item i valori dei tempi di esecuzione (vedi sez.~\ref{sec:sys_cpu_times}) e
  delle risorse usate (vedi sez.~\ref{sec:sys_resource_use}), che nel figlio
  sono posti a zero;
\item i \textit{lock} sui file (vedi sez.~\ref{sec:file_locking}) e sulla
  memoria (vedi sez.~\ref{sec:proc_mem_lock}), che non vengono ereditati dal
  figlio;
\item gli allarmi, i timer (vedi sez.~\ref{sec:sig_alarm_abort}) ed i segnali
  pendenti (vedi sez.~\ref{sec:sig_gen_beha}), che per il figlio vengono
  cancellati.
\item le operazioni di I/O asincrono in corso (vedi
  sez.~\ref{sec:file_asyncronous_io}) che non vengono ereditate dal figlio;
\item gli aggiustamenti fatti dal padre ai semafori con \func{semop} (vedi
  sez.~\ref{sec:ipc_sysv_sem}).
\item le notifiche sui cambiamenti delle directory con \textit{dnotify} (vedi
  sez.~\ref{sec:sig_notification}), che non vengono ereditate dal figlio;
\item le mappature di memoria marcate come \const{MADV\_DONTFORK} (vedi
  sez.~\ref{sec:file_memory_map}) che non vengono ereditate dal figlio;
\item l'impostazione con \func{prctl} (vedi sez.~\ref{sec:process_prctl}) che
  notifica al figlio la terminazione del padre viene cancellata se presente
  nel padre;
\item il segnale di terminazione del figlio è sempre \signal{SIGCHLD} anche
  qualora nel padre fosse stato modificato (vedi sez.~\ref{sec:process_clone}). 
\end{itemize*}

Una seconda funzione storica usata per la creazione di un nuovo processo è
\funcm{vfork}, che è esattamente identica a \func{fork} ed ha la stessa
semantica e gli stessi errori; la sola differenza è che non viene creata la
tabella delle pagine né la struttura dei task per il nuovo processo. Il
processo padre è posto in attesa fintanto che il figlio non ha eseguito una
\func{execve} o non è uscito con una \func{\_exit}. Il figlio condivide la
memoria del padre (e modifiche possono avere effetti imprevedibili) e non deve
ritornare o uscire con \func{exit} ma usare esplicitamente \func{\_exit}.

Questa funzione è un rimasuglio dei vecchi tempi in cui eseguire una
\func{fork} comportava anche la copia completa del segmento dati del processo
padre, che costituiva un inutile appesantimento in tutti quei casi in cui la
\func{fork} veniva fatta solo per poi eseguire una \func{exec}. La funzione
venne introdotta in BSD per migliorare le prestazioni.

Dato che Linux supporta il \textit{copy on write} la perdita di prestazioni è
assolutamente trascurabile, e l'uso di questa funzione, che resta un caso
speciale della \textit{system call} \func{clone} (che tratteremo in dettaglio
in sez.~\ref{sec:process_clone}) è deprecato; per questo eviteremo di
trattarla ulteriormente.


\subsection{La conclusione di un processo}
\label{sec:proc_termination}

In sez.~\ref{sec:proc_conclusion} abbiamo già affrontato le modalità con cui
chiudere un programma, ma dall'interno del programma stesso. Avendo a che fare
con un sistema \textit{multitasking} resta da affrontare l'argomento dal punto
di vista di come il sistema gestisce la conclusione dei processi.

Abbiamo visto in sez.~\ref{sec:proc_conclusion} le tre modalità con cui un
programma viene terminato in maniera normale: la chiamata di \func{exit}, che
esegue le funzioni registrate per l'uscita e chiude gli \textit{stream} e poi
esegue \func{\_exit}, il ritorno dalla funzione \code{main} equivalente alla
chiamata di \func{exit}, e la chiamata diretta a \func{\_exit}, che passa
direttamente alle operazioni di terminazione del processo da parte del kernel.

Ma abbiamo accennato che oltre alla conclusione normale esistono anche delle
modalità di conclusione anomala. Queste sono in sostanza due: il programma può
chiamare la funzione \func{abort} (vedi sez.~\ref{sec:sig_alarm_abort}) per
invocare una chiusura anomala, o essere terminato da un segnale (torneremo sui
segnali in cap.~\ref{cha:signals}).  In realtà anche la prima modalità si
riconduce alla seconda, dato che \func{abort} si limita a generare il segnale
\signal{SIGABRT}.

Qualunque sia la modalità di conclusione di un processo, il kernel esegue
comunque una serie di operazioni di terminazione: chiude tutti i file aperti,
rilascia la memoria che stava usando, e così via; l'elenco completo delle
operazioni eseguite alla chiusura di un processo è il seguente:
\begin{itemize*}
\item tutti i \textit{file descriptor} (vedi sez.~\ref{sec:file_fd}) sono
  chiusi;
\item viene memorizzato lo stato di terminazione del processo;
\item ad ogni processo figlio viene assegnato un nuovo padre (in genere
  \cmd{init});
\item viene inviato il segnale \signal{SIGCHLD} al processo padre (vedi
  sez.~\ref{sec:sig_sigchld});
\item se il processo è un leader di sessione ed il suo terminale di controllo
  è quello della sessione viene mandato un segnale di \signal{SIGHUP} a tutti i
  processi del gruppo di \textit{foreground} e il terminale di controllo viene
  disconnesso (vedi sez.~\ref{sec:sess_ctrl_term});
\item se la conclusione di un processo rende orfano un \textit{process
    group} ciascun membro del gruppo viene bloccato, e poi gli vengono
  inviati in successione i segnali \signal{SIGHUP} e \signal{SIGCONT}
  (vedi ancora sez.~\ref{sec:sess_ctrl_term}).
\end{itemize*}

\itindbeg{termination~status} 

Oltre queste operazioni è però necessario poter disporre di un meccanismo
ulteriore che consenta di sapere come la terminazione è avvenuta: dato che in
un sistema unix-like tutto viene gestito attraverso i processi, il meccanismo
scelto consiste nel riportare lo \textsl{stato di terminazione} (il cosiddetto
\textit{termination status}) al processo padre.

Nel caso di conclusione normale, abbiamo visto in
sez.~\ref{sec:proc_conclusion} che lo stato di uscita del processo viene
caratterizzato tramite il valore del cosiddetto \textit{exit status}, cioè il
valore passato come argomento alle funzioni \func{exit} o \func{\_exit} o il
valore di ritorno per \code{main}.  Ma se il processo viene concluso in
maniera anomala il programma non può specificare nessun \textit{exit status},
ed è il kernel che deve generare autonomamente il \textit{termination status}
per indicare le ragioni della conclusione anomala.

Si noti la distinzione fra \textit{exit status} e \textit{termination status}:
quello che contraddistingue lo stato di chiusura del processo e viene
riportato attraverso le funzioni \func{wait} o \func{waitpid} (vedi
sez.~\ref{sec:proc_wait}) è sempre quest'ultimo; in caso di conclusione
normale il kernel usa il primo (nel codice eseguito da \func{\_exit}) per
produrre il secondo.

La scelta di riportare al padre lo stato di terminazione dei figli, pur
essendo l'unica possibile, comporta comunque alcune complicazioni: infatti se
alla sua creazione è scontato che ogni nuovo processo abbia un padre, non è
detto che sia così alla sua conclusione, dato che il padre potrebbe essere già
terminato; si potrebbe avere cioè quello che si chiama un processo
\textsl{orfano}.

Questa complicazione viene superata facendo in modo che il processo orfano
venga \textsl{adottato} da \cmd{init}, o meglio dal processo con \ids{PID}
1,\footnote{anche se, come vedremo in sez.~\ref{sec:process_prctl}, a partire
  dal kernel 3.4 è diventato possibile delegare questo compito anche ad un
  altro processo.} cioè quello lanciato direttamente dal kernel all'avvio, che
sta alla base dell'albero dei processi visto in sez.~\ref{sec:proc_hierarchy}
e che anche per questo motivo ha un ruolo essenziale nel sistema e non può mai
terminare.\footnote{almeno non senza un blocco completo del sistema, in caso
  di terminazione o di non esecuzione di \cmd{init} infatti il kernel si
  blocca con un cosiddetto \textit{kernel panic}, dato che questo è un errore
  fatale.}

Come già accennato quando un processo termina, il kernel controlla se è il
padre di altri processi in esecuzione: in caso positivo allora il \ids{PPID}
di tutti questi processi verrà sostituito dal kernel con il \ids{PID} di
\cmd{init}, cioè con 1. In questo modo ogni processo avrà sempre un padre (nel
caso possiamo parlare di un padre \textsl{adottivo}) cui riportare il suo
stato di terminazione.  

Come verifica di questo comportamento possiamo eseguire il nostro programma
\cmd{forktest} imponendo a ciascun processo figlio due secondi di attesa prima
di uscire, il risultato è:
\begin{Console}
[piccardi@selidor sources]$ \textbf{./forktest -c2 3}
Process 1972: forking 3 child
Spawned 1 child, pid 1973 
Child 1 successfully executing
Go to next child 
Spawned 2 child, pid 1974 
Child 2 successfully executing
Go to next child 
Child 3 successfully executing
Spawned 3 child, pid 1975 
Go to next child 

[piccardi@selidor sources]$ Child 3, parent 1, exiting
Child 2, parent 1, exiting
Child 1, parent 1, exiting
\end{Console}
come si può notare in questo caso il processo padre si conclude prima dei
figli, tornando alla shell, che stampa il prompt sul terminale: circa due
secondi dopo viene stampato a video anche l'output dei tre figli che
terminano, e come si può notare in questo caso, al contrario di quanto visto
in precedenza, essi riportano 1 come \ids{PPID}.

Altrettanto rilevante è il caso in cui il figlio termina prima del padre,
perché non è detto che il padre sia in esecuzione e possa ricevere
immediatamente lo stato di terminazione, quindi il kernel deve comunque
conservare una certa quantità di informazioni riguardo ai processi che sta
terminando.

Questo viene fatto mantenendo attiva la voce nella tabella dei processi, e
memorizzando alcuni dati essenziali, come il \ids{PID}, i tempi di CPU usati
dal processo (vedi sez.~\ref{sec:sys_unix_time}) e lo stato di terminazione,
mentre la memoria in uso ed i file aperti vengono rilasciati immediatamente. 

\itindbeg{zombie}

I processi che sono terminati, ma il cui stato di terminazione non è stato
ancora ricevuto dal padre sono chiamati \textit{zombie}, essi restano presenti
nella tabella dei processi ed in genere possono essere identificati
dall'output di \cmd{ps} per la presenza di una \texttt{Z} nella colonna che ne
indica lo stato (vedi tab.~\ref{tab:proc_proc_states}). Quando il padre
effettuerà la lettura dello stato di terminazione anche questa informazione,
non più necessaria, verrà scartata ed il processo potrà considerarsi
completamente concluso.

Possiamo utilizzare il nostro programma di prova per analizzare anche questa
condizione: lanciamo il comando \cmd{forktest} in \textit{background} (vedi
sez.~\ref{sec:sess_job_control}), indicando al processo padre di aspettare 10
secondi prima di uscire. In questo caso, usando \cmd{ps} sullo stesso
terminale (prima dello scadere dei 10 secondi) otterremo:
\begin{Console}
[piccardi@selidor sources]$ \textbf{ps T}
  PID TTY      STAT   TIME COMMAND
  419 pts/0    S      0:00 bash
  568 pts/0    S      0:00 ./forktest -e10 3
  569 pts/0    Z      0:00 [forktest <defunct>]
  570 pts/0    Z      0:00 [forktest <defunct>]
  571 pts/0    Z      0:00 [forktest <defunct>]
  572 pts/0    R      0:00 ps T
\end{Console}
%$
e come si vede, dato che non si è fatto nulla per riceverne lo stato di
terminazione, i tre processi figli sono ancora presenti pur essendosi
conclusi, con lo stato di \textit{zombie} e l'indicazione che sono terminati
(la scritta \texttt{defunct}).

La possibilità di avere degli \textit{zombie} deve essere tenuta sempre
presente quando si scrive un programma che deve essere mantenuto in esecuzione
a lungo e creare molti processi figli. In questo caso si deve sempre avere
cura di far leggere al programma l'eventuale stato di uscita di tutti i
figli. Una modalità comune di farlo è attraverso l'utilizzo di un apposito
\textit{signal handler} che chiami la funzione \func{wait}, (vedi
sez.~\ref{sec:proc_wait}), ne esamineremo in dettaglio un esempio
(fig.~\ref{fig:sig_sigchld_handl}) in sez.~\ref{sec:sig_sigchld}.

La lettura degli stati di uscita è necessaria perché anche se gli
\textit{zombie} non consumano risorse di memoria o processore, occupano
comunque una voce nella tabella dei processi e se li si lasciano accumulare a
lungo quest'ultima potrebbe esaurirsi, con la conseguente impossibilità di
lanciare nuovi processi.

Si noti tuttavia che quando un processo adottato da \cmd{init} termina, non
diviene mai uno \textit{zombie}. Questo perché una delle funzioni di
\cmd{init} è appunto quella di chiamare la funzione \func{wait} per i processi
a cui fa da padre, completandone la terminazione. Questo è quanto avviene
anche quando, come nel caso del precedente esempio con \cmd{forktest}, il
padre termina con dei figli in stato di \textit{zombie}. Questi scompaiono
quando, alla terminazione del padre dopo i secondi programmati, tutti figli
che avevamo generato, e che erano diventati \textit{zombie}, vengono adottati
da \cmd{init}, il quale provvede a completarne la terminazione.

Si tenga presente infine che siccome gli \textit{zombie} sono processi già
terminati, non c'è modo di eliminarli con il comando \cmd{kill} o inviandogli
un qualunque segnale di terminazione (l'argomento è trattato in
sez.~\ref{sec:sig_termination}). Qualora ci si trovi in questa situazione
l'unica possibilità di cancellarli dalla tabella dei processi è quella di
terminare il processo che li ha generati e che non sta facendo il suo lavoro,
in modo che \cmd{init} possa adottarli e concluderne correttamente la
terminazione.

Si tenga anche presente che la presenza di \textit{zombie} nella tabella dei
processi non è sempre indice di un qualche malfunzionamento, in una macchina
con molto carico infatti può esservi una presenza temporanea dovuta al fatto
che il processo padre ancora non ha avuto il tempo di gestirli. 

\itindend{zombie}


\subsection{Le funzioni di attesa e ricezione degli stati di uscita}
\label{sec:proc_wait}

Uno degli usi più comuni delle capacità \textit{multitasking} di un sistema
unix-like consiste nella creazione di programmi di tipo server, in cui un
processo principale attende le richieste che vengono poi soddisfatte da una
serie di processi figli.

Si è già sottolineato al paragrafo precedente come in questo caso diventi
necessario gestire esplicitamente la conclusione dei figli onde evitare di
riempire di \textit{zombie} la tabella dei processi. Tratteremo in questa
sezione le funzioni di sistema deputate a questo compito; la prima è
\funcd{wait} ed il suo prototipo è:

\begin{funcproto}{ 
\fhead{sys/types.h}
\fhead{sys/wait.h}
\fdecl{pid\_t wait(int *status)}
\fdesc{Attende la terminazione di un processo.} 
}
{La funzione ritorna il \ids{PID} del figlio in caso di successo e $-1$ per un
  errore, nel qual caso \var{errno} assumerà uno dei valori:
  \begin{errlist}
  \item[\errcode{ECHILD}] il processo non ha nessun figlio di cui attendere
    uno stato di terminazione.
  \item[\errcode{EINTR}] la funzione è stata interrotta da un segnale.
  \end{errlist}}
\end{funcproto}

Questa funzione è presente fin dalle prime versioni di Unix ed è quella usata
tradizionalmente per attendere la terminazione dei figli. La funzione sospende
l'esecuzione del processo corrente e ritorna non appena un qualunque processo
figlio termina. Se un figlio è già terminato prima della sua chiamata la
funzione ritorna immediatamente, se più processi figli sono già terminati
occorrerà continuare a chiamare la funzione più volte fintanto che non si è
recuperato lo stato di terminazione di tutti quanti.

Al ritorno della funzione lo stato di terminazione del figlio viene salvato
(come \textit{value result argument}) nella variabile puntata
da \param{status} e tutte le risorse del kernel relative al processo (vedi
sez.~\ref{sec:proc_termination}) vengono rilasciate.  Nel caso un processo
abbia più figli il valore di ritorno della funzione sarà impostato al
\ids{PID} del processo di cui si è ricevuto lo stato di terminazione, cosa che
permette di identificare qual è il figlio che è terminato.

\itindend{termination~status}

Questa funzione ha il difetto di essere poco flessibile, in quanto ritorna
all'uscita di un qualunque processo figlio. Nelle occasioni in cui è
necessario attendere la conclusione di uno specifico processo fra tutti quelli
esistenti occorre predisporre un meccanismo che tenga conto di tutti processi
che sono terminati, e provveda a ripetere la chiamata alla funzione nel caso
il processo cercato non risulti fra questi. Se infatti il processo cercato è
già terminato e se è già ricevuto lo stato di uscita senza registrarlo, la
funzione non ha modo di accorgersene, e si continuerà a chiamarla senza
accorgersi che quanto interessava è già accaduto.

Per questo motivo lo standard POSIX.1 ha introdotto una seconda funzione che
effettua lo stesso servizio, ma dispone di una serie di funzionalità più
ampie, legate anche al controllo di sessione (si veda
sez.~\ref{sec:sess_job_control}).  Dato che è possibile ottenere lo stesso
comportamento di \func{wait}\footnote{in effetti il codice
  \code{wait(\&status)} è del tutto equivalente a \code{waitpid(WAIT\_ANY,
    \&status, 0)}.} si consiglia di utilizzare sempre questa nuova funzione di
sistema, \funcd{waitpid}, il cui prototipo è:

\begin{funcproto}{ 
\fhead{sys/types.h}
\fhead{sys/wait.h}
\fdecl{pid\_t waitpid(pid\_t pid, int *status, int options)}
\fdesc{Attende il cambiamento di stato di un processo figlio.} 
}
{La funzione ritorna il \ids{PID} del processo che ha cambiato stato in caso
  di successo, o 0 se è stata specificata l'opzione \const{WNOHANG} e il
  processo non è uscito e $-1$ per un errore, nel qual caso \var{errno}
  assumerà uno dei valori:
  \begin{errlist}
  \item[\errcode{ECHILD}] il processo specificato da \param{pid} non esiste o
    non è figlio del processo chiamante.
  \item[\errcode{EINTR}] non è stata specificata l'opzione \const{WNOHANG} e
    la funzione è stata interrotta da un segnale.
  \item[\errcode{EINVAL}] si è specificato un valore non valido per
    l'argomento \param{options}.
  \end{errlist}}
\end{funcproto}

La prima differenza fra le due funzioni è che con \func{waitpid} si può
specificare in maniera flessibile quale processo attendere, sulla base del
valore fornito dall'argomento \param{pid}. Questo può assumere diversi valori,
secondo lo specchietto riportato in tab.~\ref{tab:proc_waidpid_pid}, dove si
sono riportate anche le costanti definite per indicare alcuni di essi.

\begin{table}[!htb]
  \centering
  \footnotesize
  \begin{tabular}[c]{|c|c|p{8cm}|}
    \hline
    \textbf{Valore} & \textbf{Costante} &\textbf{Significato}\\
    \hline
    \hline
    $<-1$& --               & Attende per un figlio il cui \textit{process
                              group} (vedi sez.~\ref{sec:sess_proc_group}) è
                              uguale al valore assoluto di \param{pid}.\\ 
    $-1$&\constd{WAIT\_ANY}  & Attende per un figlio qualsiasi, usata in
                               questa maniera senza specificare nessuna opzione
                               è equivalente a \func{wait}.\\ 
    $ 0$&\constd{WAIT\_MYPGRP}&Attende per un figlio il cui \textit{process
                               group} (vedi sez.~\ref{sec:sess_proc_group}) è
                                uguale a quello del processo chiamante.\\ 
    $>0$& --                & Attende per un figlio il cui \ids{PID} è uguale
                              al valore di \param{pid}.\\
    \hline
  \end{tabular}
  \caption{Significato dei valori dell'argomento \param{pid} della funzione
    \func{waitpid}.}
  \label{tab:proc_waidpid_pid}
\end{table}

Il comportamento di \func{waitpid} può inoltre essere modificato passando alla
funzione delle opportune opzioni tramite l'argomento \param{options}; questo
deve essere specificato come maschera binaria delle costanti riportati nella
prima parte in tab.~\ref{tab:proc_waitpid_options} che possono essere
combinate fra loro con un OR aritmetico. Nella seconda parte della stessa
tabella si sono riportati anche alcune opzioni non standard specifiche di
Linux, che consentono un controllo più dettagliato per i processi creati con
la \textit{system call} generica \func{clone} (vedi
sez.~\ref{sec:process_clone}) e che vengono usati principalmente per la
gestione della terminazione dei \textit{thread}\unavref{ (vedi
sez.~\ref{sec:thread_xxx})}.

\begin{table}[!htb]
  \centering
  \footnotesize
  \begin{tabular}[c]{|l|p{8cm}|}
    \hline
    \textbf{Costante} & \textbf{Descrizione}\\
    \hline
    \hline
    \const{WNOHANG}   & La funzione ritorna immediatamente anche se non è
                        terminato nessun processo figlio.\\
    \const{WUNTRACED} & Ritorna anche quando un processo figlio è stato
                        fermato.\\ 
    \const{WCONTINUED}& Ritorna anche quando un processo figlio che era stato
                        fermato ha ripreso l'esecuzione (dal kernel 2.6.10).\\
    \hline
    \constd{\_\_WCLONE}& Attende solo per i figli creati con \func{clone} 
                        (vedi sez.~\ref{sec:process_clone}), vale a dire
                        processi che non emettono nessun segnale 
                        o emettono un segnale diverso da \signal{SIGCHLD} alla
                        terminazione, il default è attendere soltanto i
                        processi figli ordinari ignorando quelli creati da
                        \func{clone}.\\
    \constd{\_\_WALL}  & Attende per qualunque figlio, sia ordinario che creato
                        con \func{clone}, se specificata con
                        \const{\_\_WCLONE} quest'ultima viene ignorata. \\
    \constd{\_\_WNOTHREAD}& Non attende per i figli di altri \textit{thread}
                        dello stesso \textit{thread group}, questo era il
                        comportamento di default del kernel 2.4 che non
                        supportava la possibilità, divenuta il default a
                        partire dal 2.6, di attendere un qualunque figlio
                        appartenente allo stesso \textit{thread group}. \\
    \hline
  \end{tabular}
  \caption{Costanti che identificano i bit dell'argomento \param{options}
    della funzione \func{waitpid}.} 
  \label{tab:proc_waitpid_options}
\end{table}

\constbeg{WNOHANG}

L'uso dell'opzione \const{WNOHANG} consente di prevenire il blocco della
funzione qualora nessun figlio sia uscito o non si siano verificate le altre
condizioni per l'uscita della funzione. in tal caso. In tal caso la funzione,
invece di restituire il \ids{PID} del processo (che è sempre un intero
positivo) ritornerà un valore nullo.

\constend{WNOHANG}
\constbeg{WUNTRACED}
\constbeg{WCONTINUED}

Le altre due opzioni, \const{WUNTRACED} e \const{WCONTINUED}, consentono
rispettivamente di tracciare non la terminazione di un processo, ma il fatto
che esso sia stato fermato, o fatto ripartire, e sono utilizzate per la
gestione del controllo di sessione (vedi sez.~\ref{sec:sess_job_control}).

Nel caso di \const{WUNTRACED} la funzione ritorna, restituendone il \ids{PID},
quando un processo figlio entra nello stato \textit{stopped}\footnote{in
  realtà viene notificato soltanto il caso in cui il processo è stato fermato
  da un segnale di stop (vedi sez.~\ref{sec:sess_ctrl_term}), e non quello in
  cui lo stato \textit{stopped} è dovuto all'uso di \func{ptrace}\unavref{
    (vedi sez.~\ref{sec:process_ptrace})}.} (vedi
tab.~\ref{tab:proc_proc_states}), mentre con \const{WCONTINUED} la funzione
ritorna quando un processo in stato \textit{stopped} riprende l'esecuzione per
la ricezione del segnale \signal{SIGCONT} (l'uso di questi segnali per il
controllo di sessione è trattato in sez.~\ref{sec:sess_ctrl_term}).

\constend{WUNTRACED}
\constend{WCONTINUED}

La terminazione di un processo figlio (così come gli altri eventi osservabili
con \func{waitpid}) è chiaramente un evento asincrono rispetto all'esecuzione
di un programma e può avvenire in un qualunque momento. Per questo motivo,
come accennato nella sezione precedente, una delle azioni prese dal kernel
alla conclusione di un processo è quella di mandare un segnale di
\signal{SIGCHLD} al padre. L'azione predefinita per questo segnale (si veda
sez.~\ref{sec:sig_base}) è di essere ignorato, ma la sua generazione
costituisce il meccanismo di comunicazione asincrona con cui il kernel avverte
il processo padre che uno dei suoi figli è terminato.

Il comportamento delle funzioni è però cambiato nel passaggio dal kernel 2.4
al kernel 2.6, quest'ultimo infatti si è adeguato alle prescrizioni dello
standard POSIX.1-2001 e come da esso richiesto se \signal{SIGCHLD} viene
ignorato, o se si imposta il flag di \const{SA\_NOCLDSTOP} nella ricezione
dello stesso (si veda sez.~\ref{sec:sig_sigaction}) i processi figli che
terminano non diventano \textit{zombie} e sia \func{wait} che \func{waitpid}
si bloccano fintanto che tutti i processi figli non sono terminati, dopo di
che falliscono con un errore di \errcode{ENOCHLD}.\footnote{questo è anche il
  motivo per cui le opzioni \const{WUNTRACED} e \const{WCONTINUED} sono
  utilizzabili soltanto qualora non si sia impostato il flag di
  \const{SA\_NOCLDSTOP} per il segnale \signal{SIGCHLD}.}

Con i kernel della serie 2.4 e tutti i kernel delle serie precedenti entrambe
le funzioni di attesa ignorano questa prescrizione e si comportano sempre
nello stesso modo,\footnote{lo standard POSIX.1 originale infatti lascia
  indefinito il comportamento di queste funzioni quando \signal{SIGCHLD} viene
  ignorato.} indipendentemente dal fatto \signal{SIGCHLD} sia ignorato o meno:
attendono la terminazione di un processo figlio e ritornano il relativo
\ids{PID} e lo stato di terminazione nell'argomento \param{status}.

In generale in un programma non si vuole essere forzati ad attendere la
conclusione di un processo figlio per proseguire l'esecuzione, specie se tutto
questo serve solo per leggerne lo stato di chiusura (ed evitare eventualmente
la presenza di \textit{zombie}).  Per questo la modalità più comune di
chiamare queste funzioni è quella di utilizzarle all'interno di un
\textit{signal handler} (vedremo un esempio di come gestire \signal{SIGCHLD}
con i segnali in sez.~\ref{sec:sig_example}). In questo caso infatti, dato che
il segnale è generato dalla terminazione di un figlio, avremo la certezza che
la chiamata a \func{waitpid} non si bloccherà.

Come accennato sia \func{wait} che \func{waitpid} restituiscono lo stato di
terminazione del processo tramite il puntatore \param{status}, e se non
interessa memorizzarlo si può passare un puntatore nullo. Il valore restituito
da entrambe le funzioni dipende dall'implementazione, ma tradizionalmente gli
8 bit meno significativi sono riservati per memorizzare lo stato di uscita del
processo, e gli altri per indicare il segnale che ha causato la terminazione
(in caso di conclusione anomala), uno per indicare se è stato generato un
\textit{core dump} (vedi sez.~\ref{sec:sig_standard}), ecc.\footnote{le
  definizioni esatte si possono trovare in \file{<bits/waitstatus.h>} ma
  questo file non deve mai essere usato direttamente, esso viene incluso
  attraverso \file{<sys/wait.h>}.}

\begin{table}[!htb]
  \centering
  \footnotesize
  \begin{tabular}[c]{|l|p{10cm}|}
    \hline
    \textbf{Macro} & \textbf{Descrizione}\\
    \hline
    \hline
    \macrod{WIFEXITED}\texttt{(s)}   & Condizione vera (valore non nullo) per
                                      un processo figlio che sia terminato
                                      normalmente. \\ 
    \macrod{WEXITSTATUS}\texttt{(s)} & Restituisce gli otto bit meno
                                      significativi dello stato di uscita del
                                      processo (passato attraverso
                                      \func{\_exit}, \func{exit} o come valore
                                      di ritorno di \code{main}); può essere
                                      valutata solo se \val{WIFEXITED} ha
                                      restituito un valore non nullo.\\ 
    \macrod{WIFSIGNALED}\texttt{(s)} & Condizione vera se il processo figlio è
                                      terminato in maniera anomala a causa di
                                      un segnale che non è stato catturato
                                      (vedi sez.~\ref{sec:sig_notification}).\\ 
    \macrod{WTERMSIG}\texttt{(s)}    & Restituisce il numero del segnale che ha
                                      causato la terminazione anomala del
                                      processo; può essere valutata solo se
                                      \val{WIFSIGNALED} ha restituito un
                                      valore non nullo.\\
    \macrod{WCOREDUMP}\texttt{(s)}   & Vera se il processo terminato ha
                                      generato un file di 
                                      \textit{core dump}; può essere valutata
                                      solo se \val{WIFSIGNALED} ha restituito
                                      un valore non nullo.\footnotemark \\
    \macrod{WIFSTOPPED}\texttt{(s)}  & Vera se il processo che ha causato il
                                      ritorno di \func{waitpid} è bloccato;
                                      l'uso è possibile solo con
                                      \func{waitpid} avendo specificato
                                      l'opzione \const{WUNTRACED}.\\
    \macrod{WSTOPSIG}\texttt{(s)}    & Restituisce il numero del segnale che ha
                                      bloccato il processo; può essere
                                      valutata solo se \val{WIFSTOPPED} ha
                                      restituito un valore non nullo. \\ 
    \macrod{WIFCONTINUED}\texttt{(s)}& Vera se il processo che ha causato il
                                      ritorno è stato riavviato da un
                                      \signal{SIGCONT} (disponibile solo a
                                      partire dal kernel 2.6.10).\\
    \hline
  \end{tabular}
  \caption{Descrizione delle varie macro di preprocessore utilizzabili per 
    verificare lo stato di terminazione \var{s} di un processo.}
  \label{tab:proc_status_macro}
\end{table}

\footnotetext{questa macro non è definita dallo standard POSIX.1-2001, ma è
  presente come estensione sia in Linux che in altri Unix, deve essere
  pertanto utilizzata con attenzione (ad esempio è il caso di usarla in un
  blocco \texttt{\#ifdef WCOREDUMP ... \#endif}.}

Lo standard POSIX.1 definisce una serie di macro di preprocessore da usare per
analizzare lo stato di uscita. Esse sono definite sempre in
\file{<sys/wait.h>} ed elencate in tab.~\ref{tab:proc_status_macro}. Si tenga
presente che queste macro prevedono che gli si passi come parametro la
variabile di tipo \ctyp{int} puntata dall'argomento \param{status} restituito
da \func{wait} o \func{waitpid}.

Si tenga conto che nel caso di conclusione anomala il valore restituito da
\val{WTERMSIG} può essere confrontato con le costanti che identificano i
segnali definite in \headfile{signal.h} ed elencate in
tab.~\ref{tab:sig_signal_list}, e stampato usando le apposite funzioni
trattate in sez.~\ref{sec:sig_strsignal}.

A partire dal kernel 2.6.9, sempre in conformità allo standard POSIX.1-2001, è
stata introdotta una nuova funzione di attesa che consente di avere un
controllo molto più preciso sui possibili cambiamenti di stato dei processi
figli e più dettagli sullo stato di uscita; la funzione di sistema è
\funcd{waitid} ed il suo prototipo è:

\begin{funcproto}{ 
\fhead{sys/types.h}
\fhead{sys/wait.h}
\fdecl{int waitid(idtype\_t idtype, id\_t id, siginfo\_t *infop, int options)}
\fdesc{Attende il cambiamento di stato di un processo figlio.} 
}
{La funzione ritorna $0$ in caso di successo e $-1$ per un errore, nel qual
  caso \var{errno} assumerà uno dei valori:
  \begin{errlist}
  \item[\errcode{ECHILD}] il processo specificato da \param{pid} non esiste o
    non è figlio del processo chiamante.
  \item[\errcode{EINTR}] non è stata specificata l'opzione \const{WNOHANG} e
    la funzione è stata interrotta da un segnale.
  \item[\errcode{EINVAL}] si è specificato un valore non valido per
    l'argomento \param{options}.
  \end{errlist}}
\end{funcproto}

La funzione prevede che si specifichi quali processi si intendono osservare
usando i due argomenti \param{idtype} ed \param{id}; il primo indica se ci si
vuole porre in attesa su un singolo processo, un gruppo di processi o un
processo qualsiasi, e deve essere specificato secondo uno dei valori di
tab.~\ref{tab:proc_waitid_idtype}; il secondo indica, a seconda del valore del
primo, quale processo o quale gruppo di processi selezionare.

\begin{table}[!htb]
  \centering
  \footnotesize
  \begin{tabular}[c]{|l|p{8cm}|}
    \hline
    \textbf{Valore} & \textbf{Descrizione}\\
    \hline
    \hline
    \constd{P\_PID} & Indica la richiesta di attendere per un processo figlio
                      il cui \ids{PID} corrisponda al valore dell'argomento
                      \param{id}.\\
    \constd{P\_PGID}& Indica la richiesta di attendere per un processo figlio
                      appartenente al \textit{process group} (vedi
                      sez.~\ref{sec:sess_proc_group}) il cui \acr{pgid}
                      corrisponda al valore dell'argomento \param{id}.\\
    \constd{P\_ALL} & Indica la richiesta di attendere per un processo figlio
                      generico, il valore dell'argomento \param{id} viene
                      ignorato.\\
    \hline
  \end{tabular}
  \caption{Costanti per i valori dell'argomento \param{idtype} della funzione
    \func{waitid}.}
  \label{tab:proc_waitid_idtype}
\end{table}

Come per \func{waitpid} anche il comportamento di \func{waitid} è
controllato dall'argomento \param{options}, da specificare come maschera
binaria dei valori riportati in tab.~\ref{tab:proc_waitid_options}. Benché
alcuni di questi siano identici come significato ed effetto ai precedenti di
tab.~\ref{tab:proc_waitpid_options}, ci sono delle differenze significative:
in questo caso si dovrà specificare esplicitamente l'attesa della terminazione
di un processo impostando l'opzione \const{WEXITED}, mentre il precedente
\const{WUNTRACED} è sostituito da \const{WSTOPPED}.  Infine è stata aggiunta
l'opzione \const{WNOWAIT} che consente una lettura dello stato mantenendo il
processo in attesa di ricezione, così che una successiva chiamata possa di
nuovo riceverne lo stato.

\begin{table}[!htb]
  \centering
  \footnotesize
  \begin{tabular}[c]{|l|p{8cm}|}
    \hline
    \textbf{Valore} & \textbf{Descrizione}\\
    \hline
    \hline
    \constd{WEXITED}   & Ritorna quando un processo figlio è terminato.\\
    \constd{WNOHANG}   & Ritorna immediatamente anche se non c'è niente da
                         notificare.\\ 
    \constd{WSTOPPED}  & Ritorna quando un processo figlio è stato fermato.\\
    \const{WCONTINUED} & Ritorna quando un processo figlio che era stato
                         fermato ha ripreso l'esecuzione.\\
    \constd{WNOWAIT}   & Lascia il processo ancora in attesa di ricezione, così
                         che una successiva chiamata possa di nuovo riceverne
                         lo stato.\\
    \hline
  \end{tabular}
  \caption{Costanti che identificano i bit dell'argomento \param{options}
    della funzione \func{waitid}.} 
  \label{tab:proc_waitid_options}
\end{table}

La funzione \func{waitid} restituisce un valore nullo in caso di successo, e
$-1$ in caso di errore; viene restituito un valore nullo anche se è stata
specificata l'opzione \const{WNOHANG} e la funzione è ritornata immediatamente
senza che nessun figlio sia terminato. Pertanto per verificare il motivo del
ritorno della funzione occorre analizzare le informazioni che essa
restituisce; queste, al contrario delle precedenti \func{wait} e
\func{waitpid} che usavano un semplice valore numerico, sono ritornate in una
struttura di tipo \struct{siginfo\_t} (vedi fig.~\ref{fig:sig_siginfo_t})
all'indirizzo puntato dall'argomento \param{infop}.

Tratteremo nei dettagli la struttura \struct{siginfo\_t} ed il significato dei
suoi vari campi in sez.~\ref{sec:sig_sigaction}, per quanto ci interessa qui
basta dire che al ritorno di \func{waitid} verranno avvalorati i seguenti
campi:
\begin{basedescript}{\desclabelwidth{1.8cm}}
\item[\var{si\_pid}] con il \ids{PID} del figlio.
\item[\var{si\_uid}] con l'\textsl{user-ID reale} (vedi
  sez.~\ref{sec:proc_perms}) del figlio.
\item[\var{si\_signo}] con \signal{SIGCHLD}.
\item[\var{si\_status}] con lo stato di uscita del figlio o con il segnale che
  lo ha terminato, fermato o riavviato.
\item[\var{si\_code}] con uno fra \const{CLD\_EXITED}, \const{CLD\_KILLED},
  \const{CLD\_STOPPED}, \const{CLD\_CONTINUED}, \const{CLD\_TRAPPED} e
  \const{CLD\_DUMPED} a indicare la ragione del ritorno della funzione, il cui
  significato è, nell'ordine: uscita normale, terminazione da segnale,
  processo fermato, processo riavviato, processo terminato in
  \textit{core dump} (vedi sez.~\ref{sec:sig_standard}).
\end{basedescript}

Infine Linux, seguendo un'estensione di BSD, supporta altre due funzioni per
la lettura dello stato di terminazione di un processo, analoghe alle
precedenti ma che prevedono un ulteriore argomento attraverso il quale il
kernel può restituire al padre informazioni sulle risorse (vedi
sez.~\ref{sec:sys_res_limits}) usate dal processo terminato e dai vari figli.
Le due funzioni di sistema sono \funcd{wait3} e \funcd{wait4}, che diventano
accessibili definendo la macro \macro{\_USE\_BSD}, i loro prototipi sono:

\begin{funcproto}{ 
\fhead{sys/types.h}
\fhead{sys/times.h}
\fhead{sys/resource.h}
\fhead{sys/wait.h}
\fdecl{int wait3(int *status, int options, struct rusage *rusage)}
\fdecl{int wait4(pid\_t pid, int *status, int options, struct rusage *rusage)}
\fdesc{Attende il cambiamento di stato di un processo figlio, riportando l'uso
  delle risorse.} 
}
{La funzione ha gli stessi valori di ritorno e codici di errore di
  \func{waitpid}. }
\end{funcproto}

La funzione \func{wait4} è identica \func{waitpid} sia nel comportamento che
per i valori dei primi tre argomenti, ma in più restituisce nell'argomento
aggiuntivo \param{rusage} un sommario delle risorse usate dal processo. Questo
argomento è una struttura di tipo \struct{rusage} definita in
\headfile{sys/resource.h}, che viene utilizzata anche dalla funzione
\func{getrusage} per ottenere le risorse di sistema usate da un processo. La
sua definizione è riportata in fig.~\ref{fig:sys_rusage_struct} e ne
tratteremo in dettaglio il significato sez.~\ref{sec:sys_resource_use}. La
funzione \func{wait3} è semplicemente un caso particolare di (e con Linux
viene realizzata con la stessa \textit{system call}), ed è equivalente a
chiamare \code{wait4(-1, \&status, opt, rusage)}, per questo motivo è ormai
deprecata in favore di \func{wait4}.



\subsection{La famiglia delle funzioni \func{exec} per l'esecuzione dei
  programmi}
\label{sec:proc_exec}

Abbiamo già detto che una delle modalità principali con cui si utilizzano i
processi in Unix è quella di usarli per lanciare nuovi programmi: questo viene
fatto attraverso una delle funzioni della famiglia \func{exec}. Quando un
processo chiama una di queste funzioni esso viene completamente sostituito dal
nuovo programma, il \ids{PID} del processo non cambia, dato che non viene
creato un nuovo processo, la funzione semplicemente rimpiazza lo
\textit{stack}, i dati ed il testo del processo corrente con un nuovo
programma letto da disco, eseguendo il \textit{link-loader} con gli effetti
illustrati in sez.~\ref{sec:proc_main}.

\begin{figure}[!htb]
  \centering \includegraphics[width=8cm]{img/exec_rel}
  \caption{La interrelazione fra le sei funzioni della famiglia \func{exec}.}
  \label{fig:proc_exec_relat}
\end{figure}

Ci sono sei diverse versioni di \func{exec} (per questo la si è chiamata
famiglia di funzioni) che possono essere usate per questo compito, in realtà
(come mostrato in fig.~\ref{fig:proc_exec_relat}), tutte queste funzioni sono
tutte varianti che consentono di invocare in modi diversi, semplificando il
passaggio degli argomenti, la funzione di sistema \funcd{execve}, il cui
prototipo è:

\begin{funcproto}{
\fhead{unistd.h}
\fdecl{int execve(const char *filename, char *const argv[], char *const envp[])}
\fdesc{Esegue un programma.} 
}
{La funzione ritorna solo in caso di errore, restituendo $-1$, nel qual
 caso \var{errno} assumerà uno dei valori:
\begin{errlist}
  \item[\errcode{EACCES}] il file o l'interprete non file ordinari, o non sono
    eseguibili, o il file è su un filesystem montato con l'opzione
    \cmd{noexec}, o manca  il permesso di attraversamento di una delle
    directory del \textit{pathname}.
  \item[\errcode{EAGAIN}] dopo un cambio di \ids{UID} si è ancora  sopra il
    numero massimo di processi consentiti per l'utente (dal kernel 3.1, per i
    dettagli vedi sez.~\ref{sec:proc_setuid}).
  \item[\errcode{EINVAL}] l'eseguibile ELF ha più di un segmento
    \const{PT\_INTERP}, cioè chiede di essere eseguito da più di un
    interprete.
  \item[\errcode{ELIBBAD}] un interprete ELF non è in un formato
    riconoscibile.
  \item[\errcode{ENOENT}] il file o una delle librerie dinamiche o l'interprete
    necessari per eseguirlo non esistono.
  \item[\errcode{ENOEXEC}] il file è in un formato non eseguibile o non
    riconosciuto come tale, o compilato per un'altra architettura.
  \item[\errcode{EPERM}] il file ha i bit \acr{suid} o \acr{sgid} e l'utente
    non è root, ed il processo viene tracciato, oppure il filesystem è montato
    con l'opzione \cmd{nosuid}. 
  \item[\errcode{ETXTBSY}] l'eseguibile è aperto in scrittura da uno o più
    processi. 
  \item[\errcode{E2BIG}] la lista degli argomenti è troppo grande.
  \end{errlist}
  ed inoltre \errval{EFAULT}, \errval{EIO}, \errval{EISDIR}, \errval{ELOOP},
  \errval{EMFILE}, \errval{ENAMETOOLONG}, \errval{ENFILE}, \errval{ENOMEM},
  \errval{ENOTDIR} nel loro significato generico.  }
\end{funcproto}

La funzione \func{execve} esegue il programma o lo script indicato dal
\textit{pathname} \param{filename}, passandogli la lista di argomenti indicata
da \param{argv} e come ambiente la lista di stringhe indicata
da \param{envp}. Entrambe le liste devono essere terminate da un puntatore
nullo. I vettori degli argomenti e dell'ambiente possono essere acceduti dal
nuovo programma quando la sua funzione \code{main} è dichiarata nella forma
\code{main(int argc, char *argv[], char *envp[])}. Si tenga presente per il
passaggio degli argomenti e dell'ambiente esistono comunque dei limiti, su cui
torneremo in sez.~\ref{sec:sys_res_limits}).
% TODO aggiungere la parte sul numero massimo di argomenti, da man execve

In caso di successo la funzione non ritorna, in quanto al posto del programma
chiamante viene eseguito il nuovo programma indicato da \param{filename}. Se
il processo corrente è tracciato con \func{ptrace}\unavref{ (vedi
sez.~\ref{sec:process_ptrace})} in caso di successo viene emesso il segnale
\signal{SIGTRAP}.

Le altre funzioni della famiglia (\funcd{execl}, \funcd{execv},
\funcd{execle}, \funcd{execlp}, \funcd{execvp}) servono per fornire all'utente
una serie di possibili diverse interfacce nelle modalità di passaggio degli
argomenti all'esecuzione del nuovo programma. I loro prototipi sono:

\begin{funcproto}{ 
\fhead{unistd.h}
\fdecl{int execl(const char *path, const char *arg, ...)}
\fdecl{int execv(const char *path, char *const argv[])}
\fdecl{int execle(const char *path, const char *arg, ..., char * const envp[])}
\fdecl{int execlp(const char *file, const char *arg, ...)}
\fdecl{int execvp(const char *file, char *const argv[])}
\fdesc{Eseguono un programma.} 
}
{Le funzioni ritornano solo in caso di errore, restituendo $-1$, i codici di
  errore sono gli stessi di \func{execve}.
}
\end{funcproto}

Tutte le funzioni mettono in esecuzione nel processo corrente il programma
indicati nel primo argomento. Gli argomenti successivi consentono di
specificare gli argomenti e l'ambiente che saranno ricevuti dal nuovo
processo. Per capire meglio le differenze fra le funzioni della famiglia si può
fare riferimento allo specchietto riportato in
tab.~\ref{tab:proc_exec_scheme}. La relazione fra le funzioni è invece
illustrata in fig.~\ref{fig:proc_exec_relat}.

\begin{table}[!htb]
  \footnotesize
  \centering
  \begin{tabular}[c]{|l|c|c|c||c|c|c|}
    \hline
    \multicolumn{1}{|c|}{\textbf{Caratteristiche}} & 
    \multicolumn{6}{|c|}{\textbf{Funzioni}} \\
    \hline
    &\func{execl}\texttt{ }&\func{execlp}&\func{execle}
    &\func{execv}\texttt{ }& \func{execvp}& \func{execve} \\
    \hline
    \hline
    argomenti a lista    &$\bullet$&$\bullet$&$\bullet$&&& \\
    argomenti a vettore  &&&&$\bullet$&$\bullet$&$\bullet$\\
    \hline
    filename completo     &$\bullet$&&$\bullet$&$\bullet$&&$\bullet$\\ 
    ricerca su \var{PATH} &&$\bullet$&&&$\bullet$& \\
    \hline
    ambiente a vettore   &&&$\bullet$&&&$\bullet$ \\
    uso di \var{environ} &$\bullet$&$\bullet$&&$\bullet$&$\bullet$& \\
    \hline
  \end{tabular}
  \caption{Confronto delle caratteristiche delle varie funzioni della 
    famiglia \func{exec}.}
  \label{tab:proc_exec_scheme}
\end{table}

La prima differenza fra le funzioni riguarda le modalità di passaggio dei
valori che poi andranno a costituire gli argomenti a linea di comando (cioè i
valori di \param{argv} e \param{argc} visti dalla funzione \code{main} del
programma chiamato). Queste modalità sono due e sono riassunte dagli mnemonici
``\texttt{v}'' e ``\texttt{l}'' che stanno rispettivamente per \textit{vector}
e \textit{list}.

Nel primo caso gli argomenti sono passati tramite il vettore di puntatori
\var{argv[]} a stringhe terminate con zero che costituiranno gli argomenti a
riga di comando, questo vettore \emph{deve} essere terminato da un puntatore
nullo. Nel secondo caso le stringhe degli argomenti sono passate alla funzione
come lista di puntatori, nella forma:
\includecodesnip{listati/char_list.c}
che deve essere terminata da un puntatore nullo.  In entrambi i casi vale la
convenzione che il primo argomento (\var{arg0} o \var{argv[0]}) viene usato
per indicare il nome del file che contiene il programma che verrà eseguito.

La seconda differenza fra le funzioni riguarda le modalità con cui si
specifica il programma che si vuole eseguire. Con lo mnemonico ``\texttt{p}''
si indicano le due funzioni che replicano il comportamento della shell nello
specificare il comando da eseguire; quando l'argomento \param{file} non
contiene una ``\texttt{/}'' esso viene considerato come un nome di programma,
e viene eseguita automaticamente una ricerca fra i file presenti nella lista
di directory specificate dalla variabile di ambiente \envvar{PATH}. Il file
che viene posto in esecuzione è il primo che viene trovato. Se si ha un errore
relativo a permessi di accesso insufficienti (cioè l'esecuzione della
sottostante \func{execve} ritorna un \errcode{EACCES}), la ricerca viene
proseguita nelle eventuali ulteriori directory indicate in \envvar{PATH}; solo
se non viene trovato nessun altro file viene finalmente restituito
\errcode{EACCES}.  Le altre quattro funzioni si limitano invece a cercare di
eseguire il file indicato dall'argomento \param{path}, che viene interpretato
come il \textit{pathname} del programma.

La terza differenza è come viene passata la lista delle variabili di ambiente.
Con lo mnemonico ``\texttt{e}'' vengono indicate quelle funzioni che
necessitano di un vettore di parametri \var{envp[]} analogo a quello usato per
gli argomenti a riga di comando (terminato quindi da un \val{NULL}), le altre
usano il valore della variabile \var{environ} (vedi
sez.~\ref{sec:proc_environ}) del processo di partenza per costruire
l'ambiente.

Oltre a mantenere lo stesso \ids{PID}, il nuovo programma fatto partire da una
delle funzioni della famiglia \func{exec} mantiene la gran parte delle
proprietà del processo chiamante; una lista delle più significative è la
seguente:
\begin{itemize*}
\item il \textit{process id} (\ids{PID}) ed il \textit{parent process id}
  (\ids{PPID});
\item l'\textsl{user-ID reale}, il \textsl{group-ID reale} ed i
  \textsl{group-ID supplementari} (vedi sez.~\ref{sec:proc_access_id});
\item la directory radice (vedi sez.~\ref{sec:file_chroot}) e la directory di
  lavoro corrente (vedi sez.~\ref{sec:file_work_dir});
\item la maschera di creazione dei file (\textit{umask}, vedi
  sez.~\ref{sec:file_perm_management}) ed i \textit{lock} sui file (vedi
  sez.~\ref{sec:file_locking});
\item il valori di \textit{nice}, le priorità real-time e le affinità di
  processore (vedi sez.~\ref{sec:proc_sched_stand};
  sez.~\ref{sec:proc_real_time} e sez.~\ref{sec:proc_sched_multiprocess});
\item il \textit{session ID} (\acr{sid}) ed il \textit{process group ID}
  (\acr{pgid}), vedi sez.~\ref{sec:sess_proc_group};
\item il terminale di controllo (vedi sez.~\ref{sec:sess_ctrl_term});
\item il tempo restante ad un allarme (vedi sez.~\ref{sec:sig_alarm_abort});
\item i limiti sulle risorse (vedi sez.~\ref{sec:sys_resource_limit});
\item i valori delle variabili \var{tms\_utime}, \var{tms\_stime};
  \var{tms\_cutime}, \var{tms\_ustime} (vedi sez.~\ref{sec:sys_cpu_times});
\item la maschera dei segnali (si veda sez.~\ref{sec:sig_sigmask});
\item l'insieme dei segnali pendenti (vedi sez.~\ref{sec:sig_gen_beha}).
\end{itemize*}

Una serie di proprietà del processo originale, che non avrebbe senso mantenere
in un programma che esegue un codice completamente diverso in uno spazio di
indirizzi totalmente indipendente e ricreato da zero, vengono perse con
l'esecuzione di una \func{exec}. Lo standard POSIX.1-2001 prevede che le
seguenti proprietà non vengano preservate:
\begin{itemize*}
\item gli eventuali stack alternativi per i segnali (vedi
  sez.~\ref{sec:sig_specific_features});
\item i \textit{directory stream} (vedi sez.~\ref{sec:file_dir_read}), che
  vengono chiusi;
\item le mappature dei file in memoria (vedi sez.~\ref{sec:file_memory_map});
\item i segmenti di memoria condivisa SysV (vedi sez.~\ref{sec:ipc_sysv_shm})
  e POSIX (vedi sez.~\ref{sec:ipc_posix_shm});
\item i \textit{memory lock} (vedi sez.~\ref{sec:proc_mem_lock});
\item le funzioni registrate all'uscita (vedi sez.~\ref{sec:proc_atexit});
\item i semafori e le code di messaggi POSIX (vedi
  sez.~\ref{sec:ipc_posix_sem} e sez.~\ref{sec:ipc_posix_mq});
\item i timer POSIX (vedi sez.~\ref{sec:sig_timer_adv}).
\end{itemize*}

Inoltre i segnali che sono stati impostati per essere ignorati nel processo
chiamante mantengono la stessa impostazione pure nel nuovo programma, ma tutti
gli altri segnali, ed in particolare quelli per i quali è stato installato un
gestore vengono impostati alla loro azione predefinita (vedi
sez.~\ref{sec:sig_gen_beha}). Un caso speciale è il segnale \signal{SIGCHLD}
che, quando impostato a \const{SIG\_IGN}, potrebbe anche essere reimpostato a
\const{SIG\_DFL}. Lo standard POSIX.1-2001 prevede che questo comportamento
sia deciso dalla singola implementazione, quella di Linux è di non modificare
l'impostazione precedente.

Oltre alle precedenti, che sono completamente generali e disponibili anche su
altri sistemi unix-like, esistono altre proprietà dei processi, attinenti alle
caratteristiche specifiche di Linux, che non vengono preservate
nell'esecuzione della funzione \func{exec}, queste sono:
\begin{itemize*}
\item le operazioni di I/O asincrono (vedi sez.~\ref{sec:file_asyncronous_io})
  pendenti vengono cancellate;
\item le \textit{capabilities} vengono modificate come
  illustrato in sez.~\ref{sec:proc_capabilities};
\item tutti i \textit{thread} tranne il chiamante\unavref{ (vedi
    sez.~\ref{sec:thread_xxx})} vengono cancellati e tutti gli oggetti ad essi
  relativi\unavref{ (vedi sez.~\ref{sec:thread_xxx})} sono rimossi;
\item viene impostato il flag \const{PR\_SET\_DUMPABLE} di \func{prctl} (vedi
  sez.~\ref{sec:process_prctl}) a meno che il programma da eseguire non sia
  \acr{suid} o \acr{sgid} (vedi sez.~\ref{sec:proc_access_id} e
  sez.~\ref{sec:file_special_perm});
\item il flag \const{PR\_SET\_KEEPCAPS} di \func{prctl} (vedi
  sez.~\ref{sec:process_prctl}) viene cancellato;
\item il nome del processo viene impostato al nome del file contenente il
  programma messo in esecuzione;
\item il segnale di terminazione viene reimpostato a \signal{SIGCHLD};
\item l'ambiente viene reinizializzato impostando le variabili attinenti alla
  localizzazione al valore di default POSIX. 
\end{itemize*}

\itindbeg{close-on-exec}

La gestione dei file aperti nel passaggio al nuovo programma lanciato con
\func{exec} dipende dal valore che ha il flag di \textit{close-on-exec} (vedi
sez.~\ref{sec:file_fcntl_ioctl}) per ciascun \textit{file descriptor}. I file
per cui è impostato vengono chiusi, tutti gli altri file restano
aperti. Questo significa che il comportamento predefinito è che i file restano
aperti attraverso una \func{exec}, a meno di una chiamata esplicita a
\func{fcntl} che imposti il suddetto flag.  Per le directory, lo standard
POSIX.1 richiede che esse vengano chiuse attraverso una \func{exec}, in genere
questo è fatto dalla funzione \func{opendir} (vedi
sez.~\ref{sec:file_dir_read}) che effettua da sola l'impostazione del flag di
\textit{close-on-exec} sulle directory che apre, in maniera trasparente
all'utente.

\itindend{close-on-exec}

Il comportamento della funzione in relazione agli identificatori relativi al
controllo di accesso verrà trattato in dettaglio in sez.~\ref{sec:proc_perms},
qui è sufficiente anticipare (si faccia riferimento a
sez.~\ref{sec:proc_access_id} per la definizione di questi identificatori)
come l'\textsl{user-ID reale} ed il \textsl{group-ID reale} restano sempre gli
stessi, mentre l'\textsl{user-ID salvato} ed il \textsl{group-ID salvato}
vengono impostati rispettivamente all'\textsl{user-ID effettivo} ed il
\textsl{group-ID effettivo}. Questi ultimi normalmente non vengono modificati,
a meno che il file di cui viene chiesta l'esecuzione non abbia o il \acr{suid}
bit o lo \acr{sgid} bit impostato (vedi sez.~\ref{sec:file_special_perm}), in
questo caso l'\textsl{user-ID effettivo} ed il \textsl{group-ID effettivo}
vengono impostati rispettivamente all'utente o al gruppo cui il file
appartiene.

Se il file da eseguire è in formato \emph{a.out} e necessita di librerie
condivise, viene lanciato il \textit{linker} dinamico \cmd{/lib/ld.so} prima
del programma per caricare le librerie necessarie ed effettuare il link
dell'eseguibile; il formato è ormai in completo disuso, per cui è molto
probabile che non il relativo supporto non sia disponibile. Se il programma è
in formato ELF per caricare le librerie dinamiche viene usato l'interprete
indicato nel segmento \constd{PT\_INTERP} previsto dal formato stesso, in
genere questo è \sysfiled{/lib/ld-linux.so.1} per programmi collegati con la
\acr{libc5}, e \sysfiled{/lib/ld-linux.so.2} per programmi collegati con la
\acr{glibc}.

Infine nel caso il programma che si vuole eseguire sia uno script e non un
binario, questo deve essere un file di testo che deve iniziare con una linea
nella forma:
\begin{Example}
#!/path/to/interpreter [argomenti]
\end{Example}
dove l'interprete indicato deve essere un eseguibile binario e non un altro
script, che verrà chiamato come se si fosse eseguito il comando
\cmd{interpreter [argomenti] filename}. 

Si tenga presente che con Linux quanto viene scritto come \texttt{argomenti}
viene passato all'interprete come un unico argomento con una unica stringa di
lunghezza massima di 127 caratteri e se questa dimensione viene ecceduta la
stringa viene troncata; altri Unix hanno dimensioni massime diverse, e diversi
comportamenti, ad esempio FreeBSD esegue la scansione della riga e la divide
nei vari argomenti e se è troppo lunga restituisce un errore di
\errval{ENAMETOOLONG}; una comparazione dei vari comportamenti sui diversi
sistemi unix-like si trova su
\url{http://www.in-ulm.de/~mascheck/various/shebang/}.

Con la famiglia delle \func{exec} si chiude il novero delle funzioni su cui è
basata la gestione tradizionale dei processi in Unix: con \func{fork} si crea
un nuovo processo, con \func{exec} si lancia un nuovo programma, con
\func{exit} e \func{wait} si effettua e verifica la conclusione dei
processi. Tutte le altre funzioni sono ausiliarie e servono per la lettura e
l'impostazione dei vari parametri connessi ai processi.



\section{Il controllo di accesso}
\label{sec:proc_perms}

In questa sezione esamineremo le problematiche relative al controllo di
accesso dal punto di vista dei processi; vedremo quali sono gli identificatori
usati, come questi possono essere modificati nella creazione e nel lancio di
nuovi processi, le varie funzioni per la loro manipolazione diretta e tutte le
problematiche connesse ad una gestione accorta dei privilegi.


\subsection{Gli identificatori del controllo di accesso}
\label{sec:proc_access_id}

Come accennato in sez.~\ref{sec:intro_multiuser} il modello base\footnote{in
  realtà già esistono estensioni di questo modello base, che lo rendono più
  flessibile e controllabile, come le \textit{capabilities} illustrate in
  sez.~\ref{sec:proc_capabilities}, le ACL per i file (vedi
  sez.~\ref{sec:file_ACL}) o il \textit{Mandatory Access Control} di
  \textit{SELinux}; inoltre basandosi sul lavoro effettuato con
  \textit{SELinux}, a partire dal kernel 2.5.x, è iniziato lo sviluppo di una
  infrastruttura di sicurezza, i \textit{Linux Security Modules}, o LSM, in
  grado di fornire diversi agganci a livello del kernel per modularizzare
  tutti i possibili controlli di accesso, cosa che ha permesso di realizzare
  diverse alternative a \textit{SELinux}.} 
di sicurezza di un sistema unix-like è fondato sui concetti di utente e
gruppo, e sulla separazione fra l'amministratore (\textsl{root}, detto spesso
anche \textit{superuser}) che non è sottoposto a restrizioni, ed il resto
degli utenti, per i quali invece vengono effettuati i vari controlli di
accesso.

Abbiamo già accennato come il sistema associ ad ogni utente e gruppo due
identificatori univoci, lo \itindex{User~ID~(UID)} \textsl{User-ID}
(abbreviato in \ids{UID}) ed il \itindex{Group~ID~(GID)} \textsl{Group-ID}
(abbreviato in \ids{GID}). Questi servono al kernel per identificare uno
specifico utente o un gruppo di utenti, per poi poter controllare che essi
siano autorizzati a compiere le operazioni richieste.  Ad esempio in
sez.~\ref{sec:file_access_control} vedremo come ad ogni file vengano associati
un utente ed un gruppo (i suoi \textsl{proprietari}, indicati appunto tramite
un \ids{UID} ed un \ids{GID}) che vengono controllati dal kernel nella
gestione dei permessi di accesso.

Dato che tutte le operazioni del sistema vengono compiute dai processi, è
evidente che per poter implementare un controllo sulle operazioni occorre
anche poter identificare chi è che ha lanciato un certo programma, e pertanto
anche a ciascun processo dovrà essere associato un utente e un gruppo.

Un semplice controllo di una corrispondenza fra identificativi non garantisce
però sufficiente flessibilità per tutti quei casi in cui è necessario poter
disporre di privilegi diversi, o dover impersonare un altro utente per un
limitato insieme di operazioni. Per questo motivo in generale tutti i sistemi
unix-like prevedono che i processi abbiano almeno due gruppi di
identificatori, chiamati rispettivamente \textit{real} ed \textit{effective}
(cioè \textsl{reali} ed \textsl{effettivi}). Nel caso di Linux si aggiungono
poi altri due gruppi, il \textit{saved} (\textsl{salvati}) ed il
\textit{filesystem} (\textsl{di filesystem}), secondo la situazione illustrata
in tab.~\ref{tab:proc_uid_gid}.

\begin{table}[htb]
  \footnotesize
  \centering
  \begin{tabular}[c]{|c|c|l|p{7cm}|}
    \hline
    \textbf{Suffisso} & \textbf{Gruppo} & \textbf{Denominazione} 
                                        & \textbf{Significato} \\ 
    \hline
    \hline
    \texttt{uid} & \textit{real} & \textsl{user-ID reale} 
                 & Indica l'utente che ha lanciato il programma.\\ 
    \texttt{gid} & '' &\textsl{group-ID reale} 
                 & Indica il gruppo principale dell'utente che ha lanciato 
                   il programma.\\ 
    \hline
    \texttt{euid}& \textit{effective} &\textsl{user-ID effettivo} 
                 & Indica l'utente usato nel controllo di accesso.\\ 
    \texttt{egid}& '' & \textsl{group-ID effettivo} 
                 & Indica il gruppo usato nel controllo di accesso.\\ 
    --           & -- & \textsl{group-ID supplementari} 
                 & Indicano gli ulteriori gruppi cui l'utente appartiene.\\ 
    \hline
    --           & \textit{saved} & \textsl{user-ID salvato} 
                 & Mantiene una copia dell'\acr{euid} iniziale.\\ 
    --           & '' & \textsl{group-ID salvato} 
                 & Mantiene una copia dell'\acr{egid} iniziale.\\ 
    \hline
    \texttt{fsuid}& \textit{filesystem} &\textsl{user-ID di filesystem} 
                 & Indica l'utente effettivo per l'accesso al filesystem. \\ 
    \texttt{fsgid}& '' & \textsl{group-ID di filesystem} 
                 & Indica il gruppo effettivo per l'accesso al filesystem.\\ 
    \hline
  \end{tabular}
  \caption{Identificatori di utente e gruppo associati a ciascun processo con
    indicazione dei suffissi usati dalle varie funzioni di manipolazione.}
  \label{tab:proc_uid_gid}
\end{table}

Al primo gruppo appartengono l'\ids{UID} \textsl{reale} ed il \ids{GID}
\textsl{reale}: questi vengono impostati al login ai valori corrispondenti
all'utente con cui si accede al sistema (e relativo gruppo principale).
Servono per l'identificazione dell'utente e normalmente non vengono mai
cambiati. In realtà vedremo (in sez.~\ref{sec:proc_setuid}) che è possibile
modificarli, ma solo ad un processo che abbia i privilegi di amministratore;
questa possibilità è usata proprio dal programma \cmd{login} che, una volta
completata la procedura di autenticazione, lancia una shell per la quale
imposta questi identificatori ai valori corrispondenti all'utente che entra
nel sistema.

Al secondo gruppo appartengono l'\ids{UID} \textsl{effettivo} e il \ids{GID}
\textsl{effettivo}, a cui si aggiungono gli eventuali \ids{GID}
\textsl{supplementari} dei gruppi dei quali l'utente fa parte.  Questi sono
invece gli identificatori usati nelle verifiche dei permessi del processo e
per il controllo di accesso ai file (argomento affrontato in dettaglio in
sez.~\ref{sec:file_perm_overview}).

Questi identificatori normalmente sono identici ai corrispondenti del gruppo
\textit{real} tranne nel caso in cui, come accennato in
sez.~\ref{sec:proc_exec}, il programma che si è posto in esecuzione abbia i
bit \acr{suid} o \acr{sgid} impostati (il significato di questi bit è
affrontato in dettaglio in sez.~\ref{sec:file_special_perm}). In questo caso
essi saranno impostati all'utente e al gruppo proprietari del file. Questo
consente, per programmi in cui ci sia questa necessità, di dare a qualunque
utente i privilegi o i permessi di un altro, compreso l'amministratore.

Come nel caso del \ids{PID} e del \ids{PPID}, anche tutti questi
identificatori possono essere ottenuti da un programma attraverso altrettante
funzioni di sistema dedicate alla loro lettura, queste sono \funcd{getuid},
\funcd{geteuid}, \funcd{getgid} e \funcd{getegid}, ed i loro prototipi sono:

\begin{funcproto}{ 
\fhead{unistd.h}
\fhead{sys/types.h}
\fdecl{uid\_t getuid(void)}
\fdesc{Legge l'\ids{UID} reale del processo corrente.} 
\fdecl{uid\_t geteuid(void)}
\fdesc{Legge l'\ids{UID} effettivo del processo corrente.} 
\fdecl{gid\_t getgid(void)}
\fdesc{Legge il \ids{GID} reale del processo corrente.} 
\fdecl{gid\_t getegid(void)}
\fdesc{Legge il \ids{GID} effettivo del processo corrente.}
}
{Le funzioni ritornano i rispettivi identificativi del processo corrente, e
  non sono previste condizioni di errore.}
\end{funcproto}

In generale l'uso di privilegi superiori, ottenibile con un \ids{UID}
\textsl{effettivo} diverso da quello reale, deve essere limitato il più
possibile, per evitare abusi e problemi di sicurezza, per questo occorre anche
un meccanismo che consenta ad un programma di rilasciare gli eventuali
maggiori privilegi necessari, una volta che si siano effettuate le operazioni
per i quali erano richiesti, e a poterli eventualmente recuperare in caso
servano di nuovo.

Questo in Linux viene fatto usando altri due gruppi di identificatori, il
\textit{saved} ed il \textit{filesystem}. Il primo gruppo è lo stesso usato in
SVr4, e previsto dallo standard POSIX quando è definita
\macro{\_POSIX\_SAVED\_IDS},\footnote{in caso si abbia a cuore la portabilità
  del programma su altri Unix è buona norma controllare sempre la
  disponibilità di queste funzioni controllando se questa costante è
  definita.} il secondo gruppo è specifico di Linux e viene usato per
migliorare la sicurezza con NFS (il \textit{Network File System}, protocollo
che consente di accedere ai file via rete).

L'\ids{UID} \textsl{salvato} ed il \ids{GID} \textsl{salvato} sono copie
dell'\ids{UID} \textsl{effettivo} e del \ids{GID} \textsl{effettivo} del
processo padre, e vengono impostati dalla funzione \func{exec} all'avvio del
processo, come copie dell'\ids{UID} \textsl{effettivo} e del \ids{GID}
\textsl{effettivo} dopo che questi sono stati impostati tenendo conto di
eventuali permessi \acr{suid} o \acr{sgid} (su cui torneremo in
sez.~\ref{sec:file_special_perm}).  Essi quindi consentono di tenere traccia
di quale fossero utente e gruppo effettivi all'inizio dell'esecuzione di un
nuovo programma.

L'\ids{UID} \textsl{di filesystem} e il \ids{GID} \textsl{di filesystem} sono
un'estensione introdotta in Linux per rendere più sicuro l'uso di NFS
(torneremo sull'argomento in sez.~\ref{sec:proc_setuid}). Essi sono una
replica dei corrispondenti identificatori del gruppo \textit{effective}, ai
quali si sostituiscono per tutte le operazioni di verifica dei permessi
relativi ai file (trattate in sez.~\ref{sec:file_perm_overview}).  Ogni
cambiamento effettuato sugli identificatori effettivi viene automaticamente
riportato su di essi, per cui in condizioni normali si può tranquillamente
ignorarne l'esistenza, in quanto saranno del tutto equivalenti ai precedenti.


\subsection{Le funzioni di gestione degli identificatori dei processi}
\label{sec:proc_setuid}

Le funzioni di sistema più comuni che vengono usate per cambiare identità
(cioè utente e gruppo di appartenenza) ad un processo, e che come accennato in
sez.~\ref{sec:proc_access_id} seguono la semantica POSIX che prevede
l'esistenza dell'\ids{UID} salvato e del \ids{GID} salvato, sono
rispettivamente \funcd{setuid} e \funcd{setgid}; i loro prototipi sono:

\begin{funcproto}{
\fhead{unistd.h}
\fhead{sys/types.h}
\fdecl{int setuid(uid\_t uid)}
\fdesc{Imposta l'\ids{UID} del processo corrente.}
\fdecl{int setgid(gid\_t gid)}
\fdesc{Imposta il \ids{GID} del processo corrente.}
}
{Le funzioni ritornano $0$ in caso di successo e $-1$ per un errore, nel qual
caso \var{errno} uno dei valori: 
\begin{errlist}
\item[\errcode{EAGAIN}] (solo per \func{setuid}) la chiamata cambierebbe
  l'\ids{UID} reale ma il kernel non dispone temporaneamente delle risorse per
  farlo, oppure, per i kernel precedenti il 3.1, il cambiamento
  dell'\ids{UID} reale farebbe superare il limite per il numero dei processi
  \const{RLIMIT\_NPROC} (vedi sez.~\ref{sec:sys_resource_limit}).
\item[\errcode{EINVAL}] il valore di dell'argomento non è valido per il
    \textit{namespace} corrente (vedi sez.~\ref{sec:process_namespaces}).
\item[\errcode{EPERM}] non si hanno i permessi per l'operazione richiesta.
\end{errlist} 
}
\end{funcproto}

Il funzionamento di queste due funzioni è analogo, per cui considereremo solo
la prima, la seconda si comporta esattamente allo stesso modo facendo
riferimento al \ids{GID} invece che all'\ids{UID}. Gli eventuali \ids{GID}
supplementari non vengono modificati.

L'effetto della chiamata è diverso a seconda dei privilegi del processo; se
l'\ids{UID} effettivo è zero (cioè è quello dell'amministratore di sistema o
il processo ha la capacità \const{CAP\_SETUID}) allora tutti gli
identificatori (\textit{real}, \textit{effective} e \textit{saved}) vengono
impostati al valore specificato da \param{uid}, altrimenti viene impostato
solo l'\ids{UID} effettivo, e soltanto se il valore specificato corrisponde o
all'\ids{UID} reale o all'\ids{UID} salvato, ottenendo un errore di
\errcode{EPERM} negli altri casi.

E' importante notare che la funzione può fallire, con
\errval{EAGAIN},\footnote{non affronteremo qui l'altro caso di errore, che può
  avvenire solo quando si esegue la funzione all'interno di un diverso
  \textit{user namespace}, argomento su cui torneremo in
  sez.~\ref{sec:process_namespaces} ma la considerazione di controllare sempre
  lo stato di uscita si applica allo stesso modo.} anche quando viene invocata
da un processo con privilegi di amministratore per cambiare il proprio
l'\ids{UID} reale, sia per una temporanea indisponibilità di risorse del
kernel, sia perché l'utente di cui si vuole assumere l'\ids{UID} andrebbe a
superare un eventuale limite sul numero di processi (il limite
\const{RLIMIT\_NPROC}, che tratteremo in sez.~\ref{sec:sys_resource_limit}),
pertanto occorre sempre verificare lo stato di uscita della funzione.

Non controllare questo tipo di errori perché si presume che la funzione abbia
sempre successo quando si hanno i privilegi di amministratore può avere
conseguente devastanti per la sicurezza, in particolare quando la si usa per
cedere i suddetti privilegi ed eseguire un programma per conto di un utente
non privilegiato.

E' per diminuire l'impatto di questo tipo di disattenzioni che a partire dal
kernel 3.1 il comportamento di \func{setuid} e di tutte le analoghe funzioni
che tratteremo nel seguito di questa sezione è stato modificato e nel caso di
superamento del limite sulle risorse esse hanno comunque successo. Quando
questo avviene il processo assume comunque il nuovo \ids{UID} ed il controllo
sul superamento di \const{RLIMIT\_NPROC} viene posticipato ad una eventuale
successiva invocazione di \func{execve} (essendo questo poi il caso d'uso più
comune). In tal caso, se alla chiamata ancora sussiste la situazione di
superamento del limite, sarà \func{execve} a fallire con un errore di
\const{EAGAIN}.\footnote{che pertanto, a partire dal kernel 3.1, può
  restituire anche questo errore, non presente in altri sistemi
  \textit{unix-like}.}

Come accennato l'uso principale di queste funzioni è quello di poter
consentire ad un programma con i bit \acr{suid} o \acr{sgid} impostati (vedi
sez.~\ref{sec:file_special_perm}) di riportare l'\ids{UID} effettivo a quello
dell'utente che ha lanciato il programma, effettuare il lavoro che non
necessita di privilegi aggiuntivi, ed eventualmente tornare indietro.

Come esempio per chiarire l'uso di queste funzioni prendiamo quello con cui
viene gestito l'accesso al file \sysfiled{/var/run/utmp}.  In questo file viene
registrato chi sta usando il sistema al momento corrente; chiaramente non può
essere lasciato aperto in scrittura a qualunque utente, che potrebbe
falsificare la registrazione.

Per questo motivo questo file (e l'analogo \sysfiled{/var/log/wtmp} su cui
vengono registrati login e logout) appartengono ad un gruppo dedicato (in
genere \acr{utmp}) ed i programmi che devono accedervi (ad esempio tutti i
programmi di terminale in X, o il programma \cmd{screen} che crea terminali
multipli su una console) appartengono a questo gruppo ed hanno il bit
\acr{sgid} impostato.

Quando uno di questi programmi (ad esempio \cmd{xterm}) viene lanciato, la
situazione degli identificatori è la seguente:
\begin{eqnarray*}
  \label{eq:1}
  \textsl{group-ID reale}      &=& \textrm{\ids{GID} (del chiamante)} \\
  \textsl{group-ID effettivo}  &=& \textrm{\acr{utmp}} \\
  \textsl{group-ID salvato}    &=& \textrm{\acr{utmp}}
\end{eqnarray*}
in questo modo, dato che il \textsl{group-ID effettivo} è quello giusto, il
programma può accedere a \sysfile{/var/run/utmp} in scrittura ed aggiornarlo.
A questo punto il programma può eseguire una \code{setgid(getgid())} per
impostare il \textsl{group-ID effettivo} a quello dell'utente (e dato che il
\textsl{group-ID reale} corrisponde la funzione avrà successo), in questo modo
non sarà possibile lanciare dal terminale programmi che modificano detto file,
in tal caso infatti la situazione degli identificatori sarebbe:
\begin{eqnarray*}
  \label{eq:2}
  \textsl{group-ID reale}      &=& \textrm{\ids{GID} (invariato)}  \\
  \textsl{group-ID effettivo}  &=& \textrm{\ids{GID}} \\
  \textsl{group-ID salvato}    &=& \textrm{\acr{utmp} (invariato)}
\end{eqnarray*}
e ogni processo lanciato dal terminale avrebbe comunque \ids{GID} come
\textsl{group-ID effettivo}. All'uscita dal terminale, per poter di nuovo
aggiornare lo stato di \sysfile{/var/run/utmp} il programma eseguirà una
\code{setgid(utmp)} (dove \var{utmp} è il valore numerico associato al gruppo
\acr{utmp}, ottenuto ad esempio con una precedente \func{getegid}), dato che
in questo caso il valore richiesto corrisponde al \textsl{group-ID salvato} la
funzione avrà successo e riporterà la situazione a:
\begin{eqnarray*}
  \label{eq:3}
  \textsl{group-ID reale}      &=& \textrm{\ids{GID} (invariato)}  \\
  \textsl{group-ID effettivo}  &=& \textrm{\acr{utmp}} \\
  \textsl{group-ID salvato}    &=& \textrm{\acr{utmp} (invariato)}
\end{eqnarray*}
consentendo l'accesso a \sysfile{/var/run/utmp}.

Occorre però tenere conto che tutto questo non è possibile con un processo con
i privilegi di amministratore, in tal caso infatti l'esecuzione di una
\func{setuid} comporta il cambiamento di tutti gli identificatori associati al
processo, rendendo impossibile riguadagnare i privilegi di amministratore.
Questo comportamento è corretto per l'uso che ne fa un programma come
\cmd{login} una volta che crea una nuova shell per l'utente, ma quando si
vuole cambiare soltanto l'\ids{UID} effettivo del processo per cedere i
privilegi occorre ricorrere ad altre funzioni.

Le due funzioni di sistema \funcd{setreuid} e \funcd{setregid} derivano da BSD
che, non supportando (almeno fino alla versione 4.3+BSD) gli identificatori
del gruppo \textit{saved}, le usa per poter scambiare fra di loro
\textit{effective} e \textit{real}; i rispettivi prototipi sono:

\begin{funcproto}{ 
\fhead{unistd.h}
\fhead{sys/types.h}
\fdecl{int setreuid(uid\_t ruid, uid\_t euid)}
\fdesc{Imposta \ids{UID} reale e \ids{UID} effettivo del processo corrente.} 
\fdecl{int setregid(gid\_t rgid, gid\_t egid)}
\fdesc{Imposta \ids{GID} reale e \ids{GID} effettivo del processo corrente.} 
}
{Le funzioni ritornano $0$ in caso di successo e $-1$ per un errore, nel qual
caso \var{errno} assume i valori visti per \func{setuid}/\func{setgid}.
}
\end{funcproto}

Le due funzioni sono identiche, quanto diremo per la prima riguardo gli
\ids{UID} si applica alla seconda per i \ids{GID}.  La funzione
\func{setreuid} imposta rispettivamente l'\ids{UID} reale e l'\ids{UID}
effettivo del processo corrente ai valori specificati da \param{ruid}
e \param{euid}.

I processi non privilegiati possono impostare solo valori che corrispondano o
al loro \ids{UID} effettivo o a quello reale o a quello salvato, valori
diversi comportano il fallimento della chiamata.  L'amministratore invece può
specificare un valore qualunque.  Specificando un argomento di valore $-1$
l'identificatore corrispondente verrà lasciato inalterato.

Con queste funzioni si possono scambiare fra loro gli \ids{UID} reale ed
effettivo, e pertanto è possibile implementare un comportamento simile a
quello visto in precedenza per \func{setgid}, cedendo i privilegi con un primo
scambio, e recuperandoli, una volta eseguito il lavoro non privilegiato, con
un secondo scambio.

In questo caso però occorre porre molta attenzione quando si creano nuovi
processi nella fase intermedia in cui si sono scambiati gli identificatori, in
questo caso infatti essi avranno un \ids{UID} reale privilegiato, che dovrà
essere esplicitamente eliminato prima di porre in esecuzione un nuovo
programma, occorrerà cioè eseguire un'altra chiamata dopo la \func{fork} e
prima della \func{exec} per uniformare l'\ids{UID} reale a quello effettivo,
perché in caso contrario il nuovo programma potrebbe a sua volta effettuare
uno scambio e riottenere dei privilegi non previsti.

Lo stesso problema di propagazione dei privilegi ad eventuali processi figli
si pone anche per l'\ids{UID} salvato. Ma la funzione \func{setreuid} deriva
da un'implementazione di sistema che non ne prevede la presenza, e quindi non
è possibile usarla per correggere la situazione come nel caso precedente. Per
questo motivo in Linux tutte le volte che si imposta un qualunque valore
diverso da quello dall'\ids{UID} reale corrente, l'\ids{UID} salvato viene
automaticamente uniformato al valore dell'\ids{UID} effettivo.

Altre due funzioni di sistema, \funcd{seteuid} e \funcd{setegid}, sono
un'estensione dello standard POSIX.1, ma sono comunque supportate dalla
maggior parte degli Unix, esse vengono usate per cambiare gli identificatori
del gruppo \textit{effective} ed i loro prototipi sono:

\begin{funcproto}{ 
\fhead{unistd.h}
\fhead{sys/types.h}
\fdecl{int seteuid(uid\_t uid)}
\fdesc{Imposta l'\ids{UID} effettivo del processo corrente.} 
\fdecl{int setegid(gid\_t gid)}
\fdesc{Imposta il \ids{GID} effettivo del processo corrente.} 
}
{Le funzioni ritornano $0$ in caso di successo e $-1$ per un errore, nel qual
  caso \var{errno} assume i valori visti per \func{setuid}/\func{setgid}
  tranne \errval{EAGAIN}. 
}
\end{funcproto}

Ancora una volta le due funzioni sono identiche, e quanto diremo per la prima
riguardo gli \ids{UID} si applica allo stesso modo alla seconda per i
\ids{GID}. Con \func{seteuid} gli utenti normali possono impostare l'\ids{UID}
effettivo solo al valore dell'\ids{UID} reale o dell'\ids{UID} salvato,
l'amministratore può specificare qualunque valore. Queste funzioni sono usate
per permettere all'amministratore di impostare solo l'\ids{UID} effettivo,
dato che l'uso normale di \func{setuid} comporta l'impostazione di tutti gli
identificatori.
 
Le due funzioni di sistema \funcd{setresuid} e \funcd{setresgid} sono invece
un'estensione introdotta in Linux (a partire dal kernel 2.1.44) e permettono
un completo controllo su tutti e tre i gruppi di identificatori
(\textit{real}, \textit{effective} e \textit{saved}), i loro prototipi sono:

\begin{funcproto}{ 
\fhead{unistd.h}
\fhead{sys/types.h}
\fdecl{int setresuid(uid\_t ruid, uid\_t euid, uid\_t suid)}
\fdesc{Imposta l'\ids{UID} reale, effettivo e salvato del processo corrente.} 
\fdecl{int setresgid(gid\_t rgid, gid\_t egid, gid\_t sgid)}
\fdesc{Imposta il \ids{GID} reale, effettivo e salvato del processo corrente.} 
}
{Le funzioni ritornano $0$ in caso di successo e $-1$ per un errore, nel qual
caso \var{errno} assume i valori visti per \func{setuid}/\func{setgid}.
}
\end{funcproto}

Di nuovo le due funzioni sono identiche e quanto detto per la prima riguardo
gli \ids{UID} si applica alla seconda per i \ids{GID}.  La funzione
\func{setresuid} imposta l'\ids{UID} reale, l'\ids{UID} effettivo e
l'\ids{UID} salvato del processo corrente ai valori specificati
rispettivamente dagli argomenti \param{ruid}, \param{euid} e \param{suid}.  I
processi non privilegiati possono cambiare uno qualunque degli \ids{UID} solo
ad un valore corrispondente o all'\ids{UID} reale, o a quello effettivo o a
quello salvato, l'amministratore può specificare i valori che vuole. Un valore
di $-1$ per un qualunque argomento lascia inalterato l'identificatore
corrispondente.

Per queste funzioni di sistema esistono anche due controparti,
\funcd{getresuid} e \funcd{getresgid},\footnote{le funzioni non sono standard,
  anche se appaiono in altri kernel, su Linux sono presenti dal kernel 2.1.44
  e con le versioni della \acr{glibc} a partire dalla 2.3.2, definendo la
  macro \macro{\_GNU\_SOURCE}.} che permettono di leggere in blocco i vari
identificatori; i loro prototipi sono:

\begin{funcproto}{ 
\fhead{unistd.h}
\fhead{sys/types.h}
\fdecl{int getresuid(uid\_t *ruid, uid\_t *euid, uid\_t *suid)}
\fdesc{Legge l'\ids{UID} reale, effettivo e salvato del processo corrente.} 
\fdecl{int getresgid(gid\_t *rgid, gid\_t *egid, gid\_t *sgid)}
\fdesc{Legge il \ids{GID} reale, effettivo e salvato del processo corrente.} 
}
{Le funzioni ritornano $0$ in caso di successo e $-1$ per un errore, nel qual
  caso \var{errno} può assumere solo il valore \errcode{EFAULT} se gli
  indirizzi delle variabili di ritorno non sono validi.  }
\end{funcproto}

Anche queste funzioni sono un'estensione specifica di Linux, e non richiedono
nessun privilegio. I valori sono restituiti negli argomenti, che vanno
specificati come puntatori (è un altro esempio di \textit{value result
  argument}). Si noti che queste funzioni sono le uniche in grado di leggere
gli identificatori del gruppo \textit{saved}.

Infine le funzioni \func{setfsuid} e \func{setfsgid} servono per impostare gli
identificatori del gruppo \textit{filesystem} che sono usati da Linux per il
controllo dell'accesso ai file.  Come già accennato in
sez.~\ref{sec:proc_access_id} Linux definisce questo ulteriore gruppo di
identificatori, che in circostanze normali sono assolutamente equivalenti a
quelli del gruppo \textit{effective}, dato che ogni cambiamento di questi
ultimi viene immediatamente riportato su di essi.

C'è un solo caso in cui si ha necessità di introdurre una differenza fra gli
identificatori dei gruppi \textit{effective} e \textit{filesystem}, ed è per
ovviare ad un problema di sicurezza che si presenta quando si deve
implementare un server NFS. 

Il server NFS infatti deve poter cambiare l'identificatore con cui accede ai
file per assumere l'identità del singolo utente remoto, ma se questo viene
fatto cambiando l'\ids{UID} effettivo o l'\ids{UID} reale il server si espone
alla ricezione di eventuali segnali ostili da parte dell'utente di cui ha
temporaneamente assunto l'identità.  Cambiando solo l'\ids{UID} di filesystem
si ottengono i privilegi necessari per accedere ai file, mantenendo quelli
originari per quanto riguarda tutti gli altri controlli di accesso, così che
l'utente non possa inviare segnali al server NFS.

Le due funzioni di sistema usate appositamente per cambiare questi
identificatori sono \funcd{setfsuid} e \funcd{setfsgid} ovviamente sono
specifiche di Linux e non devono essere usate se si intendono scrivere
programmi portabili; i loro prototipi sono:

\begin{funcproto}{ 
\fhead{sys/fsuid.h}
\fdecl{int setfsuid(uid\_t fsuid)}
\fdesc{Imposta l'\ids{UID} di filesystem del processo corrente.} 
\fdecl{int setfsgid(gid\_t fsgid)}
\fdesc{Legge il \ids{GID} di filesystem del processo corrente.} 
}

{Le funzioni restituiscono sia in caso di successo che di errore il valore
  corrente dell'identificativo, e in caso di errore non viene impostato nessun
  codice in \var{errno}.}
\end{funcproto}

Le due funzioni sono analoghe ed usano il valore passato come argomento per
effettuare l'impostazione dell'identificativo.  Le funzioni hanno successo
solo se il processo chiamante ha i privilegi di amministratore o, per gli
altri utenti, se il valore specificato coincide con uno dei di quelli del
gruppo \textit{real}, \textit{effective} o \textit{saved}.

Il problema di queste funzioni è che non restituiscono un codice di errore e
non c'è modo di sapere (con una singola chiamata) di sapere se hanno avuto
successo o meno, per verificarlo occorre eseguire una chiamata aggiuntiva
passando come argomento $-1$ (un valore impossibile per un identificativo),
così fallendo si può di ottenere il valore corrente e verificare se è
cambiato.

\subsection{Le funzioni per la gestione dei gruppi associati a un processo}
\label{sec:proc_setgroups}

Le ultime funzioni che esamineremo sono quelle che permettono di operare sui
gruppi supplementari cui un utente può appartenere. Ogni processo può avere
almeno \const{NGROUPS\_MAX} gruppi supplementari\footnote{il numero massimo di
  gruppi secondari può essere ottenuto con \func{sysconf} (vedi
  sez.~\ref{sec:sys_limits}), leggendo il parametro
  \const{\_SC\_NGROUPS\_MAX}.} in aggiunta al gruppo primario; questi vengono
ereditati dal processo padre e possono essere cambiati con queste funzioni.

La funzione di sistema che permette di leggere i gruppi supplementari
associati ad un processo è \funcd{getgroups}; questa funzione è definita nello
standard POSIX.1, ed il suo prototipo è:

\begin{funcproto}{ 
\fhead{sys/types.h}
\fhead{unistd.h}
\fdecl{int getgroups(int size, gid\_t list[])}
\fdesc{Legge gli identificatori dei gruppi supplementari.} 
}
{La funzione ritorna il numero di gruppi letti in caso di successo e $-1$ per
  un errore, nel qual caso \var{errno} assumerà uno dei valori:
\begin{errlist}
\item[\errcode{EFAULT}] \param{list} non ha un indirizzo valido.
\item[\errcode{EINVAL}] il valore di \param{size} è diverso da zero ma
  minore del numero di gruppi supplementari del processo.
\end{errlist}}
\end{funcproto}

La funzione legge gli identificatori dei gruppi supplementari del processo sul
vettore \param{list} che deve essere di dimensione pari a \param{size}. Non è
specificato se la funzione inserisca o meno nella lista il \ids{GID} effettivo
del processo. Se si specifica un valore di \param{size} uguale a $0$ allora
l'argomento \param{list} non viene modificato, ma si ottiene dal valore di
ritorno il numero di gruppi supplementari.

Una seconda funzione, \funcd{getgrouplist}, può invece essere usata per
ottenere tutti i gruppi a cui appartiene utente identificato per nome; il suo
prototipo è:

\begin{funcproto}{ 
\fhead{grp.h}
\fdecl{int getgrouplist(const char *user, gid\_t group, gid\_t *groups, int
  *ngroups)} 
\fdesc{Legge i gruppi cui appartiene un utente.} 
}
{La funzione ritorna il numero di gruppi ottenuto in caso di successo e $-1$
  per un errore, che avviene solo quando il numero di gruppi è maggiore di
  quelli specificati con \param{ngroups}.}
\end{funcproto}

La funzione esegue una scansione del database dei gruppi (si veda
sez.~\ref{sec:sys_user_group}) per leggere i gruppi supplementari dell'utente
specificato per nome (e non con un \ids{UID}) nella stringa passata con
l'argomento \param{user}. Ritorna poi nel vettore \param{groups} la lista dei
\ids{GID} dei gruppi a cui l'utente appartiene. Si noti che \param{ngroups},
che in ingresso deve indicare la dimensione di \param{group}, è passato come
\textit{value result argument} perché, qualora il valore specificato sia
troppo piccolo, la funzione ritorna $-1$, passando comunque indietro il numero
dei gruppi trovati, in modo da poter ripetere la chiamata con un vettore di
dimensioni adeguate.

Infine per impostare i gruppi supplementari di un processo ci sono due
funzioni, che possono essere usate solo se si hanno i privilegi di
amministratore.\footnote{e più precisamente se si ha la \textit{capability}
  \const{CAP\_SETGID}.} La prima delle due è la funzione di sistema
\funcd{setgroups},\footnote{la funzione è definita in BSD e SRv4, ma a
  differenza di \func{getgroups} non è stata inclusa in POSIX.1-2001, per
  poterla utilizzare deve essere definita la macro \macro{\_BSD\_SOURCE}.} ed
il suo prototipo è:

\begin{funcproto}{ 
\fhead{grp.h}
\fdecl{int setgroups(size\_t size, gid\_t *list)}
\fdesc{Imposta i gruppi supplementari del processo.} 
}
{La funzione ritorna $0$ in caso di successo e $-1$ per un errore, nel qual
caso \var{errno} assumerà uno dei valori:
\begin{errlist}
\item[\errcode{EFAULT}] \param{list} non ha un indirizzo valido.
\item[\errcode{EINVAL}] il valore di \param{size} è maggiore del valore
    massimo consentito di gruppi supplementari.
\item[\errcode{EPERM}] il processo non ha i privilegi di amministratore.
\end{errlist}}
\end{funcproto}

La funzione imposta i gruppi supplementari del processo corrente ai valori
specificati nel vettore passato con l'argomento \param{list}, di dimensioni
date dall'argomento \param{size}. Il numero massimo di gruppi supplementari
che si possono impostare è un parametro di sistema, che può essere ricavato
con le modalità spiegate in sez.~\ref{sec:sys_characteristics}.

Se invece si vogliono impostare i gruppi supplementari del processo a quelli
di un utente specifico, si può usare la funzione \funcd{initgroups} il cui
prototipo è:

\begin{funcproto}{ 
\fhead{sys/types.h}
\fhead{grp.h}
\fdecl{int initgroups(const char *user, gid\_t group)}
\fdesc{Inizializza la lista dei gruppi supplementari.} 
}
{La funzione ritorna $0$ in caso di successo e $-1$ per un errore, nel qual
caso \var{errno} assumerà uno dei valori:
\begin{errlist}
\item[\errcode{ENOMEM}] non c'è memoria sufficiente per allocare lo spazio per
  informazioni dei gruppi.
\item[\errcode{EPERM}] il processo non ha i privilegi di amministratore.
\end{errlist}}
\end{funcproto}

La funzione esegue la scansione del database dei gruppi (usualmente
\conffile{/etc/group}) cercando i gruppi di cui è membro l'utente \param{user}
(di nuovo specificato per nome e non per \ids{UID}) con cui costruisce una
lista di gruppi supplementari, a cui aggiunge anche \param{group}, infine
imposta questa lista per il processo corrente usando \func{setgroups}.  Si
tenga presente che sia \func{setgroups} che \func{initgroups} non sono
definite nello standard POSIX.1 e che pertanto non è possibile utilizzarle
quando si definisce \macro{\_POSIX\_SOURCE} o si compila con il flag
\cmd{-ansi}, è pertanto meglio evitarle se si vuole scrivere codice portabile.

 
\section{La gestione della priorità dei processi}
\label{sec:proc_priority}

In questa sezione tratteremo più approfonditamente i meccanismi con il quale
lo \textit{scheduler} assegna la CPU ai vari processi attivi.  In particolare
prenderemo in esame i vari meccanismi con cui viene gestita l'assegnazione del
tempo di CPU, ed illustreremo le varie funzioni di gestione. Tratteremo infine
anche le altre priorità dei processi (come quelle per l'accesso a disco)
divenute disponibili con i kernel più recenti.


\subsection{I meccanismi di \textit{scheduling}}
\label{sec:proc_sched}

\itindbeg{scheduler}

La scelta di un meccanismo che sia in grado di distribuire in maniera efficace
il tempo di CPU per l'esecuzione dei processi è sempre una questione delicata,
ed oggetto di numerose ricerche; in generale essa dipende in maniera
essenziale anche dal tipo di utilizzo che deve essere fatto del sistema, per
cui non esiste un meccanismo che sia valido per tutti gli usi.

La caratteristica specifica di un sistema \textit{multitasking} come Linux è
quella del cosiddetto \itindex{preemptive~multitasking} \textit{preemptive
  multitasking}: questo significa che al contrario di altri sistemi (che usano
invece il cosiddetto \itindex{cooperative~multitasking} \textit{cooperative
  multitasking}) non sono i singoli processi, ma il kernel stesso a decidere
quando la CPU deve essere passata ad un altro processo. Come accennato in
sez.~\ref{sec:proc_hierarchy} questa scelta viene eseguita da una sezione
apposita del kernel, lo \textit{scheduler}, il cui scopo è quello di
distribuire al meglio il tempo di CPU fra i vari processi.

La cosa è resa ancora più complicata dal fatto che con le architetture
multi-processore si deve anche scegliere quale sia la CPU più opportuna da
utilizzare.\footnote{nei processori moderni la presenza di ampie cache può
  rendere poco efficiente trasferire l'esecuzione di un processo da una CPU ad
  un'altra, per cui effettuare la migliore scelta fra le diverse CPU non è
  banale.}  Tutto questo comunque appartiene alle sottigliezze
dell'implementazione del kernel; dal punto di vista dei programmi che girano
in \textit{user space}, anche quando si hanno più processori (e dei processi
che sono eseguiti davvero in contemporanea), le politiche di
\textit{scheduling} riguardano semplicemente l'allocazione della risorsa
\textsl{tempo di esecuzione}, la cui assegnazione sarà governata dai
meccanismi di scelta delle priorità che restano gli stessi indipendentemente
dal numero di processori.

Si tenga conto poi che i processi non devono solo eseguire del codice: ad
esempio molto spesso saranno impegnati in operazioni di I/O, o potranno
venire bloccati da un comando dal terminale, o sospesi per un certo periodo di
tempo.  In tutti questi casi la CPU diventa disponibile ed è compito dello
kernel provvedere a mettere in esecuzione un altro processo.

Tutte queste possibilità sono caratterizzate da un diverso \textsl{stato} del
processo; in Linux un processo può trovarsi in uno degli stati riportati in
tab.~\ref{tab:proc_proc_states}; ma soltanto i processi che sono nello stato
\textit{runnable} concorrono per l'esecuzione. Questo vuol dire che, qualunque
sia la sua priorità, un processo non potrà mai essere messo in esecuzione
fintanto che esso si trova in uno qualunque degli altri stati.

\begin{table}[htb]
  \footnotesize
  \centering
  \begin{tabular}[c]{|p{2.4cm}|c|p{9cm}|}
    \hline
    \textbf{Stato} & \texttt{STAT} & \textbf{Descrizione} \\
    \hline
    \hline
    \textit{runnable}& \texttt{R} & Il processo è in esecuzione o è pronto ad
                                    essere eseguito (in attesa che gli
                                    venga assegnata la CPU).\\
    \textit{sleep}   & \texttt{S} & Il processo  è in attesa di un
                                    risposta dal sistema, ma può essere 
                                    interrotto da un segnale.\\
    \textit{uninterrutible sleep}& \texttt{D} & Il  processo è in
                                    attesa di un risposta dal sistema (in 
                                    genere per I/O), e non può essere
                                    interrotto in nessuna circostanza.\\
    \textit{stopped} & \texttt{T} & Il processo è stato fermato con un
                                    \signal{SIGSTOP}, o è tracciato.\\
    \textit{zombie}  & \texttt{Z} & Il processo è terminato ma il
                                    suo stato di terminazione non è ancora
                                    stato letto dal padre.\\
    \textit{killable}& \texttt{D} & Un nuovo stato introdotto con il kernel
                                    2.6.25, sostanzialmente identico
                                    all'\textit{uninterrutible sleep} con la
                                    sola differenza che il processo può
                                    terminato con \signal{SIGKILL} (usato per
                                    lo più per NFS).\\ 
    \hline
  \end{tabular}
  \caption{Elenco dei possibili stati di un processo in Linux, nella colonna
    \texttt{STAT} si è riportata la corrispondente lettera usata dal comando 
    \cmd{ps} nell'omonimo campo.}
  \label{tab:proc_proc_states}
\end{table}

Si deve quindi tenere presente che l'utilizzo della CPU è soltanto una delle
risorse che sono necessarie per l'esecuzione di un programma, e a seconda
dello scopo del programma non è detto neanche che sia la più importante, dato
che molti programmi dipendono in maniera molto più critica dall'I/O. Per
questo motivo non è affatto detto che dare ad un programma la massima priorità
di esecuzione abbia risultati significativi in termini di prestazioni.

Il meccanismo tradizionale di \textit{scheduling} di Unix (che tratteremo in
sez.~\ref{sec:proc_sched_stand}) è sempre stato basato su delle
\textsl{priorità dinamiche}, in modo da assicurare che tutti i processi, anche
i meno importanti, potessero ricevere un po' di tempo di CPU. In sostanza
quando un processo ottiene la CPU la sua priorità viene diminuita. In questo
modo alla fine, anche un processo con priorità iniziale molto bassa, finisce
per avere una priorità sufficiente per essere eseguito.

Lo standard POSIX.1b però ha introdotto il concetto di \textsl{priorità
  assoluta}, (chiamata anche \textsl{priorità statica}, in contrapposizione
alla normale priorità dinamica), per tenere conto dei sistemi
\textit{real-time},\footnote{per sistema \textit{real-time} si intende un
  sistema in grado di eseguire operazioni in un tempo ben determinato; in
  genere si tende a distinguere fra l'\textit{hard real-time} in cui è
  necessario che i tempi di esecuzione di un programma siano determinabili con
  certezza assoluta (come nel caso di meccanismi di controllo di macchine,
  dove uno sforamento dei tempi avrebbe conseguenze disastrose), e
  \textit{soft-real-time} in cui un occasionale sforamento è ritenuto
  accettabile.} in cui è vitale che i processi che devono essere eseguiti in
un determinato momento non debbano aspettare la conclusione di altri che non
hanno questa necessità.

Il concetto di priorità assoluta dice che quando due processi si contendono
l'esecuzione, vince sempre quello con la priorità assoluta più alta.
Ovviamente questo avviene solo per i processi che sono pronti per essere
eseguiti (cioè nello stato \textit{runnable}).  La priorità assoluta viene in
genere indicata con un numero intero, ed un valore più alto comporta una
priorità maggiore. Su questa politica di \textit{scheduling} torneremo in
sez.~\ref{sec:proc_real_time}.

In generale quello che succede in tutti gli Unix moderni è che ai processi
normali viene sempre data una priorità assoluta pari a zero, e la decisione di
assegnazione della CPU è fatta solo con il meccanismo tradizionale della
priorità dinamica. In Linux tuttavia è possibile assegnare anche una priorità
assoluta, nel qual caso un processo avrà la precedenza su tutti gli altri di
priorità inferiore, che saranno eseguiti solo quando quest'ultimo non avrà
bisogno della CPU.


\subsection{Il meccanismo di \textit{scheduling} standard}
\label{sec:proc_sched_stand}

A meno che non si abbiano esigenze specifiche,\footnote{per alcune delle quali
  sono state introdotte delle varianti specifiche.} l'unico meccanismo di
\textit{scheduling} con il quale si avrà a che fare è quello tradizionale, che
prevede solo priorità dinamiche. È di questo che, di norma, ci si dovrà
preoccupare nella programmazione.  Come accennato in Linux i processi ordinari
hanno tutti una priorità assoluta nulla; quello che determina quale, fra tutti
i processi in attesa di esecuzione, sarà eseguito per primo, è la cosiddetta
\textsl{priorità dinamica}, quella che viene mostrata nella colonna
\texttt{PR} del comando \texttt{top}, che è chiamata così proprio perché varia
nel corso dell'esecuzione di un processo.

Il meccanismo usato da Linux è in realtà piuttosto complesso,\footnote{e
  dipende strettamente dalla versione di kernel; in particolare a partire
  dalla serie 2.6.x lo \textit{scheduler} è stato riscritto completamente, con
  molte modifiche susseguitesi per migliorarne le prestazioni, per un certo
  periodo ed è stata anche introdotta la possibilità di usare diversi
  algoritmi, selezionabili sia in fase di compilazione, che, nelle versioni
  più recenti, all'avvio (addirittura è stato ideato un sistema modulare che
  permette di cambiare lo \textit{scheduler} a sistema attivo).} ma a grandi
linee si può dire che ad ogni processo è assegnata una \textit{time-slice},
cioè un intervallo di tempo (letteralmente una fetta) per il quale, a meno di
eventi esterni, esso viene eseguito senza essere interrotto.  Inoltre la
priorità dinamica viene calcolata dallo \textit{scheduler} a partire da un
valore iniziale che viene \textsl{diminuito} tutte le volte che un processo è
in stato \textit{runnable} ma non viene posto in esecuzione.\footnote{in
  realtà il calcolo della priorità dinamica e la conseguente scelta di quale
  processo mettere in esecuzione avviene con un algoritmo molto più
  complicato, che tiene conto anche della \textsl{interattività} del processo,
  utilizzando diversi fattori, questa è una brutale semplificazione per
  rendere l'idea del funzionamento, per una trattazione più dettagliata dei
  meccanismi di funzionamento dello \textit{scheduler}, anche se non
  aggiornatissima, si legga il quarto capitolo di \cite{LinKernDev}.}

Lo \textit{scheduler} infatti mette sempre in esecuzione, fra tutti i processi
in stato \textit{runnable}, quello che ha il valore di priorità dinamica più
basso; con le priorità dinamiche il significato del valore numerico ad esse
associato è infatti invertito, un valore più basso significa una priorità
maggiore. Il fatto che questo valore venga diminuito quando un processo non
viene posto in esecuzione pur essendo pronto, significa che la priorità dei
processi che non ottengono l'uso del processore viene progressivamente
incrementata, così che anche questi alla fine hanno la possibilità di essere
eseguiti.

Sia la dimensione della \textit{time-slice} che il valore di partenza della
priorità dinamica sono determinate dalla cosiddetta \textit{nice} (o
\textit{niceness}) del processo.\footnote{questa è una delle tante proprietà
  che ciascun processo si porta dietro, essa viene ereditata dai processi
  figli e mantenuta attraverso una \func{exec}; fino alla serie 2.4 essa era
  mantenuta nell'omonimo campo \texttt{nice} della \texttt{task\_struct}, con
  la riscrittura dello \textit{scheduler} eseguita nel 2.6 viene mantenuta nel
  campo \texttt{static\_prio} come per le priorità statiche.} L'origine del
nome di questo parametro sta nel fatto che generalmente questo viene usato per
\textsl{diminuire} la priorità di un processo, come misura di cortesia nei
confronti degli altri.  I processi infatti vengono creati dal sistema con un
valore nullo e nessuno è privilegiato rispetto agli altri. Specificando un
valore di \textit{nice} positivo si avrà una \textit{time-slice} più breve ed
un valore di priorità dinamica iniziale più alto, mentre un valore negativo
darà una \textit{time-slice} più lunga ed un valore di priorità dinamica
iniziale più basso.

Esistono diverse funzioni che consentono di indicare un valore di
\textit{nice} di un processo; la più semplice è \funcd{nice}, che opera sul
processo corrente, il suo prototipo è:

\begin{funcproto}{ 
\fhead{unistd.h}
\fdecl{int nice(int inc)}
\fdesc{Aumenta il valore di \textit{nice} del processo corrente.} 
}
{La funzione ritorna il nuovo valore di \textit{nice} in caso di successo e
  $-1$ per un errore, nel qual caso \var{errno} assumerà uno dei valori:
\begin{errlist}
  \item[\errcode{EPERM}] non si ha il permesso di specificare un valore
    di \param{inc} negativo. 
\end{errlist}}
\end{funcproto}

\constbeg{PRIO\_MIN}
\constbeg{PRIO\_MAX}

L'argomento \param{inc} indica l'incremento da effettuare rispetto al valore
di \textit{nice} corrente, che può assumere valori compresi fra
\const{PRIO\_MIN} e \const{PRIO\_MAX}; nel caso di Linux sono fra $-20$ e
$19$,\footnote{in realtà l'intervallo varia a seconda delle versioni di
  kernel, ed è questo a partire dal kernel 1.3.43, anche se oggi si può avere
  anche l'intervallo fra $-20$ e $20$.} ma per \param{inc} si può specificare
un valore qualunque, positivo o negativo, ed il sistema provvederà a troncare
il risultato nell'intervallo consentito. Valori positivi comportano maggiore
\textit{cortesia} e cioè una diminuzione della priorità, valori negativi
comportano invece un aumento della priorità. Con i kernel precedenti il 2.6.12
solo l'amministratore\footnote{o un processo con la \textit{capability}
  \const{CAP\_SYS\_NICE}, vedi sez.~\ref{sec:proc_capabilities}.} può
specificare valori negativi di \param{inc} che permettono di aumentare la
priorità di un processo, a partire da questa versione è consentito anche agli
utenti normali alzare (entro certi limiti, che vedremo in
sez.~\ref{sec:sys_resource_limit}) la priorità dei propri processi.

\constend{PRIO\_MIN}
\constend{PRIO\_MAX}

Gli standard SUSv2 e POSIX.1 prevedono che la funzione ritorni il nuovo valore
di \textit{nice} del processo; tuttavia la \textit{system call} di Linux non
segue questa convenzione e restituisce sempre $0$ in caso di successo e $-1$
in caso di errore; questo perché $-1$ è anche un valore di \textit{nice}
legittimo e questo comporta una confusione con una eventuale condizione di
errore. La \textit{system call} originaria inoltre non consente, se non dotati
di adeguati privilegi, di diminuire un valore di \textit{nice} precedentemente
innalzato.
 
Fino alla \acr{glibc} 2.2.4 la funzione di libreria riportava direttamente il
risultato dalla \textit{system call}, violando lo standard, per cui per
ottenere il nuovo valore occorreva una successiva chiamata alla funzione
\func{getpriority}. A partire dalla \acr{glibc} 2.2.4 \func{nice} è stata
reimplementata e non viene più chiamata la omonima \textit{system call}, con
questa versione viene restituito come valore di ritorno il valore di
\textit{nice}, come richiesto dallo standard.\footnote{questo viene fatto
  chiamando al suo interno \func{setpriority}, che tratteremo a breve.}  In
questo caso l'unico modo per rilevare in maniera affidabile una condizione di
errore è quello di azzerare \var{errno} prima della chiamata della funzione e
verificarne il valore quando \func{nice} restituisce $-1$.

Per leggere il valore di \textit{nice} di un processo occorre usare la
funzione di sistema \funcd{getpriority}, derivata da BSD; il suo prototipo è:

\begin{funcproto}{ 
\fhead{sys/time.h}
\fhead{sys/resource.h}
\fdecl{int getpriority(int which, int who)}
\fdesc{Legge un valore di \textit{nice}.} 
}
{La funzione ritorna $0$ in caso di successo e $-1$ per un errore, nel qual
caso \var{errno} assumerà uno dei valori:
\begin{errlist}
\item[\errcode{EINVAL}] il valore di \param{which} non è uno di quelli
    elencati in tab.~\ref{tab:proc_getpriority}.
\item[\errcode{ESRCH}] non c'è nessun processo che corrisponda ai valori di
  \param{which} e \param{who}.
\end{errlist}}
\end{funcproto}

La funzione permette, a seconda di quanto specificato nell'argomento
\param{which}, di leggere il valore di \textit{nice} o di un processo, o di un
gruppo di processi (vedi sez.~\ref{sec:sess_proc_group}) o di un utente,
indicati con l'argomento \param{who}. Nelle vecchie versioni può essere
necessario includere anche \headfiled{sys/time.h}, questo non è più necessario
con versioni recenti delle librerie, ma è comunque utile per portabilità.

I valori possibili per \param{which}, ed il tipo di valore che occorre usare
in corrispondenza per \param{who}, solo elencati nella legenda di
tab.~\ref{tab:proc_getpriority} insieme ai relativi significati. Usare un
valore nullo per \param{who} indica, a seconda della corrispondente
indicazione usata per \param{which}, il processo, il gruppo di processi o
l'utente correnti.

\begin{table}[htb]
  \centering
  \footnotesize
  \begin{tabular}[c]{|c|c|l|}
    \hline
    \param{which} & \param{who} & \textbf{Significato} \\
    \hline
    \hline
    \constd{PRIO\_PROCESS} & \type{pid\_t} & processo  \\
    \constd{PRIO\_PRGR}    & \type{pid\_t} & \textit{process group} (vedi
                                             sez.~\ref{sec:sess_proc_group})\\
    \constd{PRIO\_USER}    & \type{uid\_t} & utente \\
    \hline
  \end{tabular}
  \caption{Legenda del valore dell'argomento \param{which} e del tipo
    dell'argomento \param{who} delle funzioni \func{getpriority} e
    \func{setpriority} per le tre possibili scelte.}
  \label{tab:proc_getpriority}
\end{table}

In caso di una indicazione che faccia riferimento a più processi, la funzione
restituisce la priorità più alta (cioè il valore più basso) fra quelle dei
processi corrispondenti. Come per \func{nice}, $-1$ è un possibile valore
corretto, per cui di nuovo per poter rilevare una condizione di errore è
necessario cancellare sempre \var{errno} prima della chiamata alla funzione e
quando si ottiene un valore di ritorno uguale a $-1$ per verificare che essa
resti uguale a zero.

Analoga a \func{getpriority} è la funzione di sistema \funcd{setpriority} che
permette di impostare la priorità di uno o più processi; il suo prototipo è:

\begin{funcproto}{ 
\fhead{sys/time.h}
\fhead{sys/resource.h}
\fdecl{int setpriority(int which, int who, int prio)}
\fdesc{Imposta un valore di \textit{nice}.} 
}
{La funzione ritorna $0$ in caso di successo e $-1$ per un errore, nel qual
caso \var{errno} assumerà uno dei valori:
\begin{errlist}
\item[\errcode{EACCES}] si è richiesto un aumento di priorità senza avere
  sufficienti privilegi.
\item[\errcode{EINVAL}] il valore di \param{which} non è uno di quelli
  elencati in tab.~\ref{tab:proc_getpriority}.
\item[\errcode{EPERM}] un processo senza i privilegi di amministratore ha
  cercato di modificare la priorità di un processo di un altro utente.
\item[\errcode{ESRCH}] non c'è nessun processo che corrisponda ai valori di
  \param{which} e \param{who}.
\end{errlist}}
\end{funcproto}

La funzione imposta la priorità dinamica al valore specificato da \param{prio}
per tutti i processi indicati dagli argomenti \param{which} e \param{who}, per
i quali valgono le stesse considerazioni fatte per \func{getpriority} e lo
specchietto di tab.~\ref{tab:proc_getpriority}. 

In questo caso come valore di \param{prio} deve essere specificato il valore
di \textit{nice} da assegnare nell'intervallo fra \const{PRIO\_MIN} ($-20$) e
\const{PRIO\_MAX} ($19$), e non un incremento (positivo o negativo) come nel
caso di \func{nice}. La funzione restituisce il valore di \textit{nice}
assegnato in caso di successo e $-1$ in caso di errore, e come per \func{nice}
anche in questo caso per rilevare un errore occorre sempre porre a zero
\var{errno} prima della chiamata della funzione, essendo $-1$ un valore di
\textit{nice} valido.

Si tenga presente che solo l'amministratore\footnote{o più precisamente un
  processo con la \textit{capability} \const{CAP\_SYS\_NICE}, vedi
  sez.~\ref{sec:proc_capabilities}.} ha la possibilità di modificare
arbitrariamente le priorità di qualunque processo. Un utente normale infatti
può modificare solo la priorità dei suoi processi ed in genere soltanto
diminuirla.  Fino alla versione di kernel 2.6.12 Linux ha seguito le
specifiche dello standard SUSv3, e come per tutti i sistemi derivati da SysV
veniva richiesto che l'\ids{UID} reale o quello effettivo del processo
chiamante corrispondessero all'\ids{UID} reale (e solo a quello) del processo
di cui si intendeva cambiare la priorità. A partire dalla versione 2.6.12 è
stata adottata la semantica in uso presso i sistemi derivati da BSD (SunOS,
Ultrix, *BSD), in cui la corrispondenza può essere anche con l'\ids{UID}
effettivo.

Sempre a partire dal kernel 2.6.12 è divenuto possibile anche per gli utenti
ordinari poter aumentare la priorità dei propri processi specificando un
valore di \param{prio} negativo. Questa operazione non è possibile però in
maniera indiscriminata, ed in particolare può essere effettuata solo
nell'intervallo consentito dal valore del limite \const{RLIMIT\_NICE}
(torneremo su questo in sez.~\ref{sec:sys_resource_limit}).

Infine nonostante i valori siano sempre rimasti gli stessi, il significato del
valore di \textit{nice} è cambiato parecchio nelle progressive riscritture
dello \textit{scheduler} di Linux, ed in particolare a partire dal kernel
2.6.23 l'uso di diversi valori di \textit{nice} ha un impatto molto più forte
nella distribuzione della CPU ai processi. Infatti se viene comunque calcolata
una priorità dinamica per i processi che non ricevono la CPU, così che anche
essi possano essere messi in esecuzione, un alto valore di \textit{nice}
corrisponde comunque ad una \textit{time-slice} molto piccola che non cresce
comunque, per cui un processo a bassa priorità avrà davvero scarse possibilità
di essere eseguito in presenza di processi attivi a priorità più alta.


\subsection{Il meccanismo di \textit{scheduling real-time}}
\label{sec:proc_real_time}

Come spiegato in sez.~\ref{sec:proc_sched} lo standard POSIX.1b ha introdotto
le priorità assolute per permettere la gestione di processi
\textit{real-time}. In realtà nel caso di Linux non si tratta di un vero
\textit{hard real-time}, in quanto in presenza di eventuali interrupt il
kernel interrompe l'esecuzione di un processo, qualsiasi sia la sua
priorità,\footnote{questo a meno che non si siano installate le patch di
  RTLinux, RTAI o Adeos, con i quali è possibile ottenere un sistema
  effettivamente \textit{hard real-time}. In tal caso infatti gli interrupt
  vengono intercettati dall'interfaccia \textit{real-time} (o nel caso di
  Adeos gestiti dalle code del nano-kernel), in modo da poterli controllare
  direttamente qualora ci sia la necessità di avere un processo con priorità
  più elevata di un \textit{interrupt handler}.} mentre con l'incorrere in un
\textit{page fault} si possono avere ritardi non previsti.  Se l'ultimo
problema può essere aggirato attraverso l'uso delle funzioni di controllo
della memoria virtuale (vedi sez.~\ref{sec:proc_mem_lock}), il primo non è
superabile e può comportare ritardi non prevedibili riguardo ai tempi di
esecuzione di qualunque processo.

Nonostante questo, ed in particolare con una serie di miglioramenti che sono
stati introdotti nello sviluppo del kernel,\footnote{in particolare a partire
  dalla versione 2.6.18 sono stati inserite nel kernel una serie di modifiche
  che consentono di avvicinarsi sempre di più ad un vero e proprio sistema
  \textit{real-time} estendendo il concetto di \textit{preemption} alle
  operazioni dello stesso kernel; esistono vari livelli a cui questo può
  essere fatto, ottenibili attivando in fase di compilazione una fra le
  opzioni \texttt{CONFIG\_PREEMPT\_NONE}, \texttt{CONFIG\_PREEMPT\_VOLUNTARY}
  e \texttt{CONFIG\_PREEMPT\_DESKTOP}.} si può arrivare ad una ottima
approssimazione di sistema \textit{real-time} usando le priorità assolute.
Occorre farlo però con molta attenzione: se si dà ad un processo una priorità
assoluta e questo finisce in un loop infinito, nessun altro processo potrà
essere eseguito, ed esso sarà mantenuto in esecuzione permanentemente
assorbendo tutta la CPU e senza nessuna possibilità di riottenere l'accesso al
sistema. Per questo motivo è sempre opportuno, quando si lavora con processi
che usano priorità assolute, tenere attiva una shell cui si sia assegnata la
massima priorità assoluta, in modo da poter essere comunque in grado di
rientrare nel sistema.

Quando c'è un processo con priorità assoluta lo \textit{scheduler} lo metterà
in esecuzione prima di ogni processo normale. In caso di più processi sarà
eseguito per primo quello con priorità assoluta più alta. Quando ci sono più
processi con la stessa priorità assoluta questi vengono tenuti in una coda e
tocca al kernel decidere quale deve essere eseguito.  Il meccanismo con cui
vengono gestiti questi processi dipende dalla politica di \textit{scheduling}
che si è scelta; lo standard ne prevede due:
\begin{basedescript}{\desclabelwidth{1.2cm}\desclabelstyle{\nextlinelabel}}
\item[\textit{First In First Out} (FIFO)] Il processo viene eseguito
  fintanto che non cede volontariamente la CPU (con la funzione
  \func{sched\_yield}), si blocca, finisce o viene interrotto da un processo a
  priorità più alta. Se il processo viene interrotto da uno a priorità più
  alta esso resterà in cima alla lista e sarà il primo ad essere eseguito
  quando i processi a priorità più alta diverranno inattivi. Se invece lo si
  blocca volontariamente sarà posto in coda alla lista (ed altri processi con
  la stessa priorità potranno essere eseguiti).
\item[\textit{Round Robin} (RR)] Il comportamento è del tutto analogo a quello
  precedente, con la sola differenza che ciascun processo viene eseguito al
  massimo per un certo periodo di tempo (la cosiddetta \textit{time-slice})
  dopo di che viene automaticamente posto in fondo alla coda dei processi con
  la stessa priorità. In questo modo si ha comunque una esecuzione a turno di
  tutti i processi, da cui il nome della politica. Solo i processi con la
  stessa priorità ed in stato \textit{runnable} entrano nel
  \textsl{girotondo}.
\end{basedescript}

Lo standard POSIX.1-2001 prevede una funzione che consenta sia di modificare
le politiche di \textit{scheduling}, passando da \textit{real-time} a
ordinarie o viceversa, che di specificare, in caso di politiche
\textit{real-time}, la eventuale priorità statica; la funzione di sistema è
\funcd{sched\_setscheduler} ed il suo prototipo è:

\begin{funcproto}{ 
\fhead{sched.h}
\fdecl{int sched\_setscheduler(pid\_t pid, int policy, const struct
  sched\_param *p)}
\fdesc{Imposta priorità e politica di \textit{scheduling}.} 
}
{La funzione ritorna $0$ in caso di successo e $-1$ per un errore, nel qual
caso \var{errno} assumerà uno dei valori:
\begin{errlist}
    \item[\errcode{EINVAL}] il valore di \param{policy} non esiste o il
      valore di \param{p} non è valido per la politica scelta.
    \item[\errcode{EPERM}] il processo non ha i privilegi per attivare la
      politica richiesta.
    \item[\errcode{ESRCH}] il processo \param{pid} non esiste.
    \end{errlist}}
\end{funcproto}

La funzione esegue l'impostazione per il processo specificato dall'argomento
\param{pid}; un valore nullo di questo argomento esegue l'impostazione per il
processo corrente.  La politica di \textit{scheduling} è specificata
dall'argomento \param{policy} i cui possibili valori sono riportati in
tab.~\ref{tab:proc_sched_policy}; la parte alta della tabella indica le
politiche \textit{real-time}, quella bassa le politiche ordinarie. Un valore
negativo per \param{policy} mantiene la politica di \textit{scheduling}
corrente.

\begin{table}[htb]
  \centering
  \footnotesize
  \begin{tabular}[c]{|l|p{6cm}|}
    \hline
    \textbf{Politica}  & \textbf{Significato} \\
    \hline
    \hline
    \constd{SCHED\_FIFO} & \textit{Scheduling real-time} con politica
                           \textit{FIFO}. \\
    \constd{SCHED\_RR}   & \textit{Scheduling real-time} con politica
                           \textit{Round Robin}. \\ 
    \hline
    \constd{SCHED\_OTHER}& \textit{Scheduling} ordinario.\\
    \constd{SCHED\_BATCH}& \textit{Scheduling} ordinario con l'assunzione
                           ulteriore di lavoro \textit{CPU
                           intensive} (dal kernel 2.6.16).\\ 
    \constd{SCHED\_IDLE} & \textit{Scheduling} di priorità estremamente
                           bassa (dal kernel 2.6.23).\\
    \hline
  \end{tabular}
  \caption{Valori dell'argomento \param{policy} per la funzione
    \func{sched\_setscheduler}.}
  \label{tab:proc_sched_policy}
\end{table}

% TODO Aggiungere SCHED_DEADLINE, sulla nuova politica di scheduling aggiunta
% con il kernel 3.14, vedi anche Documentation/scheduler/sched-deadline.txt e
% http://lwn.net/Articles/575497/
% vedi anche man 7 sched, man sched_setattr

Con le versioni più recenti del kernel sono state introdotte anche delle
varianti sulla politica di \textit{scheduling} tradizionale per alcuni carichi
di lavoro specifici, queste due nuove politiche sono specifiche di Linux e non
devono essere usate se si vogliono scrivere programmi portabili.

La politica \const{SCHED\_BATCH} è una variante della politica ordinaria con
la sola differenza che i processi ad essa soggetti non ottengono, nel calcolo
delle priorità dinamiche fatto dallo \textit{scheduler}, il cosiddetto bonus
di interattività che mira a favorire i processi che si svegliano dallo stato
di \textit{sleep}.\footnote{cosa che accade con grande frequenza per i
  processi interattivi, dato che essi sono per la maggior parte del tempo in
  attesa di dati in ingresso da parte dell'utente.} La si usa pertanto, come
indica il nome, per processi che usano molta CPU (come programmi di calcolo)
che in questo modo, pur non perdendo il loro valore di \textit{nice}, sono
leggermente sfavoriti rispetto ai processi interattivi che devono rispondere a
dei dati in ingresso.

La politica \const{SCHED\_IDLE} invece è una politica dedicata ai processi che
si desidera siano eseguiti con la più bassa priorità possibile, ancora più
bassa di un processo con il minimo valore di \textit{nice}. In sostanza la si
può utilizzare per processi che devono essere eseguiti se non c'è niente altro
da fare. Va comunque sottolineato che anche un processo \const{SCHED\_IDLE}
avrà comunque una sua possibilità di utilizzo della CPU, sia pure in
percentuale molto bassa.

Qualora si sia richiesta una politica \textit{real-time} il valore della
priorità statica viene impostato attraverso la struttura
\struct{sched\_param}, riportata in fig.~\ref{fig:sig_sched_param}, il cui
solo campo attualmente definito è \var{sched\_priority}. Il campo deve
contenere il valore della priorità statica da assegnare al processo; lo
standard prevede che questo debba essere assegnato all'interno di un
intervallo fra un massimo ed un minimo che nel caso di Linux sono
rispettivamente 1 e 99.

\begin{figure}[!htb]
  \footnotesize \centering
  \begin{minipage}[c]{0.5\textwidth}
    \includestruct{listati/sched_param.c}
  \end{minipage} 
  \normalsize 
  \caption{La struttura \structd{sched\_param}.} 
  \label{fig:sig_sched_param}
\end{figure}

I processi con politica di \textit{scheduling} ordinaria devono sempre
specificare un valore nullo di \var{sched\_priority} altrimenti si avrà un
errore \errcode{EINVAL}, questo valore infatti non ha niente a che vedere con
la priorità dinamica determinata dal valore di \textit{nice}, che deve essere
impostato con le funzioni viste in precedenza.

Lo standard POSIX.1b prevede che l'intervallo dei valori delle priorità
statiche possa essere ottenuto con le funzioni di sistema
\funcd{sched\_get\_priority\_max} e \funcd{sched\_get\_priority\_min}, i cui
prototipi sono:

\begin{funcproto}{ 
\fhead{sched.h}
\fdecl{int sched\_get\_priority\_max(int policy)}
\fdesc{Legge il valore massimo di una priorità statica.} 
\fdecl{int sched\_get\_priority\_min(int policy)}
\fdesc{Legge il valore minimo di una priorità statica.} 
}
{Le funzioni ritornano il valore della priorità in caso di successo e $-1$ per
  un errore, nel qual caso \var{errno} assumerà il valore:
\begin{errlist}
\item[\errcode{EINVAL}] il valore di \param{policy} non è valido.
\end{errlist}}
\end{funcproto}

Le funzioni ritornano rispettivamente il valore massimo e minimo usabile per
la priorità statica di una delle politiche di \textit{scheduling}
\textit{real-time} indicata dall'argomento \param{policy}.

Si tenga presente che quando si imposta una politica di \textit{scheduling}
real-time per un processo o se ne cambia la priorità statica questo viene
messo in cima alla lista dei processi con la stessa priorità; questo comporta
che verrà eseguito subito, interrompendo eventuali altri processi con la
stessa priorità in quel momento in esecuzione.

Il kernel mantiene i processi con la stessa priorità assoluta in una lista, ed
esegue sempre il primo della lista, mentre un nuovo processo che torna in
stato \textit{runnable} viene sempre inserito in coda alla lista. Se la
politica scelta è \const{SCHED\_FIFO} quando il processo viene eseguito viene
automaticamente rimesso in coda alla lista, e la sua esecuzione continua
fintanto che non viene bloccato da una richiesta di I/O, o non rilascia
volontariamente la CPU (in tal caso, tornando nello stato \textit{runnable}
sarà in coda alla lista); l'esecuzione viene ripresa subito solo nel caso che
esso sia stato interrotto da un processo a priorità più alta.

Solo un processo con i privilegi di amministratore\footnote{più precisamente
  con la capacità \const{CAP\_SYS\_NICE}, vedi
  sez.~\ref{sec:proc_capabilities}.} può impostare senza restrizioni priorità
assolute diverse da zero o politiche \const{SCHED\_FIFO} e
\const{SCHED\_RR}. Un utente normale può modificare solo le priorità di
processi che gli appartengono; è cioè richiesto che l'\ids{UID} effettivo del
processo chiamante corrisponda all'\ids{UID} reale o effettivo del processo
indicato con \param{pid}.

Fino al kernel 2.6.12 gli utenti normali non potevano impostare politiche
\textit{real-time} o modificare la eventuale priorità statica di un loro
processo. A partire da questa versione è divenuto possibile anche per gli
utenti normali usare politiche \textit{real-time} fintanto che la priorità
assoluta che si vuole impostare è inferiore al limite \const{RLIMIT\_RTPRIO}
(vedi sez.~\ref{sec:sys_resource_limit}) ad essi assegnato. 

Unica eccezione a questa possibilità sono i processi \const{SCHED\_IDLE}, che
non possono cambiare politica di \textit{scheduling} indipendentemente dal
valore di \const{RLIMIT\_RTPRIO}. Inoltre, in caso di processo già sottoposto
ad una politica \textit{real-time}, un utente può sempre, indipendentemente
dal valore di \const{RLIMIT\_RTPRIO}, diminuirne la priorità o portarlo ad una
politica ordinaria.

Se si intende operare solo sulla priorità statica di un processo si possono
usare le due funzioni di sistema \funcd{sched\_setparam} e
\funcd{sched\_getparam} che consentono rispettivamente di impostarne e
leggerne il valore, i loro prototipi sono:

\begin{funcproto}{
\fhead{sched.h}
\fdecl{int sched\_setparam(pid\_t pid, const struct sched\_param *param)}
\fdesc{Imposta la priorità statica di un processo.} 
\fdecl{int sched\_getparam(pid\_t pid, struct sched\_param *param)}
\fdesc{Legge la priorità statica di un processo.} 
}
{Le funzioni ritornano $0$ in caso di successo e $-1$ per un errore, nel qual
caso \var{errno} assumerà uno dei valori:
\begin{errlist}
\item[\errcode{EINVAL}] il valore di \param{param} non ha senso per la
  politica usata dal processo.
\item[\errcode{EPERM}] non si hanno privilegi sufficienti per eseguire
  l'operazione.
\item[\errcode{ESRCH}] il processo \param{pid} non esiste.
\end{errlist}}
\end{funcproto}

Le funzioni richiedono di indicare nell'argomento \param{pid} il processo su
cui operare e usano l'argomento \param{param} per mantenere il valore della
priorità dinamica. Questo è ancora una struttura \struct{sched\_param} ed
assume gli stessi valori già visti per \func{sched\_setscheduler}.

L'uso di \func{sched\_setparam}, compresi i controlli di accesso che vi si
applicano, è del tutto equivalente a quello di \func{sched\_setscheduler} con
argomento \param{policy} uguale a $-1$. Come per \func{sched\_setscheduler}
specificando $0$ come valore dell'argomento \param{pid} si opera sul processo
corrente. Benché la funzione sia utilizzabile anche con processi sottoposti a
politica ordinaria essa ha senso soltanto per quelli \textit{real-time}, dato
che per i primi la priorità statica può essere soltanto nulla.  La
disponibilità di entrambe le funzioni può essere verificata controllando la
macro \macro{\_POSIX\_PRIORITY\_SCHEDULING} che è definita nell'\textit{header
  file} \headfiled{sched.h}.

Se invece si vuole sapere quale è politica di \textit{scheduling} di un
processo si può usare la funzione di sistema \funcd{sched\_getscheduler}, il
cui prototipo è:

\begin{funcproto}{ 
\fhead{sched.h}
\fdecl{int sched\_getscheduler(pid\_t pid)}
\fdesc{Legge la politica di \textit{scheduling}.} 
}
{La funzione ritorna la politica di \textit{scheduling}  in caso di successo e
  $-1$ per un errore, nel qual caso \var{errno} assumerà uno dei valori:
\begin{errlist}
    \item[\errcode{EPERM}] non si hanno privilegi sufficienti per eseguire
      l'operazione.
    \item[\errcode{ESRCH}] il processo \param{pid} non esiste.
\end{errlist}}
\end{funcproto}

La funzione restituisce il valore, secondo quanto elencato in
tab.~\ref{tab:proc_sched_policy}, della politica di \textit{scheduling} per il
processo specificato dall'argomento \param{pid}, se questo è nullo viene
restituito il valore relativo al processo chiamante.

L'ultima funzione di sistema che permette di leggere le informazioni relative
ai processi real-time è \funcd{sched\_rr\_get\_interval}, che permette di
ottenere la lunghezza della \textit{time-slice} usata dalla politica
\textit{round robin}; il suo prototipo è:

\begin{funcproto}{ 
\fhead{sched.h}
\fdecl{int sched\_rr\_get\_interval(pid\_t pid, struct timespec *tp)}
\fdesc{Legge la durata della \textit{time-slice} per lo \textit{scheduling}
  \textit{round robin}.}  
}
{La funzione ritorna $0$ in caso di successo e $-1$ per un errore, nel qual
caso \var{errno} assumerà uno dei valori:
\begin{errlist}
\item[\errcode{EINVAL}] l'argomento \param{pid} non è valido. 
\item[\errcode{ENOSYS}] la \textit{system call} non è presente (solo per
  kernel arcaici).
\item[\errcode{ESRCH}] il processo \param{pid} non esiste.
\end{errlist}
ed inoltre anche \errval{EFAULT} nel suo significato generico.}
\end{funcproto}

La funzione restituisce nell'argomento \param{tp} come una struttura
\struct{timespec}, (la cui definizione si può trovare in
fig.~\ref{fig:sys_timeval_struct}) il valore dell'intervallo di tempo usato
per la politica \textit{round robin} dal processo indicato da \ids{PID}. Il
valore dipende dalla versione del kernel, a lungo infatti questo intervallo di
tempo era prefissato e non modificabile ad un valore di 150 millisecondi,
restituito indipendentemente dal \ids{PID} indicato. 

Con kernel recenti però è possibile ottenere una variazione della
\textit{time-slice}, modificando il valore di \textit{nice} del processo
(anche se questo non incide assolutamente sulla priorità statica) che come
accennato in precedenza modifica il valore assegnato alla \textit{time-slice}
di un processo ordinario, che però viene usato anche dai processi
\textit{real-time}.

Come accennato ogni processo può rilasciare volontariamente la CPU in modo da
consentire agli altri processi di essere eseguiti; la funzione di sistema che
consente di fare tutto questo è \funcd{sched\_yield}, il cui prototipo è:

\begin{funcproto}{ 
\fhead{sched.h}
\fdecl{int sched\_yield(void)}
\fdesc{Rilascia volontariamente l'esecuzione.} 
}
{La funzione ritorna $0$ in caso di successo e teoricamente $-1$ per un
  errore, ma su Linux ha sempre successo.}
\end{funcproto}


Questa funzione ha un utilizzo effettivo soltanto quando si usa lo
\textit{scheduling} \textit{real-time}, e serve a far sì che il processo
corrente rilasci la CPU, in modo da essere rimesso in coda alla lista dei
processi con la stessa priorità per permettere ad un altro di essere eseguito;
se però il processo è l'unico ad essere presente sulla coda l'esecuzione non
sarà interrotta. In genere usano questa funzione i processi con politica
\const{SCHED\_FIFO}, per permettere l'esecuzione degli altri processi con pari
priorità quando la sezione più urgente è finita.

La funzione può essere utilizzata anche con processi che usano lo
\textit{scheduling} ordinario, ma in questo caso il comportamento non è ben
definito, e dipende dall'implementazione. Fino al kernel 2.6.23 questo
comportava che i processi venissero messi in fondo alla coda di quelli attivi,
con la possibilità di essere rimessi in esecuzione entro breve tempo, con
l'introduzione del \textit{Completely Fair Scheduler} questo comportamento è
cambiato ed un processo che chiama la funzione viene inserito nella lista dei
processi inattivi, con un tempo molto maggiore.\footnote{è comunque possibile
  ripristinare un comportamento analogo al precedente scrivendo il valore 1
  nel file \sysctlfiled{kernel/sched\_compat\_yield}.}

L'uso delle funzione nella programmazione ordinaria può essere utile e
migliorare le prestazioni generali del sistema quando si è appena rilasciata
una risorsa contesa con altri processi, e si vuole dare agli altri una
possibilità di approfittarne mettendoli in esecuzione, ma chiamarla senza
necessità, specie se questo avviene ripetutamente all'interno di un qualche
ciclo, può avere invece un forte impatto negativo per la generazione di
\textit{context switch} inutili.


\subsection{Il controllo dello \textit{scheduler} per i sistemi
  multiprocessore}
\label{sec:proc_sched_multiprocess}

\index{effetto~ping-pong|(} 

Con il supporto dei sistemi multiprocessore sono state introdotte delle
funzioni che permettono di controllare in maniera più dettagliata la scelta di
quale processore utilizzare per eseguire un certo programma. Uno dei problemi
che si pongono nei sistemi multiprocessore è infatti quello del cosiddetto
\textsl{effetto ping-pong}. Può accadere cioè che lo \textit{scheduler},
quando riavvia un processo precedentemente interrotto scegliendo il primo
processore disponibile, lo faccia eseguire da un processore diverso rispetto a
quello su cui era stato eseguito in precedenza. Se il processo passa da un
processore all'altro in questo modo, cosa che avveniva abbastanza di frequente
con i kernel della seria 2.4.x, si ha l'effetto ping-pong.

Questo tipo di comportamento può generare dei seri problemi di prestazioni;
infatti tutti i processori moderni utilizzano una memoria interna (la
\textit{cache}) contenente i dati più usati, che permette di evitare di
eseguire un accesso (molto più lento) alla memoria principale sulla scheda
madre.  Chiaramente un processo sarà favorito se i suoi dati sono nella cache
del processore, ma è ovvio che questo può essere vero solo per un processore
alla volta, perché in presenza di più copie degli stessi dati su più
processori, non si potrebbe determinare quale di questi ha la versione dei
dati aggiornata rispetto alla memoria principale.

Questo comporta che quando un processore inserisce un dato nella sua cache,
tutti gli altri processori che hanno lo stesso dato devono invalidarlo, e
questa operazione è molto costosa in termini di prestazioni. Il problema
diventa serio quando si verifica l'effetto ping-pong, in tal caso infatti un
processo \textsl{rimbalza} continuamente da un processore all'altro e si ha
una continua invalidazione della cache, che non diventa mai disponibile.

\itindbeg{CPU~affinity}

Per ovviare a questo tipo di problemi è nato il concetto di \textsl{affinità
  di processore} (o \textit{CPU affinity}); la possibilità cioè di far sì che
un processo possa essere assegnato per l'esecuzione sempre allo stesso
processore. Lo \textit{scheduler} dei kernel della serie 2.4.x aveva una
scarsa \textit{CPU affinity}, e l'effetto ping-pong era comune; con il nuovo
\textit{scheduler} dei kernel della 2.6.x questo problema è stato risolto ed
esso cerca di mantenere il più possibile ciascun processo sullo stesso
processore.

\index{effetto~ping-pong|)}

In certi casi però resta l'esigenza di poter essere sicuri che un processo sia
sempre eseguito dallo stesso processore,\footnote{quella che viene detta
  \textit{hard CPU affinity}, in contrasto con quella fornita dallo
  \textit{scheduler}, detta \textit{soft CPU affinity}, che di norma indica
  solo una preferenza, non un requisito assoluto.} e per poter risolvere
questo tipo di problematiche nei nuovi kernel\footnote{le due \textit{system
    call} per la gestione della \textit{CPU affinity} sono state introdotte
  nel kernel 2.5.8, e le corrispondenti funzioni di sistema nella
  \textsl{glibc} 2.3.} è stata introdotta l'opportuna infrastruttura ed una
nuova \textit{system call} che permette di impostare su quali processori far
eseguire un determinato processo attraverso una \textsl{maschera di
  affinità}. La corrispondente funzione di sistema è
\funcd{sched\_setaffinity} ed il suo prototipo è:

\index{insieme~di~processori|(}

\begin{funcproto}{ 
\fhead{sched.h}
\fdecl{int sched\_setaffinity(pid\_t pid, size\_t setsize, 
  cpu\_set\_t *mask)}
\fdesc{Imposta la maschera di affinità di un processo.} 
}
{La funzione ritorna $0$ in caso di successo e $-1$ per un errore, nel qual
caso \var{errno} assumerà uno dei valori:
\begin{errlist}
\item[\errcode{EINVAL}] il valore di \param{mask} contiene riferimenti a
  processori non esistenti nel sistema o a cui non è consentito l'accesso.
\item[\errcode{EPERM}] il processo non ha i privilegi sufficienti per
  eseguire l'operazione.
\item[\errcode{ESRCH}] il processo \param{pid} non esiste.
\end{errlist}
ed inoltre anche \errval{EFAULT} nel suo significato generico.}
\end{funcproto}

Questa funzione e la corrispondente \func{sched\_getaffinity} hanno una storia
abbastanza complessa, la sottostante \textit{system call} infatti prevede
l'uso di due soli argomenti (per il pid e l'indicazione della maschera dei
processori), che corrispondono al fatto che l'implementazione effettiva usa
una semplice maschera binaria. Quando le funzioni vennero incluse nella
\acr{glibc} assunsero invece un prototipo simile a quello mostrato però con il
secondo argomento di tipo \ctyp{unsigned int}. A complicare la cosa si
aggiunge il fatto che nella versione 2.3.3 della \acr{glibc} detto argomento
venne stato eliminato, per poi essere ripristinato nella versione 2.3.4 nella
forma attuale.\footnote{pertanto se la vostra pagina di manuale non è
  aggiornata, o usate quella particolare versione della \acr{glibc}, potrete
  trovare indicazioni diverse, il prototipo illustrato è quello riportato
  nella versione corrente (maggio 2008) delle pagine di manuale e
  corrispondente alla definizione presente in \headfile{sched.h}.}

La funzione imposta, con l'uso del valore contenuto all'indirizzo
\param{mask}, l'insieme dei processori sui quali deve essere eseguito il
processo identificato tramite il valore passato in \param{pid}. Come in
precedenza il valore nullo di \param{pid} indica il processo corrente.  Per
poter utilizzare questa funzione sono richiesti i privilegi di amministratore
(è necessaria la capacità \const{CAP\_SYS\_NICE}) altrimenti essa fallirà con
un errore di \errcode{EPERM}. Una volta impostata una maschera di affinità,
questa viene ereditata attraverso una \func{fork}, in questo modo diventa
possibile legare automaticamente un gruppo di processi ad un singolo
processore.

Nell'uso comune, almeno con i kernel successivi alla serie 2.6.x, utilizzare
questa funzione non è necessario, in quanto è lo \textit{scheduler} stesso che
provvede a mantenere al meglio l'affinità di processore. Esistono però
esigenze particolari, ad esempio quando un processo (o un gruppo di processi)
è utilizzato per un compito importante (ad esempio per applicazioni
\textit{real-time} o la cui risposta è critica) e si vuole la massima
velocità; con questa interfaccia diventa possibile selezionare gruppi di
processori utilizzabili in maniera esclusiva.  Lo stesso dicasi quando
l'accesso a certe risorse (memoria o periferiche) può avere un costo diverso a
seconda del processore, come avviene nelle architetture NUMA
(\textit{Non-Uniform Memory Access}).

Infine se un gruppo di processi accede alle stesse risorse condivise (ad
esempio una applicazione con più \textit{thread}) può avere senso usare lo
stesso processore in modo da sfruttare meglio l'uso della sua cache; questo
ovviamente riduce i benefici di un sistema multiprocessore nell'esecuzione
contemporanea dei \textit{thread}, ma in certi casi (quando i \textit{thread}
sono inerentemente serializzati nell'accesso ad una risorsa) possono esserci
sufficienti vantaggi nell'evitare la perdita della cache da rendere
conveniente l'uso dell'affinità di processore.

Dato che il numero di processori può variare a seconda delle architetture, per
semplificare l'uso dell'argomento \param{mask} la \acr{glibc} ha introdotto un
apposito dato di tipo, \typed{cpu\_set\_t},\footnote{questa è una estensione
  specifica della \acr{glibc}, da attivare definendo la macro
  \macro{\_GNU\_SOURCE}, non esiste infatti una standardizzazione per questo
  tipo di interfaccia e POSIX al momento non prevede nulla al riguardo.} che
permette di identificare un insieme di processori. Il dato è normalmente una
maschera binaria: nei casi più comuni potrebbe bastare un intero a 32 bit, in
cui ogni bit corrisponde ad un processore, ma oggi esistono architetture in
cui questo numero può non essere sufficiente, e per questo è stato creato
questo tipo opaco e una interfaccia di gestione che permette di usare a basso
livello un tipo di dato qualunque rendendosi indipendenti dal numero di bit e
dalla loro disposizione.  Per questo le funzioni di libreria richiedono che
oltre all'insieme di processori si indichi anche la dimensione dello stesso
con l'argomento \param{setsize}, per il quale, se non si usa l'allocazione
dinamica che vedremo a breve, è in genere sufficiente passare il valore
\code{sizeof(cpu\_set\_t)}.

L'interfaccia di gestione degli insiemi di processori, oltre alla definizione
del tipo \type{cpu\_set\_t}, prevede una serie di macro di preprocessore per
la manipolazione degli stessi. Quelle di base, che consentono rispettivamente
di svuotare un insieme, di aggiungere o togliere un processore o di verificare
se esso è già presente in un insieme, sono le seguenti:

{\centering
\vspace{3pt}
\begin{funcbox}{ 
\fhead{sched.h}
\fdecl{void \macrod{CPU\_ZERO}(cpu\_set\_t *set)}
\fdesc{Inizializza un insieme di processori vuoto \param{set}.} 
\fdecl{void \macrod{CPU\_SET}(int cpu, cpu\_set\_t *set)}
\fdesc{Inserisce il processore \param{cpu} nell'insieme di
  processori \param{set}.}  
\fdecl{void \macrod{CPU\_CLR}(int cpu, cpu\_set\_t *set)}
\fdesc{Rimuove il processore \param{cpu} nell'insieme di
  processori \param{set}.}  
\fdecl{int \macrod{CPU\_ISSET}(int cpu, cpu\_set\_t *set)}
\fdesc{Controlla se il processore \param{cpu} è nell'insieme di
  processori \param{set}.}  
}
\end{funcbox}}

Queste macro che sono ispirate dalle analoghe usate per gli insiemi di
\textit{file descriptor} (vedi sez.~\ref{sec:file_select}) e sono state
introdotte con la versione 2.3.3 della \acr{glibc}. Tutte richiedono che si
specifichi il numero di una CPU nell'argomento \param{cpu}, ed un insieme su
cui operare. L'unica che ritorna un risultato è \macro{CPU\_ISSET}, che
restituisce un intero da usare come valore logico (zero se la CPU non è
presente, diverso da zero se è presente).

\itindbeg{side~effects}
Si tenga presente che trattandosi di macro l'argomento \param{cpu} può essere
valutato più volte. Questo significa ad esempio che non si può usare al suo
posto una funzione o un'altra macro, altrimenti queste verrebbero eseguite più
volte; l'argomento cioè non deve avere \textsl{effetti collaterali} (in gergo
 \textit{side effects}).\footnote{nel linguaggio C si
  parla appunto di \textit{side effects} quando si usano istruzioni la cui
  valutazione comporta effetti al di fuori dell'istruzione stessa, come il
  caso indicato in cui si passa una funzione ad una macro che usa l'argomento
  al suo interno più volte, o si scrivono espressioni come \code{a=a++} in cui
  non è chiaro se prima avvenga l'incremento e poi l'assegnazione, ed il cui
  risultato dipende dall'implementazione del compilatore.}
\itindend{side~effects}


Le CPU sono numerate da zero (che indica la prima disponibile) fino ad un
numero massimo che dipende dall'architettura hardware. La costante
\constd{CPU\_SETSIZE} indica il numero massimo di processori che possono far
parte di un insieme (al momento vale sempre 1024), e costituisce un limite
massimo al valore dell'argomento \param{cpu}.  Dalla versione 2.6 della
\acr{glibc} alle precedenti macro è stata aggiunta, per contare il numero di
processori in un insieme, l'ulteriore:

{\centering
\vspace{3pt}
\begin{funcbox}{ 
\fhead{sched.h}
\fdecl{int \macrod{CPU\_COUNT}(cpu\_set\_t *set)}
\fdesc{Conta il numero di processori presenti nell'insieme \param{set}.} 
}
\end{funcbox}}

A partire dalla versione 2.7 della \acr{glibc} sono state introdotte altre
macro che consentono ulteriori manipolazioni, in particolare si possono
compiere delle operazioni logiche sugli insiemi di processori con:

{\centering
\vspace{3pt}
\begin{funcbox}{ 
\fhead{sched.h}
\fdecl{void \macrod{CPU\_AND}(cpu\_set\_t *destset, cpu\_set\_t *srcset1, cpu\_set\_t *srcset2)}
\fdesc{Esegue l'AND logico di due insiemi di processori.} 
\fdecl{void \macrod{CPU\_OR}(cpu\_set\_t *destset, cpu\_set\_t *srcset1, cpu\_set\_t *srcset2)}
\fdesc{Esegue l'OR logico di due insiemi di processori.} 
\fdecl{void \macrod{CPU\_XOR}(cpu\_set\_t *destset, cpu\_set\_t *srcset1, cpu\_set\_t *srcset2)}
\fdesc{Esegue lo XOR logico di due insiemi di processori.} 
\fdecl{int \macrod{CPU\_EQUAL}(cpu\_set\_t *set1, cpu\_set\_t *set2)}
\fdesc{Verifica se due insiemi di processori sono uguali.} 
}
\end{funcbox}}

Le prime tre macro richiedono due insiemi di partenza, \param{srcset1} e
\param{srcset2} e forniscono in un terzo insieme \param{destset} (che può
essere anche lo stesso di uno dei precedenti) il risultato della rispettiva
operazione logica sui contenuti degli stessi. In sostanza con \macro{CPU\_AND}
si otterrà come risultato l'insieme che contiene le CPU presenti in entrambi
gli insiemi di partenza, con \macro{CPU\_OR} l'insieme che contiene le CPU
presenti in uno qualunque dei due insiemi di partenza, e con \macro{CPU\_XOR}
l'insieme che contiene le CPU presenti in uno solo dei due insiemi di
partenza. Infine \macro{CPU\_EQUAL} confronta due insiemi ed è l'unica che
restituisce un intero, da usare come valore logico che indica se sono identici
o meno.

Inoltre, sempre a partire dalla versione 2.7 della \acr{glibc}, è stata
introdotta la possibilità di una allocazione dinamica degli insiemi di
processori, per poterli avere di dimensioni corrispondenti al numero di CPU
effettivamente in gioco, senza dover fare riferimento necessariamente alla
precedente dimensione preimpostata di 1024. Per questo motivo sono state
definite tre ulteriori macro, che consentono rispettivamente di allocare,
disallocare ed ottenere la dimensione in byte di un insieme di processori:

{\centering
\vspace{3pt}
\begin{funcbox}{ 
\fhead{sched.h}
\fdecl{cpu\_set\_t * \macrod{CPU\_ALLOC}(num\_cpus)}
\fdesc{Alloca dinamicamente un insieme di processori di dimensione voluta.} 
\fdecl{void \macrod{CPU\_FREE}(cpu\_set\_t *set)}
\fdesc{Disalloca un insieme di processori allocato dinamicamente.} 
\fdecl{size\_t \macrod{CPU\_ALLOC\_SIZE}(num\_cpus)}
\fdesc{Ritorna la dimensione di un insieme di processori allocato dinamicamente.} 
}
\end{funcbox}}

La prima macro, \macro{CPU\_ALLOC}, restituisce il puntatore ad un insieme di
processori in grado di contenere almeno \param{num\_cpus} che viene allocato
dinamicamente. Ogni insieme così allocato dovrà essere disallocato con
\macro{CPU\_FREE} passandogli un puntatore ottenuto da una precedente
\macro{CPU\_ALLOC}. La terza macro, \macro{CPU\_ALLOC\_SIZE}, consente di
ottenere la dimensione in byte di un insieme allocato dinamicamente che
contenga \param{num\_cpus} processori.

Dato che le dimensioni effettive possono essere diverse le macro di gestione e
manipolazione che abbiamo trattato in precedenza non si applicano agli insiemi
allocati dinamicamente, per i quali dovranno sono state definite altrettante
macro equivalenti contraddistinte dal suffisso \texttt{\_S}, che effettuano le
stesse operazioni, ma richiedono in più un argomento
aggiuntivo \param{setsize} che deve essere assegnato al valore ottenuto con
\macro{CPU\_ALLOC\_SIZE}. Questo stesso valore deve essere usato per l'omonimo
argomento delle funzioni \func{sched\_setaffinity} o \func{sched\_getaffinity}
quando si vuole usare per l'argomento che indica la maschera di affinità un
insieme di processori allocato dinamicamente.

\index{insieme~di~processori|)}

A meno di non aver utilizzato \func{sched\_setaffinity}, in condizioni
ordinarie la maschera di affinità di un processo è preimpostata dal sistema in
modo che esso possa essere eseguito su qualunque processore. Se ne può
comunque ottenere il valore corrente usando la funzione di sistema
\funcd{sched\_getaffinity}, il cui prototipo è:

\begin{funcproto}{ 
\fhead{sched.h}
\fdecl{int sched\_getaffinity (pid\_t pid, size\_t setsize, 
  cpu\_set\_t *mask)}
\fdesc{Legge la maschera di affinità di un processo.} 
}
{La funzione ritorna $0$ in caso di successo e $-1$ per un errore, nel qual
caso \var{errno} assumerà uno dei valori:
\begin{errlist}
\item[\errcode{EINVAL}] \param{setsize} è più piccolo delle dimensioni
  della maschera di affinità usata dal kernel.
\item[\errcode{ESRCH}] il processo \param{pid} non esiste.
\end{errlist}
ed inoltre anche \errval{EFAULT} nel suo significato generico.}
\end{funcproto}

La funzione restituirà all'indirizzo specificato da \param{mask} il valore
della maschera di affinità del processo indicato dall'argomento \param{pid}
(al solito un valore nullo indica il processo corrente) così da poterla
riutilizzare per una successiva reimpostazione.

È chiaro che queste funzioni per la gestione dell'affinità hanno significato
soltanto su un sistema multiprocessore, esse possono comunque essere
utilizzate anche in un sistema con un processore singolo, nel qual caso però
non avranno alcun risultato effettivo.

\itindend{scheduler}
\itindend{CPU~affinity}


\subsection{Le priorità per le operazioni di I/O}
\label{sec:io_priority}

A lungo l'unica priorità usata per i processi è stata quella relativa
all'assegnazione dell'uso del processore. Ma il processore non è l'unica
risorsa che i processi devono contendersi, un'altra, altrettanto importante
per le prestazioni, è quella dell'accesso a disco. Per questo motivo nello
sviluppo del kernel sono stati introdotti diversi \textit{I/O scheduler} in
grado di distribuire in maniera opportuna questa risorsa ai vari processi.

Fino al kernel 2.6.17 era possibile soltanto differenziare le politiche
generali di gestione, scegliendo di usare un diverso \textit{I/O scheduler}. A
partire da questa versione, con l'introduzione dello \textit{scheduler} CFQ
(\textit{Completely Fair Queuing}) è divenuto possibile, qualora si usi questo
\textit{scheduler}, impostare anche delle diverse priorità di accesso per i
singoli processi.\footnote{al momento (kernel 2.6.31), le priorità di I/O sono
  disponibili soltanto per questo \textit{scheduler}.}

La scelta di uno \textit{scheduler} di I/O si può fare in maniera generica per
tutto il sistema all'avvio del kernel con il parametro di avvio
\texttt{elevator},\footnote{per la trattazione dei parametri di avvio del
  kernel si rimanda al solito alla sez.~5.3 di \cite{AGL}.} cui assegnare il
nome dello \textit{scheduler}, ma se ne può anche indicare uno specifico per
l'accesso al singolo disco scrivendo nel file
\texttt{/sys/block/\textit{<dev>}/queue/scheduler} (dove
\texttt{\textit{<dev>}} è il nome del dispositivo associato al disco).

Gli \textit{scheduler} disponibili sono mostrati dal contenuto dello stesso
file che riporta fra parentesi quadre quello attivo, il default in tutti i
kernel recenti è proprio il \texttt{cfq}, nome con cui si indica appunto lo
\textit{scheduler} CFQ, che supporta le priorità. Per i dettagli sulle
caratteristiche specifiche degli altri \textit{scheduler}, la cui discussione
attiene a problematiche di ambito sistemistico, si consulti la documentazione
nella directory \texttt{Documentation/block/} dei sorgenti del kernel.

Una volta che si sia impostato lo \textit{scheduler} CFQ ci sono due
specifiche \textit{system call}, specifiche di Linux, che consentono di
leggere ed impostare le priorità di I/O.\footnote{se usate in corrispondenza
  ad uno \textit{scheduler} diverso il loro utilizzo non avrà alcun effetto.}
Dato che non esiste una interfaccia diretta nella \acr{glibc} per queste due
funzioni\footnote{almeno al momento della scrittura di questa sezione, con la
  versione 2.11 della \acr{glibc}.} occorrerà invocarle tramite la funzione
\func{syscall} (come illustrato in sez.~\ref{sec:proc_syscall}). Le due
\textit{system call} sono \funcd{ioprio\_get} ed \funcd{ioprio\_set}; i
rispettivi prototipi sono:

\begin{funcproto}{ 
\fhead{linux/ioprio.h}
\fdecl{int ioprio\_get(int which, int who)}
\fdesc{Legge la priorità di I/O di un processo.} 
\fdecl{int ioprio\_set(int which, int who, int ioprio)}
\fdesc{Imposta la priorità di I/O di un processo.} 
}
{Le funzioni ritornano rispettivamente un intero positivo o 0 in caso di
  successo e $-1$ per un errore, nel qual caso \var{errno} assumerà uno dei
  valori:
\begin{errlist}
\item[\errcode{EINVAL}] i valori di \param{which} o di \param{ioprio} non
  sono validi. 
\item[\errcode{EPERM}] non si hanno i privilegi per eseguire
  l'impostazione (solo per \func{ioprio\_set}). 
\item[\errcode{ESRCH}] non esiste un processo corrispondente alle indicazioni.
\end{errlist}}
\end{funcproto}

Le funzioni leggono o impostano la priorità di I/O sulla base dell'indicazione
dei due argomenti \param{which} e \param{who} che hanno lo stesso significato
già visto per gli omonimi argomenti di \func{getpriority} e
\func{setpriority}. Anche in questo caso si deve specificare il valore di
\param{which} tramite le opportune costanti riportate in
tab.~\ref{tab:ioprio_args} che consentono di indicare un singolo processo, i
processi di un \textit{process group} (vedi sez.~\ref{sec:sess_proc_group}) o
tutti i processi di un utente.

\begin{table}[htb]
  \centering
  \footnotesize
  \begin{tabular}[c]{|c|c|l|}
    \hline
    \param{which} & \param{who} & \textbf{Significato} \\
    \hline
    \hline
    \constd{IPRIO\_WHO\_PROCESS} & \type{pid\_t} & processo\\
    \constd{IPRIO\_WHO\_PRGR}    & \type{pid\_t} & \textit{process group}\\ 
    \constd{IPRIO\_WHO\_USER}    & \type{uid\_t} & utente\\
    \hline
  \end{tabular}
  \caption{Legenda del valore dell'argomento \param{which} e del tipo
    dell'argomento \param{who} delle funzioni \func{ioprio\_get} e
    \func{ioprio\_set} per le tre possibili scelte.}
  \label{tab:ioprio_args}
\end{table}

In caso di successo \func{ioprio\_get} restituisce un intero positivo che
esprime il valore della priorità di I/O, questo valore è una maschera binaria
composta da due parti, una che esprime la \textsl{classe} di
\textit{scheduling} di I/O del processo, l'altra che esprime, quando la classe
di \textit{scheduling} lo prevede, la priorità del processo all'interno della
classe stessa. Questo stesso formato viene utilizzato per indicare il valore
della priorità da impostare con l'argomento \param{ioprio} di
\func{ioprio\_set}.

\begin{table}[htb]
  \centering
  \footnotesize
  \begin{tabular}[c]{|l|p{8cm}|}
    \hline
    \textbf{Macro} & \textbf{Significato}\\
    \hline
    \hline
    \macrod{IOPRIO\_PRIO\_CLASS}\texttt{(\textit{value})}
                                & Dato il valore di una priorità come
                                  restituito da \func{ioprio\_get} estrae il
                                  valore della classe.\\
    \macrod{IOPRIO\_PRIO\_DATA}\texttt{(\textit{value})}
                                & Dato il valore di una priorità come
                                  restituito da \func{ioprio\_get} estrae il
                                  valore della priorità.\\
    \macrod{IOPRIO\_PRIO\_VALUE}\texttt{(\textit{class},\textit{prio})}
                                & Dato un valore di priorità ed una classe
                                  ottiene il valore numerico da passare a
                                  \func{ioprio\_set}.\\
    \hline
  \end{tabular}
  \caption{Le macro per la gestione dei valori numerici .}
  \label{tab:IOsched_class_macro}
\end{table}


Per la gestione dei valori che esprimono le priorità di I/O sono state
definite delle opportune macro di preprocessore, riportate in
tab.~\ref{tab:IOsched_class_macro}. I valori delle priorità si ottengono o si
impostano usando queste macro.  Le prime due si usano con il valore restituito
da \func{ioprio\_get} e per ottenere rispettivamente la classe di
\textit{scheduling}\footnote{restituita dalla macro con i valori di
  tab.~\ref{tab:IOsched_class}.} e l'eventuale valore della priorità. La terza
macro viene invece usata per creare un valore di priorità da usare come
argomento di \func{ioprio\_set} per eseguire una impostazione.

Le classi di \textit{scheduling} previste dallo \textit{scheduler} CFQ sono
tre, e ricalcano tre diverse modalità di distribuzione delle risorse, analoghe
a quelle già adottate anche nel funzionamento dello \textit{scheduler} del
processore. Ciascuna di esse è identificata tramite una opportuna costante,
secondo quanto riportato in tab.~\ref{tab:IOsched_class}.

\begin{table}[htb]
  \centering
  \footnotesize
  \begin{tabular}[c]{|l|l|}
    \hline
    \textbf{Classe}  & \textbf{Significato} \\
    \hline
    \hline
    \const{IOPRIO\_CLASS\_RT}  & \textit{Scheduling} di I/O
                                 \textit{real-time}.\\  
    \const{IOPRIO\_CLASS\_BE}  & \textit{Scheduling} di I/O ordinario.\\ 
    \const{IOPRIO\_CLASS\_IDLE}& \textit{Scheduling} di I/O di priorità
                                  minima.\\
    \hline
  \end{tabular}
  \caption{Costanti che identificano le classi di \textit{scheduling} di I/O.}
  \label{tab:IOsched_class}
\end{table}

La classe di priorità più bassa è \constd{IOPRIO\_CLASS\_IDLE}; i processi in
questa classe riescono ad accedere a disco soltanto quando nessun altro
processo richiede l'accesso. Occorre pertanto usarla con molta attenzione,
perché un processo in questa classe può venire completamente bloccato quando
ci sono altri processi in una qualunque delle altre due classi che stanno
accedendo al disco. Quando si usa questa classe non ha senso indicare un
valore di priorità, dato che in questo caso non esiste nessuna gerarchia e la
priorità è identica, la minima possibile, per tutti i processi che la usano.

La seconda classe di priorità di I/O è \constd{IOPRIO\_CLASS\_BE} (il nome sta
per \textit{best-effort}), che è quella usata ordinariamente da tutti
processi. In questo caso esistono priorità diverse che consentono di
assegnazione di una maggiore banda passante nell'accesso a disco ad un
processo rispetto agli altri, con meccanismo simile a quello dei valori di
\textit{nice} in cui si evita che un processo a priorità più alta possa
bloccare indefinitamente quelli a priorità più bassa. In questo caso però le
diverse priorità sono soltanto otto, indicate da un valore numerico fra 0 e 7
e come per \textit{nice} anche in questo caso un valore più basso indica una
priorità maggiore.

Infine la classe di priorità di I/O \textit{real-time}
\constd{IOPRIO\_CLASS\_RT} ricalca le omonime priorità di processore: un
processo in questa classe ha sempre la precedenza nell'accesso a disco
rispetto a tutti i processi delle altre classi e di un processo nella stessa
classe ma con priorità inferiore, ed è pertanto in grado di bloccare
completamente tutti gli altri. Anche in questo caso ci sono 8 priorità diverse
con un valore numerico fra 0 e 7, con una priorità più elevata per valori più
bassi.

In generale nel funzionamento ordinario la priorità di I/O di un processo
viene impostata in maniera automatica nella classe \const{IOPRIO\_CLASS\_BE}
con un valore ottenuto a partire dal corrispondente valore di \textit{nice}
tramite la formula: $\mathtt{\mathit{prio}}=(\mathtt{\mathit{nice}}+20)/5$. Un
utente ordinario può modificare con \func{ioprio\_set} soltanto le priorità
dei processi che gli appartengono,\footnote{per la modifica delle priorità di
  altri processi occorrono privilegi amministrativi, ed in particolare la
  capacità \const{CAP\_SYS\_NICE} (vedi sez.~\ref{sec:proc_capabilities}).}
cioè quelli il cui \ids{UID} reale corrisponde all'\ids{UID} reale o effettivo
del chiamante. Data la possibilità di ottenere un blocco totale del sistema,
solo l'amministratore\footnote{o un processo con la capacità
  \const{CAP\_SYS\_ADMIN} (vedi sez.~\ref{sec:proc_capabilities}).} può
impostare un processo ad una priorità di I/O nella classe
\const{IOPRIO\_CLASS\_RT}, lo stesso privilegio era richiesto anche per la
classe \const{IOPRIO\_CLASS\_IDLE} fino al kernel 2.6.24, ma dato che in
questo caso non ci sono effetti sugli altri processi questo limite è stato
rimosso a partire dal kernel 2.6.25.

%TODO verificare http://lwn.net/Articles/355987/


\section{Problematiche di programmazione \textit{multitasking}}
\label{sec:proc_multi_prog}

Benché i processi siano strutturati in modo da apparire il più possibile come
indipendenti l'uno dall'altro, nella programmazione in un sistema
\textit{multitasking} occorre tenere conto di una serie di problematiche che
normalmente non esistono quando si ha a che fare con un sistema in cui viene
eseguito un solo programma alla volta.

Per questo motivo, essendo questo argomento di carattere generale, ci è parso
opportuno introdurre sinteticamente queste problematiche, che ritroveremo a
più riprese in capitoli successivi, in questa sezione conclusiva del capitolo
in cui abbiamo affrontato la gestione dei processi, sottolineando come esse
diventino cogenti quando invece si usano i \textit{thread}.


\subsection{Le operazioni atomiche}
\label{sec:proc_atom_oper}

La nozione di \textsl{operazione atomica} deriva dal significato greco della
parola atomo, cioè indivisibile; si dice infatti che un'operazione è atomica
quando si ha la certezza che, qualora essa venga effettuata, tutti i passaggi
che devono essere compiuti per realizzarla verranno eseguiti senza possibilità
di interruzione in una fase intermedia.

In un ambiente \textit{multitasking} il concetto è essenziale, dato che un
processo può essere interrotto in qualunque momento dal kernel che mette in
esecuzione un altro processo o dalla ricezione di un segnale. Occorre pertanto
essere accorti nei confronti delle possibili \textit{race condition} (vedi
sez.~\ref{sec:proc_race_cond}) derivanti da operazioni interrotte in una fase
in cui non erano ancora state completate.

Nel caso dell'interazione fra processi la situazione è molto più semplice, ed
occorre preoccuparsi della atomicità delle operazioni solo quando si ha a che
fare con meccanismi di intercomunicazione (che esamineremo in dettaglio in
cap.~\ref{cha:IPC}) o nelle operazioni con i file (vedremo alcuni esempi in
sez.~\ref{sec:file_shared_access}). In questi casi in genere l'uso delle
appropriate funzioni di libreria per compiere le operazioni necessarie è
garanzia sufficiente di atomicità in quanto le \textit{system call} con cui
esse sono realizzate non possono essere interrotte (o subire interferenze
pericolose) da altri processi.

Nel caso dei segnali invece la situazione è molto più delicata, in quanto lo
stesso processo, e pure alcune \textit{system call}, possono essere interrotti
in qualunque momento, e le operazioni di un eventuale \textit{signal handler}
sono compiute nello stesso spazio di indirizzi del processo. Per questo, anche
il solo accesso o l'assegnazione di una variabile possono non essere più
operazioni atomiche (torneremo su questi aspetti in
sez.~\ref{sec:sig_adv_control}).

Qualora invece si usino i \textit{thread}, in cui lo spazio degli indirizzi è
condiviso, il problema è sempre presente, perché qualunque \textit{thread} può
interromperne un altro in qualunque momento e l'atomicità di qualunque
operazione è messa in discussione, per cui l'assenza di eventuali \textit{race
  condition} (vedi sez.~\ref{sec:proc_race_cond}) deve essere sempre
verificata nei minimi dettagli.

In questo caso il sistema provvede un tipo di dato, il \typed{sig\_atomic\_t},
il cui accesso è assicurato essere atomico.  In pratica comunque si può
assumere che, in ogni piattaforma su cui è implementato Linux, il tipo
\ctyp{int}, gli altri interi di dimensione inferiore ed i puntatori sono
atomici. Non è affatto detto che lo stesso valga per interi di dimensioni
maggiori (in cui l'accesso può comportare più istruzioni in assembler) o per
le strutture di dati. In tutti questi casi è anche opportuno marcare come
\dirct{volatile} le variabili che possono essere interessate ad accesso
condiviso, onde evitare problemi con le ottimizzazioni del codice.



\subsection{Le \textit{race condition} ed i \textit{deadlock}}
\label{sec:proc_race_cond}

\itindbeg{race~condition}

Si definiscono \textit{race condition} tutte quelle situazioni in cui processi
diversi operano su una risorsa comune, ed in cui il risultato viene a
dipendere dall'ordine in cui essi effettuano le loro operazioni. Il caso
tipico è quello di un'operazione che viene eseguita da un processo in più
passi, e può essere compromessa dall'intervento di un altro processo che
accede alla stessa risorsa quando ancora non tutti i passi sono stati
completati.

Dato che in un sistema \textit{multitasking} ogni processo può essere
interrotto in qualunque momento per farne subentrare un altro in esecuzione,
niente può assicurare un preciso ordine di esecuzione fra processi diversi o
che una sezione di un programma possa essere eseguita senza interruzioni da
parte di altri. Queste situazioni comportano pertanto errori estremamente
subdoli e difficili da tracciare, in quanto nella maggior parte dei casi tutto
funzionerà regolarmente, e solo occasionalmente si avranno degli errori.

Per questo occorre essere ben consapevoli di queste problematiche, e del fatto
che l'unico modo per evitarle è quello di riconoscerle come tali e prendere
gli adeguati provvedimenti per far sì che non si verifichino. Casi tipici di
\textit{race condition} si hanno quando diversi processi accedono allo stesso
file, o nell'accesso a meccanismi di intercomunicazione come la memoria
condivisa. 

\index{sezione~critica|(}

In questi casi, se non si dispone della possibilità di eseguire atomicamente
le operazioni necessarie, occorre che quelle parti di codice in cui si
compiono le operazioni sulle risorse condivise, quelle che in genere vengono
denominate ``\textsl{sezioni critiche}'' del programma, siano opportunamente
protette da meccanismi di sincronizzazione (vedremo alcune problematiche di
questo tipo in cap.~\ref{cha:IPC}).

\index{sezione~critica|)}

Nel caso dei \textit{thread} invece la situazione è molto più delicata e
sostanzialmente qualunque accesso in memoria (a buffer, variabili o altro) può
essere soggetto a \textit{race condition} dato potrebbe essere interrotto in
qualunque momento da un altro \textit{thread}. In tal caso occorre pianificare
con estrema attenzione l'uso delle variabili ed utilizzare i vari meccanismi
di sincronizzazione che anche in questo caso sono disponibili (torneremo su
queste problematiche di questo tipo in cap.~\ref{sec:pthread_sync})

\itindbeg{deadlock} 

Un caso particolare di \textit{race condition} sono poi i cosiddetti
\textit{deadlock} (traducibile in \textsl{condizione di stallo}), che
particolarmente gravi in quanto comportano spesso il blocco completo di un
servizio, e non il fallimento di una singola operazione. Per definizione un
\textit{deadlock} è una situazione in cui due o più processi non sono più in
grado di proseguire perché ciascuno aspetta il risultato di una operazione che
dovrebbe essere eseguita dall'altro.

L'esempio tipico di una situazione che può condurre ad un
\textit{deadlock} è quello in cui un flag di
``\textsl{occupazione}'' viene rilasciato da un evento asincrono (come un
segnale o un altro processo) fra il momento in cui lo si è controllato
(trovandolo occupato) e la successiva operazione di attesa per lo sblocco. In
questo caso, dato che l'evento di sblocco del flag è avvenuto senza che ce ne
accorgessimo proprio fra il controllo e la messa in attesa, quest'ultima
diventerà perpetua (da cui il nome di \textit{deadlock}).

In tutti questi casi è di fondamentale importanza il concetto di atomicità
visto in sez.~\ref{sec:proc_atom_oper}; questi problemi infatti possono essere
risolti soltanto assicurandosi, quando essa sia richiesta, che sia possibile
eseguire in maniera atomica le operazioni necessarie.

\itindend{race~condition}
\itindend{deadlock}


\subsection{Le funzioni rientranti}
\label{sec:proc_reentrant}

\index{funzioni!rientranti|(}

Si dice \textsl{rientrante} una funzione che può essere interrotta in
qualunque punto della sua esecuzione ed essere chiamata una seconda volta da
un altro \textit{thread} di esecuzione senza che questo comporti nessun
problema nell'esecuzione della stessa. La problematica è comune nella
programmazione con i \textit{thread}, ma si hanno gli stessi problemi quando
si vogliono chiamare delle funzioni all'interno dei gestori dei segnali.

Fintanto che una funzione opera soltanto con le variabili locali è rientrante;
queste infatti vengono allocate nello \textit{stack}, ed un'altra invocazione
non fa altro che allocarne un'altra copia. Una funzione può non essere
rientrante quando opera su memoria che non è nello \textit{stack}.  Ad esempio
una funzione non è mai rientrante se usa una variabile globale o statica.

Nel caso invece la funzione operi su un oggetto allocato dinamicamente, la
cosa viene a dipendere da come avvengono le operazioni: se l'oggetto è creato
ogni volta e ritornato indietro la funzione può essere rientrante, se invece
esso viene individuato dalla funzione stessa due chiamate alla stessa funzione
potranno interferire quando entrambe faranno riferimento allo stesso oggetto.
Allo stesso modo una funzione può non essere rientrante se usa e modifica un
oggetto che le viene fornito dal chiamante: due chiamate possono interferire
se viene passato lo stesso oggetto; in tutti questi casi occorre molta cura da
parte del programmatore.

In genere le funzioni di libreria non sono rientranti, molte di esse ad
esempio utilizzano variabili statiche, la \acr{glibc} però mette a
disposizione due macro di compilatore, \macro{\_REENTRANT} e
\macro{\_THREAD\_SAFE}, la cui definizione attiva le versioni rientranti di
varie funzioni di libreria, che sono identificate aggiungendo il suffisso
\code{\_r} al nome della versione normale.

\index{funzioni!rientranti|)}


% LocalWords:  multitasking like VMS child process identifier pid sez shell fig
% LocalWords:  parent kernel init pstree keventd kswapd table struct linux call
% LocalWords:  nell'header scheduler system interrupt timer HZ asm Hertz clock
% LocalWords:  l'alpha tick fork wait waitpid exit exec image glibc int pgid ps
% LocalWords:  sid thread Ingo Molnar ppid getpid getppid sys unistd LD threads
% LocalWords:  void tempnam pathname sibling cap errno EAGAIN ENOMEM context
% LocalWords:  stack read only copy write tab client spawn forktest sleep PATH
% LocalWords:  source LIBRARY scheduling race condition printf descriptor dup
% LocalWords:  close group session tms lock vfork execve BSD stream main abort
% LocalWords:  SIGABRT SIGCHLD SIGHUP foreground SIGCONT termination signal ANY
% LocalWords:  handler kill EINTR POSIX options WNOHANG ECHILD option WUNTRACED
% LocalWords:  dump bits rusage getrusage heap const filename argv envp EACCES
% LocalWords:  filesystem noexec EPERM suid sgid root nosuid ENOEXEC ENOENT ELF
% LocalWords:  ETXTBSY EINVAL ELIBBAD BIG EFAULT EIO ENAMETOOLONG ELOOP ENOTDIR
% LocalWords:  ENFILE EMFILE argc execl path execv execle execlp execvp vector
% LocalWords:  list environ NULL umask utime cutime ustime fcntl linker Posix
% LocalWords:  opendir libc interpreter FreeBSD capabilities mandatory access
% LocalWords:  control MAC SELinux security modules LSM superuser uid gid saved
% LocalWords:  effective euid egid dell' fsuid fsgid getuid geteuid getgid SVr
% LocalWords:  getegid IDS NFS setuid setgid all' logout utmp screen xterm TODO
% LocalWords:  setreuid setregid FIXME ruid rgid seteuid setegid setresuid size
% LocalWords:  setresgid getresuid getresgid value result argument setfsuid DAC
% LocalWords:  setfsgid NGROUPS sysconf getgroups getgrouplist groups ngroups
% LocalWords:  setgroups initgroups patch LIDS CHOWN OVERRIDE Discrectionary 
% LocalWords:  SEARCH chattr sticky NOATIME socket domain immutable append mmap
% LocalWords:  broadcast multicast multicasting memory locking mlock mlockall
% LocalWords:  shmctl ioperm iopl chroot ptrace accounting swap reboot hangup
% LocalWords:  vhangup mknod lease permitted inherited inheritable bounding AND
% LocalWords:  capability capget capset header ESRCH undef version obj clear PT
% LocalWords:  pag ssize length proc capgetp preemptive cache runnable idled
% LocalWords:  SIGSTOP soft slice nice niceness counter which SC switch side
% LocalWords:  getpriority who setpriority RTLinux RTAI Adeos fault FIFO  COUNT
% LocalWords:  yield Robin setscheduler policy param OTHER priority setparam to
% LocalWords:  min getparam getscheduler interval robin ENOSYS fifo ping long
% LocalWords:  affinity setaffinity unsigned mask cpu NUMA CLR ISSET SETSIZE RR
% LocalWords:  getaffinity assembler deadlock REENTRANT SAFE tgz MYPGRP l'OR rr
% LocalWords:  WIFEXITED WEXITSTATUS WIFSIGNALED WTERMSIG WCOREDUMP WIFSTOPPED
% LocalWords:  WSTOPSIG opt char INTERP arg SIG IGN DFL mascheck grp FOWNER RAW
% LocalWords:  FSETID SETPCAP BIND SERVICE ADMIN PACKET IPC OWNER MODULE RAWIO
% LocalWords:  PACCT RESOURCE TTY CONFIG SETFCAP hdrp datap libcap lcap text tp
% LocalWords:  get ncap caps CapInh CapPrm fffffeff CapEff getcap STAT dall'I
% LocalWords:  inc PRIO SUSv PRGR prio SysV SunOS Ultrix sched timespec len sig
% LocalWords:  cpusetsize cpuset atomic tickless redirezione WCONTINUED stopped
% LocalWords:  waitid NOCLDSTOP ENOCHLD WIFCONTINUED ifdef endif idtype siginfo
% LocalWords:  infop ALL WEXITED WSTOPPED WNOWAIT signo CLD EXITED KILLED page
% LocalWords:  CONTINUED sources forking Spawned successfully executing exiting
% LocalWords:  next cat for COMMAND pts bash defunct TRAPPED DUMPED PR effects
% LocalWords:  SIGKILL static RLIMIT preemption PREEMPT VOLUNTARY IDLE RTPRIO
% LocalWords:  completely fair compat uniform CFQ queuing elevator dev cfq RT
% LocalWords:  documentation block syscall ioprio IPRIO CLASS class best effort
% LocalWords:  refresh semop dnotify MADV DONTFORK prctl WCLONE WALL big mount
% LocalWords:  WNOTHREAD DUMPABLE KEEPCAPS IRIX CAPBSET endianness endian flags
% LocalWords:  little PPC PowerPC FPEMU NOPRINT SIGFPE FPEXC point FP SW malloc
% LocalWords:  exception EXC ENABLE OVF overflow UND underflow RES INV DISABLED
% LocalWords:  NONRECOV ASYNC KEEP securebits NAME NUL PDEATHSIG SECCOMP VM FS
% LocalWords:  secure computing sigreturn TIMING STATISTICAL TSC MCE conditions
% LocalWords:  timestamp Stamp SIGSEGV UNALIGN SIGBUS MCEERR AO failure early
% LocalWords:  namespace vsyscall SETTID FILES NEWIPC NEWNET NEWNS NEWPID ptid
% LocalWords:  NEWUTS SETTLS SIGHAND SYSVSEM UNTRACED tls ctid CLEARTID panic
% LocalWords:  loader EISDIR SIGTRAP uninterrutible killable EQUAL sizeof XOR
% LocalWords:  destset srcset ALLOC num cpus setsize emacs pager getty TID
%  LocalWords:  reaper SUBREAPER Library futex klogd named rpc statd NPROC
%  LocalWords:  EACCESS EBADF EBUSY ENXIO EOPNOTSUPP DISABLE tracer Yama
 
%%% Local Variables: 
%%% mode: latex
%%% TeX-master: "gapil"
%%% End: 
